%% LyX 2.4.0-alpha3 created this file.  For more info, see https://www.lyx.org/.
%% Do not edit unless you really know what you are doing.
\documentclass[english,tableposition=top]{report}
\usepackage{lmodern}
\renewcommand{\sfdefault}{lmss}
\renewcommand{\ttdefault}{lmtt}
\usepackage[T1]{fontenc}
\usepackage{textcomp}
\usepackage[utf8]{inputenc}
\setcounter{secnumdepth}{3}
\setcounter{tocdepth}{3}
\usepackage{color}
\definecolor{shadecolor}{rgb}{0.667969, 1, 1}
\usepackage{babel}
\usepackage{array}
\usepackage{cprotect}
\usepackage{wrapfig}
\usepackage{booktabs}
\usepackage{framed}
\usepackage{url}
\usepackage{amsmath}
\usepackage{amssymb}
\usepackage[unicode=true,pdfusetitle,
 bookmarks=true,bookmarksnumbered=true,bookmarksopen=true,bookmarksopenlevel=2,
 breaklinks=true,pdfborder={0 0 1},backref=section,colorlinks=true]
 {hyperref}

\makeatletter

%%%%%%%%%%%%%%%%%%%%%%%%%%%%%% LyX specific LaTeX commands.
\providecommand{\LyX}{\texorpdfstring{\ensureascii{%
  L\kern-.1667em\lower.25em\hbox{Y}\kern-.125emX\@}}{LyX}}
\DeclareRobustCommand*{\lyxarrow}{%
\@ifstar
{\leavevmode\,$\triangleleft$\,\allowbreak}
{\leavevmode\,$\triangleright$\,\allowbreak}}
%% Because html converters don't know tabularnewline
\providecommand{\tabularnewline}{\\}

%%%%%%%%%%%%%%%%%%%%%%%%%%%%%% Textclass specific LaTeX commands.
\newenvironment{centred}%
  {\begin{center}\baselineskip=13pt\parskip=1pt}{\end{center}}
\newenvironment{lyxcode}
	{\par\begin{list}{}{
		\setlength{\rightmargin}{\leftmargin}
		\setlength{\listparindent}{0pt}% needed for AMS classes
		\raggedright
		\setlength{\itemsep}{0pt}
		\setlength{\parsep}{0pt}
		\normalfont\ttfamily}%
	 \item[]}
	{\end{list}}

%%%%%%%%%%%%%%%%%%%%%%%%%%%%%% User specified LaTeX commands.
\usepackage{numerica}

\usepackage{upquote}
\newcommand\rel{\,\varrho\;}
\DeclareMathOperator{\erf}{erf}
\DeclareMathOperator{\gd}{gd}

\reuse{}

\makeatother

\begin{document}
\title{\texttt{numerica}}
\author{Andrew Parsloe\\
(\url{ajparsloe@gmail.com})\\
}
\maketitle
\begin{abstract}
The \verb`numerica` package defines a command to numerically evaluate
mathematical expressions in the LaTeX form in which they are typeset.
For programs like LyX with a preview facility, or compile-as-you-go
systems, interactive back-of-envelope calculations and numerical exploration
are possible within the document being worked on. The package requires
the bundles \verb`l3kernel` and \verb`l3packages`, and the \verb`amsmath`
and \verb`mathtools` packages. \\
\\
\noindent\begin{minipage}[t]{1\columnwidth}%
\begin{shaded}%

\paragraph*{Note:}
\begin{itemize}
\item {\normalsize This document applies to version 2.0.0 of }{\normalsize\texttt{numerica.sty}}{\normalsize .}{\small\par}
\item {\normalsize Reasonably recent versions of the \LaTeX 3 bundles }{\normalsize\texttt{l3kernel}}{\normalsize{}
and }{\normalsize\texttt{l3packages}}{\normalsize{} are required (although
much of }{\normalsize\verb`l3kernel`}{\normalsize{} has
been incorporated into \LaTeXe{} since February 2020).}{\small\par}
\item {\normalsize The package requires }{\normalsize\texttt{amsmath}}{\normalsize{}
and }{\normalsize\texttt{mathtools}}{\normalsize .}{\small\par}
\item {\normalsize I refer many times in this document (especially §\ref{sec:Argument-parsing})
to }{\normalsize\emph{Handbook of Mathematical Functions}}{\normalsize ,
edited by Milton Abramowitz and Irene A. Segun, Dover, 1965. This
is abbreviated to }{\normalsize\emph{HMF}}{\normalsize , often followed
by a number like 1.2.3 to locate the actual expression referenced.}{\small\par}
\item {\normalsize Version 2.0.0 of }{\normalsize\texttt{numerica}}{\small\par}
\begin{itemize}
\item {\normalsize splits into distinct packages the additional functionality
previously available with the }{\normalsize\texttt{plus}}{\normalsize{}
and }{\normalsize\texttt{tables}}{\normalsize{} package options of version
1;}{\small\par}
\item {\normalsize allows for user-defined macros and constants (with the
}{\normalsize\texttt{\textbackslash nmcMacros }}{\normalsize and }{\normalsize\texttt{\textbackslash nmcConstants}}{\normalsize{}
commands) to be used in expressions to be evaluated;}{\small\par}
\item {\normalsize rewrites the code and changes the behaviour of }{\normalsize\texttt{\textbackslash nmcReuse}}{\normalsize{}
to maintain uniformity across all commands (}{\normalsize\texttt{\textbackslash nmcEvaluate}}{\normalsize ,
}{\normalsize\texttt{\textbackslash nmcInfo}}{\normalsize , }{\normalsize\texttt{\textbackslash nmcMacros}}{\normalsize ,
}{\normalsize\texttt{\textbackslash nmcConstants}}{\normalsize , }{\normalsize\texttt{\textbackslash nmcReuse}}{\normalsize );
this command is no longer compatible with its use in v.1;}{\small\par}
\item {\normalsize changes the behaviour of }{\normalsize\texttt{\textbackslash text}}{\normalsize{}
and }{\normalsize\texttt{\textbackslash mbox}}{\normalsize{} in the
}{\normalsize\texttt{\textbackslash eval}}{\normalsize{} command; adds
}{\normalsize\texttt{\textbackslash textrm}}{\normalsize , }{\normalsize\texttt{\textbackslash textsf}}{\normalsize ,
and }{\normalsize\texttt{\textbackslash texttt}}{\normalsize{} to compensate;}{\small\par}
\item {\normalsize includes many adjustments to the code, including around
nesting of commands;}{\small\par}
\item {\normalsize adds to and amends documentation.}{\small\par}
\end{itemize}
\end{itemize}
\end{shaded}%
\end{minipage}
\end{abstract}
\begin{center}
\tableofcontents{}
\par\end{center}

\chapter{Introduction}

\texttt{numerica} is a \LaTeX{} package offering the ability to numerically
evaluate mathematical expressions in the \LaTeX{} form in which they
are typeset.\footnote{\texttt{numerica} evolved from the author's \texttt{calculyx} package
that was designed for use with the document processor \LyX{} (and available
for download from a link on the \LyX{} wiki website but not CTAN). }

There are a number of packages which can do calculations in \LaTeX ,\footnote{A simple search finds the venerable \texttt{calc} in the \LaTeX{} base,
\texttt{calculator }(including an associated \texttt{calculus} package),
\texttt{fltpoint}, \texttt{fp} (\emph{fixed} rather than floating
point), \texttt{spreadtab} (using either \texttt{fp} or \texttt{l3fp}
as its calculational engine) if you want simple spreadsheeting with
your calculations, the elaborate \texttt{xint}, \texttt{pst-calculate}
(a limited interface to \texttt{l3fp}), \texttt{l3fp} in the \LaTeX 3
kernel, and \texttt{xfp}, the \LaTeX 3 interface to \texttt{l3fp}.
Other packages include a calculational element but are restricted
in their scope. (\texttt{longdivision} for instance is elegant, but
limited only to long division.) } but those I am aware of all require the mathematical expressions
they operate on to be changed to an appropriate syntax. Of these packages
\texttt{xfp} comes closest to my objective with \texttt{numerica}.
For instance, given a formula
\begin{centred}
\verb`\frac{\sin (3.5)}{2} + 2\cdot 10^{-3}`
\end{centred}
(in a math environment), this can be evaluated using \texttt{xfp}
by transforming the expression to \verb`sin(3.5)/2 + 2e-3` and wrapping
this in the command \verb`\fpeval`. In \texttt{numerica} you don't
need to transform the formula, just wrap it in an \verb`\eval` command:
\begin{centred}
\verb`\eval{ \frac{\sin (3.5)}{2} + 2\cdot 10^{-3} }`.
\end{centred}
(for the acutal calculation see §\ref{subsec:introSimple-examples}).\texttt{ }

\texttt{numerica}, like \texttt{xfp} and a number of other packages,
uses \texttt{l3fp} (the \LaTeX 3 floating point module in \texttt{l3kernel})
as its calculational engine. To some extent the main command, \verb`\nmcEvaluate`,
short-name form \verb`\eval`, is a pre-processor to \texttt{l3fp},
converting mathematical expressions written in the \LaTeX\texttt{
}form in which they will be typeset into an `fp-ified' form that
is digestible by \texttt{l3fp}. The aim is to make the command act
as a wrapper around such formulas. Ideally, one should not have to
make \emph{any} adjustment to them, although any text on Fourier series
suggests that hope in full generality is delusional. Surprisingly
often however it \emph{is} possible. We shall see shortly that even
complicated formulas like

\[
\cos\tfrac{m}{n}\pi-(1-4\sin^{2}\tfrac{m}{3n}\pi)\frac{\sin\tfrac{1}{n}\pi\sin\tfrac{m-1}{n}\pi}{2\sin^{2}\tfrac{m}{3n}\pi},
\]
and
\[
\left(\frac{1-4\sin^{2}\tfrac{m}{3n}\pi}{2\sin^{2}\tfrac{m}{3n}\pi}\right)\sin\tfrac{2m-3}{3n}\pi\sin\tfrac{m-3}{3n}\pi,
\]

\noindent can be evaluated `as is' (see below, §\ref{subsec:introReassurance}).
There is no need to shift the position of the superscript $2$ on
the sines, no need to parenthesize the arguments of $\sin$ and $\cos$,
no need to insert asterisks to indicate multiplication, no need to
change the \verb`\frac` and \verb`\tfrac`-s to slashes, \texttt{/},
no need to delete the \verb`\left` and \verb`\right` that qualify
the big parentheses (in the underlying \LaTeX ) in the second expression.
Of course, if there are variables in an expression, as in these examples,
they will need to be assigned values. And how the result of the evaluation
is presented also requires specifying, but the aim is always: to evaluate
mathematical expressions in \LaTeX{} with as little adjustment as possible
to the form in which they are typeset.

\texttt{numerica} is written in \texttt{expl3}, the programming language
of the \LaTeX 3 project. It uses the \LaTeX 3 module \texttt{l3fp}
(part of \texttt{l3kernel}) as its calculational engine. This enables
floating point operations to 16 significant figures, with exponents
ranging between $-10000$ and $+10000$. Many functions and operations
are built into \texttt{l3fp} – arithmetic operations, trigonometric,
exponential and logarithm functions, factorials, absolute value, max
and min. Others have been constructed for \texttt{numerica }from \texttt{l3fp}
ingredients – binomial coefficients, hyperbolic functions, sums and
products – but to the user there should be no discernible difference. 

Associated packages provide for additional operations: iteration of
functions, finding zeros of functions, recurrence relations, mathematical
table building; others are planned (e.g. calculus).

\section{How to use \texttt{numerica}}

The package is invoked in the usual way: put
\begin{lyxcode}
\textbackslash usepackage\{numerica\}
\end{lyxcode}
\noindent in the \LaTeX{} preamble. \texttt{numerica} requires the
\texttt{amsmath} and \texttt{mathtools} packages and loads these automatically.
\texttt{numerica} will also accept use of some relational symbols
from the \texttt{amssymb} package provided that package is loaded
by the user; see §\ref{subsec:evalBoolean-output}. 

\subsection{Package options}

\label{subsec:introPackagesOptions}Currently there are none. With
version 2.0.0 of \verb`numerica` a change has been made to how additional
functionality for the package is invoked; see §\ref{subsec:Related-packages}
below. This means that the options available in version 1 have been
discontinued.

\subsection{Associated packages}

\label{subsec:Related-packages}In version 1 of \verb`numerica` some
additional functionality for the package could be gained by specifying
package options – for instance the ability to create tables of function
values or to iterate or find fixed points of functions. However this
manner of invoking the addtional functionality makes the maintaining
of semantic version numbering across the whole \verb`numerica` package
difficult. With version 2.0.0 of the package, the addtional functionality
has been separated into separate \LaTeX{} packages. Currently there
are two of these, \verb`numerica-plus` and \verb`numerica-tables`.
They are loaded with the familiar \verb`\usepackage` command in the
document preamble and require the availability of the \verb`numerica`
package in your \TeX{} distribution. Neither package requires a \verb`\usepackage{numerica}`
statement; they take care of that themselves. So, if you enter
\begin{lyxcode}
\textbackslash usepackage\{numerica-plus\}
\end{lyxcode}
in the preamble of your document you gain access not only to the commands
in the \verb`numerica` package but also to the commands \verb`\nmcIterate`,
\verb`\nmcSolve`, and \verb`\nmcRecur`. \verb`\nmcIterate` enables
the iteration of functions of a single variable, including finding
fixed points. \verb`\nmcSolve` enables the solving of equations of
the form $f(x)=0$ (i.e. finding the zeros of $f$), or the finding
of local maxima or minima of a function of one variable. \verb`\nmcRecur`
enables the calculation of terms in recurrence relations, like the
terms of the Fibonacci series, or othogonal polynomials defined recurrently.
See the associated document \verb`numerica-plus.pdf` for details.

If you enter
\begin{lyxcode}
\textbackslash usepackage\{numerica-tables\}
\end{lyxcode}
in the preamble of your document you gain access not only to the commands
in the \verb`numerica` package but also to the command \verb`\nmcTabulate`
which enables the creation of (possibly multi-column) tables of function
values and makes available most of the table formats evident in \emph{HMF}.
See the associated document \verb`numerica-tables.pdf` for details. 

A package \verb`numerica-calculus` is currently being developed.

\subsection{Simple examples of use}

\label{subsec:introSimple-examples}A simple example of use is provided
by the document
\begin{verbatim}
    \documentclass{article} 
    \usepackage{numerica}
    \begin{document}

        \eval{$ mc^2 $}[m=70,c=299 792 458][8x]

    \end{document}

\end{verbatim}
\noindent We have a formula between math delimiters: \verb`$ mc^2 $`.
We have wrapped a command \verb`\eval` around the lot, added an optional
argument in parentheses specifying numericaal values for the quantities
\texttt{m} and \texttt{c}, and concluded it all with a trailing optional
argument specifying that the result should be presented to $8$ places
of decimals and in scientific notation (the \verb`x`). Running \texttt{pdflatex}
on this document generates a pdf displaying
\begin{centred}
\eval{$ mc^2 $}[m=70,c=299 792 458][8x]
\end{centred}
\noindent where the formula ($mc^{2})$ is equated to the numerical
value resulting from substituting the given values of $m$ and $c$.
Those values are displayed in a list following the result. The calculation
is presented to $8$ decimal places in scientific notation. (According
to Einstein's famous equation $E=mc^{2}$ this is the enormous energy
content, in joules, of what was once considered an average adult Caucasian
male. Only a minute fraction is ever available.) 

A second example is provided by the formula in earlier remarks:
\begin{verbatim}
    \documentclass{article} 
    \usepackage{numerica}
    \begin{document}

        \eval{\[ \frac{\sin(3.5)}{2} + 2\cdot 10^{-3} \]}

    \end{document}

\end{verbatim}
Running \texttt{pdflatex} on this document produces the result

\eval{\[ \frac{\sin(3.5)}{2} + 2\cdot 10^{-3} \]}

The \verb`\eval` command used in these examples is the main command
of the \texttt{numerica} package. It is discussed in full in the next
chapter, but first some preliminaries.

\subsection{Display of the result}

\label{subsec:introResultDisplay}In what follows I shall write things
like (but generally more complicated than)
\begin{centred}
\verb`$ \eval{ 1+1 } $` $\Longrightarrow \eval{ 1 + 1 } $
\end{centred}
to mean: run \texttt{pdflatex} on a document containing \verb`\eval{1+1}`
in the document body to generate a pdf containing the calculated result
($2$ in this instance). The reader will note that I have used dollar
signs to delimit the math environment. I could (and perhaps should)
have used the more \LaTeX -pure \verb`\( \)`, which will do equally
well, but habit has won out. In the example the \verb`\eval` command
is used \emph{within} a math environment (delimited by the dollar
signs). It is not limited to this behaviour. The command can also
wrap \emph{around} the math delimiters (as we saw in the previous
examples):
\begin{centred}
\verb`\eval{$ 1+1 $}` $\Longrightarrow$ \eval{$ 1+1 $}.
\end{centred}
As you can see, the display that results is different. 
\begin{itemize}
\item When the \verb`\eval` command is used\emph{ within} a math environment,
only the \emph{result,} followed possibly by the \emph{variable~=~value
list} (see §\ref{subsec:evalVv-list}) is displayed.
\end{itemize}
Environments may include the various AMS environments as well as the
standard \LaTeX{} inline ( \verb`$ $` or \verb`\( \)` ), \texttt{equation}
( \verb`\[ \]` ) and \texttt{eqnarray} environments. For an example
of \verb`\eval` within an \verb`align*` environment see §\ref{subsec:introExploring}
below. 
\begin{itemize}
\item When the \verb`\eval` command is wrapped\emph{ around} a math environment,
the result is displayed in the form, \emph{formula = result} (followed
possibly by the \emph{variable~=~value list}) within that environment,
\begin{itemize}
\item If the formula is long or contains many variables then it may be desirable
to split the display over two lines; see §\ref{subsec:evalChanging-display-format}
and §\ref{subsec:settings New-line-display},
\end{itemize}
\end{itemize}
the whole presented as an inline expression if \verb`$` delimiters
are used, or as a display-style expression otherwise. (See the $mc^{2}$
example for an illustration.) 

It is not clear to me that wrapping \verb`\eval` \emph{around} the
AMS environments, except for \texttt{multline}, makes much sense,
although it can be done. Here is an example of \verb`\eval` wrapped
around a \texttt{multline{*}} environment (the phantom is there so
that the hanging $+$ sign spaces correctly),\medskip{}

\begin{minipage}[t]{0.8\columnwidth}%
\begin{verbatim}
\eval{ \begin{multline*}
          1+2+3+4+5+6+7+8+9+10+\phantom{0}\\
            11+12+13+14+15+16+17+18+19 
       \end{multline*} }
\end{verbatim}
%
\end{minipage}

\medskip{}$\Longrightarrow$ %
\begin{minipage}[t]{0.8\columnwidth}%
\vspace{-2ex}\eval{ \begin{multline*}       
		1+2+3+4+5+6+7+8+9+10+\phantom{0}\\       		  11+12+13+14+15+16+17+18+19      
\end{multline*} }%
\end{minipage}
\begin{itemize}
\item It is also possible to dispense with math delimiters entirely, neither
wrapped within nor wrapped around the \verb`\eval` command, but in
that case \texttt{numerica} acts as if \verb`\eval` had been used
within \verb`\[` and \verb`\]` and displays the result accordingly.
\end{itemize}

\subsection{Checking}

A question I found on the internet that caught my attention was to
simplify $\sqrt{220-30\sqrt{35}}$. I found myself intrigued. After
some bumbling and fumbling, I let 
\[
x=\sqrt{220-30\sqrt{35}},\qquad y=\sqrt{220+30\sqrt{35}},
\]
(which seems an obvious thiing to do). Then 
\[
xy=10\sqrt{484-315}=10\sqrt{13^{2}}=130.
\]
Since $x^{2}+y^{2}=440$ it was then easy to form both $(x+y)^{2}$
and $(x-y)^{2}$ and by separating the resulting numbers into their
prime factors, to work out that $x=5\sqrt{7}-3\sqrt{5}$. Was I right? 
\begin{centred}
\verb`\eval{$ \sqrt{220-30\sqrt{35}} $}` $\Longrightarrow$ \eval{$ \sqrt{220-30\sqrt{35}} $},

\verb`\eval{$ 5\sqrt{7}-3\sqrt{5} $}` $\Longrightarrow$ \eval{$ 5\sqrt{7}-3\sqrt{5} $}.
\end{centred}
Yes, the simplification was correct. And indeed $y=5\sqrt{7}+3\sqrt{5}$:
\begin{centred}
\verb`\eval{$ \sqrt{220+30\sqrt{35}} $}` $\Longrightarrow$ \eval{$ \sqrt{220+30\sqrt{35}} $},

\verb`\eval{$ 5\sqrt{7}+3\sqrt{5} $}` $\Longrightarrow$ \eval{$ 5\sqrt{7}+3\sqrt{5} $}.
\end{centred}
As a final flourish,
\begin{verbatim}
    \eval{$ xy $}
       [ x=5\sqrt{7}-3\sqrt{5}, 
         y=5\sqrt{7}+3\sqrt{5} ]
\end{verbatim}
$\Longrightarrow$     \eval{$ xy $}
       [x=5\sqrt{7}-3\sqrt{5}, 
        y=5\sqrt{7}+3\sqrt{5}].

\subsection{Exploring}

\label{subsec:introExploring}When working on \texttt{numerica}'s
predecessor package, I constantly tested it against known results
to check for coding errors. One test was to ensure that 
\[
\left(1+\frac{1}{n}\right)^{n}
\]
did indeed converge to the number $e$ as $n$ increased.\texttt{ }Let's
do that here. Try first $n=10$:
\begin{center}
\verb`\eval{$ e-(1+1/n)^n $}[n=10][x]` $\Longrightarrow$ \eval{$ e-(1+1/n)^n $}[n=10][x].
\par\end{center}

\noindent (The default number of decimal places displayed is $6$.)
The difference between $e$ and $(1+1/n)^{n}$ is about an eighth
($0.125$) when $n=10$, which is encouraging but hardly decisive.
The obvious thing to do is increase the value of $n$. I'll use an
\verb`align*` environment to `prettify' the presentation of the
results:
\begin{verbatim}
  \begin{align*}
    e-(1+1/n)^{n} & =\eval{e-(1+1/n)^n}[n=1\times10^5][*x],\\
    e-(1+1/n)^{n} & =\eval{e-(1+1/n)^n}[n=1\times10^6][*x],\\
    e-(1+1/n)^{n} & =\eval{e-(1+1/n)^n}[n=1\times10^7][*x],\\
    e-(1+1/n)^{n} & =\eval{e-(1+1/n)^n}[n=1\times10^8][*x].
  \end{align*}
\end{verbatim}
(most of which was written using copy and paste) which produces
\begin{align*}
e-(1+1/n)^{n} & =\eval{e-(1+1/n)^{n}}[n=1\times10^{5}][x],\\
e-(1+1/n)^{n} & =\eval{e-(1+1/n)^{n}}[n=1\times10^{6}][*x],\\
e-(1+1/n)^{n} & =\eval{e-(1+1/n)^{n}}[n=1\times10^{7}][x],\\
e-(1+1/n)^{n} & =\eval{e-(1+1/n)^{n}}[n=1\times10^{8}][x].
\end{align*}
Clearly $(1+1/n)^{n}$ converges to $e$, the difference between them
being of order $1/n$, but that is not what catches the eye. There
is an unanticipated regularity here. 1.35914? Double the number: \verb`$\eval{2\times 1.35914}[5]$`\texttt{
}$\Longrightarrow\eval{2\times1.35914}[5]$ which is close enough
to $e$ to suggest a relationship, namely,
\[
\lim_{n\to\infty}n\left(e-\left(1+\frac{1}{n}\right)^{n}\right)=\tfrac{1}{2}e.
\]

This was new to me. Is it true? From the familiar expansion of the
logarithm

\begin{align*}
\ln\left(1+\frac{1}{n}\right)^{n} & =n\ln\left(1+\frac{1}{n}\right)\\
 & =n\left(\frac{1}{n}-\frac{1}{2}\frac{1}{n^{2}}+\frac{1}{3}\frac{1}{n^{3}}-\ldots\right)\\
 & =1-\frac{1}{2n}\left(1-\frac{2}{3}\frac{1}{n}+\frac{2}{4}\frac{1}{n^{2}}-\right)\\
 & \equiv1-\frac{1}{2n}E_{n},
\end{align*}
say. Since $E_{n}$ is an alternating series and the magnitudes of
the terms of the series tend to $0$ monotonically, $1>E_{n}>1-2/3n$.
From this and the inequalities $1/(1-x)>e^{x}>1+x$ when $x<1$ it
proved a straightforward matter to verify the proposed limit.

\subsection{Reassuring}

\label{subsec:introReassurance}In the course of some hobbyist investigations
in plane hyperbolic geometry I derived the formula
\[
\Phi_{1}(m,n)=\cos\tfrac{m}{n}\pi-(1-4\sin^{2}\tfrac{m}{3n}\pi)\frac{\sin\tfrac{1}{n}\pi\sin\tfrac{m-1}{n}\pi}{2\sin^{2}\tfrac{m}{3n}\pi},
\]
for $m=2,3,\ldots$ and integral $n\ge2m+1$. A key concern was: when
is $\Phi_{1}$ positive? After an embarrassingly laborious struggle,
I managed to work this expression into the form
\[
\Phi_{2}(m,n)=\left(\frac{1-4\sin^{2}\tfrac{m}{3n}\pi}{2\sin^{2}\tfrac{m}{3n}\pi}\right)\sin\tfrac{2m-3}{3n}\pi\sin\tfrac{m-3}{3n}\pi,
\]
in which the conditions for positivity are clear: with $n\ge2m+1$,
so that $m\pi/3n<\pi/6$, the first factor is always positive, the
second is positive for $m\ge2$, and the third is positive for $m\ge4$.
All well and good, but given the struggle to derive $\Phi_{2}$, was
I confident that $\Phi_{1}$ and $\Phi_{2}$ really\emph{ }are equal?
It felt all too likely that I had made a mistake. 

The simplest way to check was to see if the two expressions gave the
same numerical answers for a number of $m,\thinspace n$ values. I
wrote \verb`\eval{\[ \]}[m=2,n=5]` twice and between the delimiters
pasted the already composed expressions for $\Phi_{1}$ and $\Phi_{2}$,
namely: 
\begin{verbatim}
    \eval{\[ 
            \cos\tfrac{m}{n}\pi-(1-4\sin^{2}\tfrac{m}{3n}\pi)
            \frac{\sin\tfrac{1}{n}\pi\sin\tfrac{m-1}{n}\pi}
            {2\sin^{2}\tfrac{m}{3n}\pi} 
          \]}[m=2,n=5]
    \eval{\[ 
            \left(
              \frac{1-4\sin^{2}\tfrac{m}{3n}\pi}
              {2\sin^{2}\tfrac{m}{3n}\pi}
            \right)
            \sin\tfrac{2m-3}{3n}\pi\sin\tfrac{m-3}{3n}\pi
          \]}[m=2,n=5]
\end{verbatim}
I have added some formatting – indenting, line breaks – to make the
formulas more readable for the present document but otherwise left
them unaltered. The \verb`\eval` command can be used for even quite
complicated expressions without needing to tinker with their \LaTeX{}
form, but you may wish – as here – to adjust white space to clarify
the component parts of the formula. Running \texttt{pdflatex} on these
expressions, the results were

\eval{\[ 
        \cos\tfrac{m}{n}\pi-(1-4\sin^{2}\tfrac{m}{3n}\pi)
        \frac{\sin\tfrac{1}{n}\pi\sin\tfrac{m-1}{n}\pi}
        {2\sin^{2}\tfrac{m}{3n}\pi} 
      \]}[m=2,n=5]

\eval{\[ 
        \left(
          \frac{1-4\sin^{2}\tfrac{m}{3n}\pi}
          {2\sin^{2}\tfrac{m}{3n}\pi}
        \right)
       \sin\tfrac{2m-3}{3n}\pi\sin\tfrac{m-3}{3n}\pi
      \]}[m=2,n=5]

\noindent which was reassuring. Doing it again but with different
values of $m$ and $n$, again the results coincided:

\eval{\[ 
        \cos\tfrac{m}{n}\pi-(1-4\sin^{2}\tfrac{m}{3n}\pi)
        \frac{\sin\tfrac{1}{n}\pi\sin\tfrac{m-1}{n}\pi}
        {2\sin^{2}\tfrac{m}{3n}\pi} 
      \]}[m=5,n=13]

\eval{\[ 
        \left(
          \frac{1-4\sin^{2}\tfrac{m}{3n}\pi}
          {2\sin^{2}\tfrac{m}{3n}\pi}
        \right)
       \sin\tfrac{2m-3}{3n}\pi\sin\tfrac{m-3}{3n}\pi
      \]}[m=5,n=13]

\noindent Thus reassured that there was \emph{not }an error in my
laborious derivation of $\Phi_{2}$ from $\Phi_{1}$, it was not difficult
to work back from $\Phi_{2}$ to $\Phi_{1}$ then reverse the argument
to find a straightforward derivation. 

\chapter{\texttt{\textbackslash nmcEvaluate} (\texttt{\textbackslash eval)}}

The main calculational command in \texttt{numerica} is \verb`\nmcEvaluate`.
Unlike some other commands which are loaded optionally, \verb`\nmcEvaluate`
is \emph{always} loaded, and therefore always available. Because \verb`\nmcEvaluate`
would be tiresome to write too frequently,\texttt{ }particularly for
back-of-envelope calculations, there is an equivalent short-name form,
\verb`\eval`, used almost exclusively in the following. But note:
wherever you see the command \verb`\eval`, you can substitute \verb`\nmcEvaluate`
and obtain the same result.

\verb`\eval` (like other short-name forms of other commands in the
\texttt{numerica} suite) is defined using \verb`\ProvideDocumentCommand`
from the \texttt{xparse} package. Hence if \verb`\eval` has already
been defined in some other package already loaded, it will not be
redefined by \texttt{numerica}. It will retain its meaning in the
other package. Its consequent absence from \texttt{numerica} may be
an irritant, but only that; \verb`\nmcEvaluate` is defined using
\texttt{xparse}'s \verb`\DeclareDocumentCommand` which would override
any (freakishly unlikely) previous definition of \verb`\nmcEvaluate`
in another package and would therefore still be available.

\section{Syntax of \texttt{\textbackslash nmcEvaluate (\textbackslash eval)}}

There are five arguments to the \verb`\nmcEvaluate` (or \verb`\eval`)
command, of which only one, the third, is mandatory. All others are
optional. If all are deployed the command looks like
\begin{centred}
\noindent \verb`\nmcEvaluate*[settings]{expr.}[vv-list][num. format]`
\end{centred}
I discuss the various arguments in the referenced sections.
\begin{enumerate}
\item \verb`*` optional switch; if present ensures display of only the
numerical result (suppresses display of the formula and vv-list);
see §\ref{subsec:evalVvSuppresList}
\item \verb`[settings]` optional comma-separated list of \emph{key=value
}settings for this particular calculation; see §\ref{sec:settingsOption}
\item \verb`{expr.}` the only mandatory argument; the mathematical expression/formula
in \LaTeX{} form that is to be evaluated 
\item \verb`[vv-list]` optional comma-separated list of \emph{variable=value
}items; see §\ref{subsec:evalVv-list}
\item \verb`[num. format]` optional format specification for presentation
of the numerical result (rounding, padding with zeros, scientific
notation, boolean output); see~§\ref{subsec:evalRoundingEtc}
\end{enumerate}
Note that arguments 4 and 5 are both square-bracket delimited optional
arguments. Should only one such argument be used, \texttt{numerica}
determines which is intended by looking for an equals sign within
the argument. Its presence indicates the argument is the vv-list;
its absence indicates the argument is the number format specification. 

The vv-list and number-format specification are \emph{trailing} optional
arguments. They do not need to be hard against their preceding arguments;
intervening spaces are allowed. This means there is a possibility
that should the \verb`\eval` command be followed by a square-bracketed
mathematical expression \texttt{numerica} might confuse it with one
of its trailing arguments. Experience using \texttt{numerica} suggests
that this will be a (very) rare occurrence and is easily prevented
by inserting an empty brace pair (\verb`{}`) before the offending
square-bracketed expression. Allowing spaces between the arguments
enables complicated expressions and large vv-lists to be formatted
with new lines and white space to aid clarity – without requiring
the insertion of comment characters (\verb`%`).

Recommended practice is to minimise the number of optional arguments
used in \LaTeX{} commands by consolidating such arguments into a single
\emph{key=value} list. Although \texttt{numerica} uses such an argument,
the vv-list does not fit naturally into that scheme. And practice
suggests that separating out the elements of the number format specification
(rounding value, padding with zeros, scientific notation, boolean
output) and placing them in a trailing argument is both convenient
and intuitive for the kind of back-of-envelope calculations envisaged
for \texttt{numerica}. 

\section{The variable=value list}

\label{subsec:evalVv-list}To evaluate algebraic, trigonometric and
other formulas that involve \emph{variables} we need to give those
variables values. This is done in the \emph{variable=value list} –
or \emph{vv-list} for short. This is the fourth argument of the \texttt{\textbackslash nmcEvaluate}
command and is a square-bracket delimited optional argument (optional
because an expression may depend only on constants and numbers).

\subsection{Variable names}

\label{subsec:evalVariableNames}In mathematical practice, variable
names are generally single letters of the Roman or Greek alphabets,
sometimes also from other alphabets, in a variety of fonts, and often
with subscripts or primes or other decorations. In \texttt{numerica}
a variable name is \emph{what lies to the left of the equals sign
in an item} of the vv-list. Thus variables can be multi-token affairs:
$x',x'',x^{iv},x_{n},x'_{n},x''_{mn}$, $^{k}C_{n},var,\mathrm{var},Fred,\mathbf{Fred},\mathcal{FRED}\ldots$
(This criterion for what makes a variable name means a name may contain
spaces – for instance \verb`x x` should not cause a \verb`numerica`
error – but such names are not part of mathematical practice.) Usually,
for the kind of back-of-envelope calculations envisaged for \verb`numerica`,
and for ease of typing, most variables will be single letters from
the Roman or Greek alphabets. 

Because equals signs and commas give structure to the vv-list, it
should also be clear that a variable name should not contain a \emph{naked}
equals sign or a \emph{naked} comma. They can be incorporated in a
variable name but only when decently wrapped in braces, like \verb`R_{=}`
displaying as $R_{=}$ or \verb`X_{,i}` displaying as $X_{,i}$. 

Note that $x$ and $\mathrm{x}$ will be treated by \verb`numerica`
as \emph{different} variables since, in the underlying \LaTeX , one
is \texttt{x} and the other \verb`\mathrm{x}`. Even names that look
identical in the pdf may well be distinct in \LaTeX . This is true
particularly of superscripts and subscripts: \verb`x_0` and \verb`x_{0}`
appear identical in the pdf but in the underlying \LaTeX{} they are
distinct, and will be treated as distinct variables by \verb`numerica`. 

Although multi-token variables are perfectly acceptable, \emph{internally}
\verb`numerica` works with single tokens. Variable names can be so
different in structure, one from another, that to ease the parsing
of formulas, all \emph{internal} variable names are assumed to be
single tokens. Hence a necessary initial step for the package is to
map all multi-token variable names in the vv-list and the formula
to single tokens. \verb`numerica` does this by turning the multi-token
variable names into control sequences with names in the sequence \verb`\nmc_a`,
\verb`\nmc_b`, \verb`\nmc_c`, etc., then searches through the vv-list
and the formula for every occurrence of the multi-token names and
replaces them with the relevant control sequences. It does this in
order of decreasing size of name, working from the names that contain
most tokens down to names containing only two tokens. (Doing the replacing
in this order prevents \emph{parts} of longer names possibly being
mistaken for shorter variable names.) 

The conversion process uses computer resources. Even if there are
no multi-token variables present, \verb`numerica` still needs to
check that this is so – unless the user alerts the program to the
fact. This can be done by making a brief entry \texttt{xx=0 }in the
settings option (the second optional argument of \verb`\nmcEvaluate`);
see §\ref{subsec:settingsMultitokSwitch}. If the user never (or hardly
ever) uses multi-token variables, then a more permanent solution is
to create a file \verb`numerica.cfg`\texttt{ }with the line \texttt{multitoken-variables
= false}; see §\ref{sec:settingsDefaults} for this.

\subsection{The vv-list and its use}

A vv-list is a comma-separated list where each item is of the form
\emph{variable=value}. It might be something simple like
\begin{lyxcode}
{[}g=9.81,t=2{]}
\end{lyxcode}
or something more complicated like
\begin{lyxcode}
{[}V\_S=\textbackslash tfrac43\textbackslash pi~r\textasciicircum 3,V\_C=2\textbackslash pi~r\textasciicircum 2h,h=3/2,r=2{]}.
\end{lyxcode}
Spaces around the equals signs or the commas are stripped away during
processing so that 
\begin{lyxcode}
{[}g=9.81,t=2{]}~\textrm{and}~{[}~g~=~9.81~,~t~=~2{]}
\end{lyxcode}
are the \emph{same} variable=value list.

\subsubsection{Evaluation from right to left}

In these examples, with variables depending on other variables, there
is an implication: that the list is evaluated \emph{from the right}.
Recall how a function of a function is evaluated, say\emph{ }$y=f(g(h(x)))$.
To evaluate $y$, first $x$ is assigned a value then $h(x)$ is calculated,
then $g(h(x))$ then $f(g(h(x)))=y$. We work from right to left,
from the innermost to the outermost element. Or consider an example
like calculating the area of a triangle by means of the formula
\[
A=\sqrt{s(s-a)(s-b)(s-c)}.
\]
First we write the formula; then we state how $s$ depends on $a,b,c$,
namely $s=\frac{1}{2}(a+b+c)$, then we give values to $a,b,c$. In
\texttt{numerica} this is mirrored in the layout of the \verb`\eval`
command:
\begin{verbatim}
    \eval{$ \sqrt{s(s-a)(s-b)(s-c)} $}
                 [s=\tfrac12(a+b+c),a=3,b=4,c=5]
\end{verbatim}
The formula in a sense is the leftmost extension of the vv-list. The
entire evaluation occurs from right to left.

This means that the rightmost variable in the vv-list can depend only
on constants and numbers – although it may be a complicated expression
of those elements. Other variables in the vv-list can depend on variables
\emph{to their right} but not to their left.

\subsubsection{Expressions in the variable=value list}

Suppose our expression is $\tfrac{4}{3}\pi r^{3}$, the volume $V_{S}$
of a sphere in terms of its radius $r$, and we want to calculate
the volume for different values of $r$ to get a sense of how rapidly
volume increases with radius.
\begin{centred}
\verb`$ V_S=\eval{ \tfrac43\pi r^3 }[r=1] $` $\Longrightarrow$ $ V_S=\eval{ \tfrac43\pi r^3 }[r=1] $.
\end{centred}
Having set up this calculation it is now an easy matter to change
the value of $r$ in the vv-list:
\begin{centred}
\verb`$ V_S=\eval{ \tfrac43\pi r^3 }[r=1.5] $` $\Longrightarrow$
$ V_S= \eval{ \tfrac43\pi r^3 }[r=1.5] $.

\verb`$ V_S=\eval{ \tfrac43\pi r^3 }[r=2] $` $  \Longrightarrow $ $V_S= \eval{ \tfrac43\pi r^3 }[r=2] $.
\end{centred}
To compute the volume $V_{C}=\pi r^{2}h$ of a cylinder, we have two
variables to assign values to:
\begin{centred}
\verb`$ V_C=\eval{ \pi r^2h }[h=4/3,r=1] $` $\Longrightarrow$ $ V_C=\eval{ \pi r^2h }[h=4/3,r=1] $.
\end{centred}
Although values in the vv-list are generally either numbers or simple
expressions (like \texttt{4/3}), that is not essential. A little more
complicated is
\begin{centred}
\verb`$ V_C=\eval{ hA_C }[A_C=\pi r^2,h=4/3,r=1] $` $\Longrightarrow$
$ V_C=\eval{ hA_C }[A_C=\pi r^2,h=4/3,r=1] $.
\end{centred}
where calculation of the volume of the cylinder has been split into
two: first calculate the area $A_{C}$ of its circular base and then,
once that has been effected, calculate the volume.

A second example is provided by Brahmagupta's formula for the area
of a triangle in terms of its semi-perimeter. In a triangle ABC, the
sides are $a=3$, $b=4$ and $c=5$. (Of course we know this is a
right-angled triangle with area $\tfrac{1}{2}ab=6$.) The semi-perimeter
$s=\tfrac{1}{2}(a+b+c)$ and the area of ABC is \medskip{}
\begin{verbatim}
    \eval{$ \sqrt{s(s-a)(s-b)(s-c) $}
                 [s=\tfrac12(a+b+c),a=3,b=4,c=5] 
\end{verbatim}
$\Longrightarrow$ \eval{$ \sqrt{s(s-a)(s-b)(s-c)} $}
    [s=\tfrac12(a+b+c),a=3,b=4,c=5].

\subsubsection{Constants}

\label{subsec:Built-in-Constants}\verb`numerica` has five built-in
constants and can also accept user-defined constants. For the latter,
see §\ref{sec:supplConstants}. The five built-in constants known
to \texttt{numerica} are \verb`\pi`, the ratio of circumference to
diameter of a circle; \verb`e`, the base of natural logarithms; Euler's
constant \verb`\gamma`, the limit of $\left(\sum_{1}^{N}1/n\right)-\ln N$
as $N\to\infty$; the golden ratio \verb`\phi`, equal to $\tfrac{1}{2}(1+\surd5)$;
and the utilitarian constant \verb`\deg`, the number of radians in
a degree.
\begin{centred}
\verb`\eval{$ \pi $}` $\Longrightarrow$ \eval{$ \pi $},

\verb`\eval{$ e $}` $\Longrightarrow$ \eval{$ e $},

\verb`\eval{$ \gamma $}` $\Longrightarrow$ \eval{$ \gamma $},

\verb`\eval{$ \phi $}` $\Longrightarrow$ \eval{$ \phi $},

\verb`\eval{$ \deg $}` $\Longrightarrow$ \eval{$ \deg $},
\end{centred}
so that \verb`\eval{$ 180\deg $}` $\Longrightarrow$ \eval{$ 180\deg $}
(as it should).

Let's combine two of these in a formula:
\begin{centred}
\verb`\eval{$ e^\pi-\pi^e $}` $\Longrightarrow$ \eval{$ e^\pi-\pi^e $},
\end{centred}
which is close-ish to $\tfrac{1}{4}e$: \verb`\eval{$ \tfrac14e $}`
$\Longrightarrow$ \eval{$ \tfrac14e $}.

\noindent\begin{minipage}[t]{1\columnwidth}%
\begin{shaded}%
In some contexts it may feel natural to use any or all of \verb`\pi`,
\verb`e`, \verb`\gamma` and \verb`\phi` as variables by assigning
values to them in the vv-list. \cprotect\texttt{numerica} does not
object. The values assigned in this way override the constants' values.
For example, if the triangle we labelled ABC previously was instead
labelled CDE then it has sides $c=3,d=4$ and (note!) $e=5$. It's
area therefore is\medskip{}
\begin{verbatim}
    \eval{$ \sqrt{s(s-c)(s-d)(s-e)} $}
              [s=\tfrac12(c+d+e),c=3,d=4,e=5]       
\end{verbatim}
$\Longrightarrow$ 

\eval{$ \sqrt{s(s-c)(s-d)(s-e)} $}
    [s=\tfrac12(c+d+e),c=3,d=4,e=5].\medskip{}

\noindent Since this is the correct area we see that \cprotect\texttt{e}
has been treated as a variable with the assigned value $5$, not as
the constant. But if \cprotect\texttt{e} (or \verb`\pi` or \verb`\gamma`
or \verb`\phi`) is not assigned a value in the vv-list then it has,
by default, the value of the constant. In the case of \cprotect\texttt{e},
if you wish to use it as a variable, the constant is always available
as \verb`\exp(1)`. There is no similar alternative available for
\verb`\pi`, \verb`\gamma` or \verb`\phi`. \end{shaded}%
\end{minipage}

\subsection{Display of the vv-list}

By default, the vv-list is displayed with (in fact following) the
numerical result. That and the format of the display can both be changed.

\subsubsection{Star option: suppressing display of the vv-list}

\label{subsec:evalVvSuppresList}If display of the vv-list is not
wanted at all, only the numerical result, it suffices to attach an
asterisk (star) to the \texttt{\textbackslash eval} command:
\begin{centred}
\verb`$ V_C=\eval*{ hA_C }[A_C=\pi r^2,h=4/3,r=1] $` $\Longrightarrow$
$ V_C=\eval*{ hA_C }[A_C=\pi r^2,h=4/3,r=1] $,
\end{centred}
or simply the naked result:
\begin{centred}
\verb`\eval*{$ hA_C $}[A_C=\pi r^2,h=4/3,r=1]` $\Longrightarrow$
\eval*{$ hA_C $}[A_C=\pi r^2,h=4/3,r=1].
\end{centred}
In the latter case, note that a negative result will display with
a hyphen for the minus sign unless you, the user, explicitly write
math delimiters around the \verb`\eval*` command as a whole. Wrapping
them around the formula has no effect:
\begin{centred}
\verb`\eval*{$ y $}[y=ax+b,x=2,a=-2,b=2]` $\Longrightarrow$ \eval*{$ y $}[y=ax+b,x=2,a=-2,b=2],

\verb`$ \eval*{ y }[y=ax+b,x=2,a=-2,b=2] $` $\Longrightarrow$ $ \eval*{ y }[y=ax+b,x=2,a=-2,b=2] $.
\end{centred}
The star option delivers a number as result, pure and simple.

\subsubsection{Suppressing display of items}

\label{subsec:evalVvSuppressVars}You may wish to retain some variables
in the vv-list display, but not all. For those variables you wish
omitted from the display, wrap each variable (but not the equals sign
or value) in braces. When calculating the volume of a cylinder in
the previous examples, the base area $A_{C}$ has a different status
from the `fundamental' variables $r$ and $h$. It is an intermediate
value, one that we pass through on the way to the final result. To
suppress it from display enclose the variable in braces:
\begin{centred}
\verb`$ V_C=\eval{ hA_C }[{A_C}=\pi r^2,h=4/3,r=1] $` $\Longrightarrow$
$ V_C=\eval{ hA_C }[{A_C}=\pi r^2,h=4/3,r=1] $.
\end{centred}
As you can see, $A_{C}$ no longer appears in the displayed vv-list.
Of course the name and its value are still recorded `behind the scenes'
and can still be used in calculations.

\subsubsection{Empty vv-list suppressed}

Should the vv-list be empty, or display of \emph{all} variables is
suppressed by wrapping each in braces, then \emph{nothing} is displayed
where the vv-list would normally be, not even any punctuation:
\begin{centred}
\verb`$ V_C=\eval{ hA_C }[{A_C}=\pi r^2,{h}=4/3,{r}=1] $` $\Longrightarrow$
$ V_C=\eval{ hA_C }[{A_C}=\pi r^2,{h}=4/3,{r}=1] $
\end{centred}
If you want a full stop after the result then you will need to add
it by hand or use the \verb`p=.` setting of §\ref{subsec:settingsPunctuation}.

\subsubsection{Changing the display format}

\label{subsec:evalChanging-display-format}In two examples above,
we have calculated the area of a triangle using Brahmagupta's formula.
Display of the result is crowded. Two remedies have just been suggested,
but a third one and preferable in this case would be to force display
of the vv-list and result to a new line. This can be done through
the settings option to the \verb`\eval` command, discussed in §\ref{subsec:settings New-line-display}.
However, if \verb`\eval` is wrapped around an \emph{appropriate}
environment (like \verb`multline`, but not \verb`equation`) it can
also be done simply by including \texttt{\textbackslash\textbackslash}
at the end of the formula. 

In the following example I use Brahmagupta's formula for calculating
the area of a cyclic quadrilateral (of which his formula for a triangle
is a special case). The cyclic quadrilateral in the example is formed
by a 45-45-90 triangle of hypotenuse 2 joined along the hypotenuse
to a 30-60-90 triangle. The sides are therefore $\surd2,\surd2,\surd3,1$.
Adding the areas of the two triangles, the area of the quadrilateral
is $A=1+\tfrac{1}{2}\surd3$, or in decimal form, \verb`$\eval{1+\tfrac12\surd3}$`
$\Longrightarrow$ $\eval{1+\tfrac12\surd3}$. Let's check with Brahmagupta's
formula:
\begin{verbatim}
    \eval{
      \begin{multline*}
        \sqrt{(s-a)(s-b)(s-c)(s-d)}\\
      \end{multline*} 
         }[s=\tfrac12(a+b+c+d),
           a=\surd2,b=\surd2,c=\surd3,d=1] 
\end{verbatim}
$\Longrightarrow$     \eval{
      \begin{multline*}
        \sqrt{(s-a)(s-b)(s-c)(s-d)}\\
      \end{multline*} 
         }[s=\tfrac12(a+b+c+d),
           a=\surd2,b=\surd2,c=\surd3,d=1]

\noindent\begin{minipage}[t]{1\columnwidth}%
\begin{shaded}%

\subsubsection{Abusing multi-token variable names}

\label{subsec:evalDon't-do-this!}A variable name is what lies to
the left of the equals sign of an item in the vv-list. Since multi-token
variables are converted to single tokens (like \verb`\nmc_a`) before
any calculating is done, it is possible to sin. Thus :
\begin{centred}
\verb`\eval{$ \sin\pi $}[{\sin\pi}=1]` $\Longrightarrow$ \eval{$ \sin\pi $}[{\sin\pi}=1];
\end{centred}
and (more?) egregiously,
\begin{centred}
\verb`\eval{$ 10 $}[{10}=20]` $\Longrightarrow$ \eval{$ 10 $}[{10}=20].
\end{centred}
What is happening here is that the multi-token `variables' \verb`\sin\pi`
and \verb`10` are being converted, right at the start of proceedings,
to single tokens like \verb`\nmc_a`, which in \TeX -speak are macros
containing their respective multiple tokens. For display purposes
they expand to those multiple tokens, but for calculating within \verb`numerica`
the single token is used. By this means one can easily create further
grotesqueries:
\begin{centred}
\verb`\eval{$ ++ + ++ $}[{++}=1]` $\Longrightarrow$ \eval{$ ++ + ++ $}[{++}=1],

\verb`\eval{$ 2(1+1) $}[{2(1}=3,{+1)}=5]` $\Longrightarrow$ \eval{$ 2(1+1) $}[{2(1}=3,{+1)}=5],

\verb`\eval{$ 1!! $}[{!!}=42]` $\Longrightarrow$ \eval{$ 1!! $}[{!!}=42].
\end{centred}
Should \verb`numerica` try to check variable names to avoid consequences
like this? I don't see any reasonable way of doing that. Symbols like
\verb`(` and \verb`+` can easily be part of valid variable names
– $k^{+},\,k^{-}$, $C_{n}^{(0)}$ and so on. It is left to the user,
in any \emph{public} document, to avoid such sins. (And they could
easily construct the displayed expressions in \LaTeX{} if they so wished
without recourse to \verb`\eval` at all.) See also §\ref{subsec:supplMacrosDisplay}
where a similar issue arises with user-defined macros.\end{shaded}%
\end{minipage}

\section{Formatting the numerical result}

\label{subsec:evalRoundingEtc}Internally, values are stored to $16$
significant figures (if available), calculations are carried out to
$16$ significant figures, but only rarely do we want to view the
result to $16$ figures. Generally, we round to some smaller number
of figures. The default rounding value is $6$, meaning by default
at most $6$ decimal places are shown. So far, all results have been
rounded to this figure, although not all digits are always displayed
– for instance if the sixth one is $0$, or the result is an integer. 

Like other elements of the display, both rounding value and the (dis)appearance
of trailing zeros can be customized, in this case by means of an optional
argument following the vv-list (or the formula if there is no vv-list).
This optional argument may contain up to four juxtaposed items from
six possibilities:
\begin{itemize}
\item a question mark ?, which gives boolean output, or 
\item an integer, the \emph{rounding value}, positive, negative or zero,
specifying how many decimal places to display the result to, or
\item an asterisk {*}, which pads the result with zeros should it not have
as many decimal places as the rounding value specifies, or
\item the character \texttt{x} (lower case!) which presents the result in
`proper' scientific notation (a form like $1.2345\times10^{5}$
for 123450), or
\item the character \texttt{t} (lower case!) which presents the result in
a bastardized scientific notation useful in tables (a form like $(5)1.2345$
for 123450), or
\item a character other than \texttt{?}, \texttt{{*}}, \texttt{x}, \texttt{t}
or an integer, usually  one of the letters\texttt{ e d} \texttt{E
D}, which presents the result in scientific notation with that character
as the exponent mark (a form like $1.2345\text{e}5$ for $123450$).
\end{itemize}
If you use \texttt{?} in the same specification as some other character,
the \texttt{?} prevails; if you use \texttt{x} in the same specification
as some other character except for \texttt{?}, the \texttt{x} prevails;
if you use \texttt{t} in the same specification as some other character
except for \texttt{?} or \texttt{x}, the \texttt{t} prevails.

If you repeat the character serving as the exponent mark in scientific
notation – say \verb`xx` or \verb`dd` – then scientific notation
extends to numbers in the interval \verb`[1,10)`.

If you repeat a question mark specifying boolean output, then the
formatting of that output is changed from \verb`1/0` to $T/F$ or
\verb`T/F` depending as there are two or three question marks used.

\subsection{Rounding value }

\label{subsec:evalRounding-value}If the number is displayed as a
decimal, the rounding value specifies the number of decimal places
displayed. If a number is displayed in scientific notation (see below
§\ref{subsec:evalScientificNotation}) that is still true, but it
can mean differences in the overall number of digits displayed. For
the moment, I show the effect of rounding in a purely decimal display:
\begin{centred}
\verb`$ \eval{ 1/3 }[4] $` $\Longrightarrow$ $ \eval{ 1/3 }[4] $
\end{centred}
In this case \verb`4` was entered in the number-format option and
the result is displayed to four decimal places. The default rounding
value is $6$:
\begin{centred}
\verb`$ \eval{ 35/3 } $` $\Longrightarrow$ $ \eval{ 35/3 } $
\end{centred}
Following the default behaviour in \verb`l3fp`, the calculational
engine which \verb`numerica` uses, `ties' are rounded to the nearest
\emph{even} digit. Thus a number ending $55$ with a `choice' of
rounding to $5$ or $6$ rounds up to the even digit $6$, and a number
ending $65$ with a `choice' of rounding to $6$ or $7$ rounds
down to the even digit $6$:
\begin{centred}
\verb`$ \eval{ 0.1234555 } $`  $\Longrightarrow\eval{0.1234555}$ 

\verb`$ \eval{ 0.1234565 } $`  $\Longrightarrow\eval{0.1234565}$ 
\end{centred}
\verb`l3fp` works to 16 significant figures and never displays more
than that number (and often fewer). 
\begin{itemize}
\item In the first of the following although I have specified a rounding
value of $19$ only $16$ decimal places are displayed, with the final
digit rounded up to $7$; 
\item in the second I have added $10$ zeros after the decimal point, meaning
that all $19$ decimal places specified by the rounding value can
be displayed since the 10 initial zeros do not contribute to the significant
figures; 
\item in the third I have changed the figure \emph{before} the decimal point
to $1$ so that the $10$ added zeros are now included among the significant
figures; 
\item and in the fourth, I have added $9$ digits before the decimal point:
\end{itemize}
\begin{centred}
\verb`$ \eval{ 0.1234567890123456789 }[19] $` $\Longrightarrow$
$\eval{ 0.1234567890123456789 }[19]$

\verb`$ \eval{ 0.00000000001234567890123456789 }[19] $` $\Longrightarrow$
$\eval{ 0.00000000001234567890123456789 }[19]$

\verb`$ \eval{ 1.00000000001234567890123456789 }[19] $` $\Longrightarrow$
$\eval{ 1.00000000001234567890123456789 }[19]$

\verb`$ \eval{ 987654321.1234567890123456789 }[19] $` $\Longrightarrow$
$\eval{ 987654321.1234567890123456789 }[19]$
\end{centred}
In all cases, no more than $16$ \emph{significant} figures are displayed,
although the number of decimal places displayed may exceed $16$ as
in the second example.

It is possible to use \emph{negative} rounding values. Such a value
zeroes the specified number of digits \emph{before} the decimal point.
\begin{centred}
\verb`$ \eval{ 987654321.123456789 }[-4] $` $\Longrightarrow$ $\eval{ 987654321.123456789 }[-4]$
\end{centred}
A rounding value of $0$ rounds to the nearest integer:
\begin{centred}
\verb`$ \eval{ 987654321.123456789 }[0] $` $\Longrightarrow$ $\eval{ 987654321.123456789 }[0]$
\end{centred}
If you wish to change the \emph{default} rounding value from $6$
to some other value, this can be done by creating or editing a file
\texttt{numerica.cfg} in a text editor; see §\ref{sec:settingsDefaults}.

\subsection{Padding with zeros }

\label{subsec:evalPadding-with-zeros}A result may contain fewer decimal
places than the rounding value specifies, the trailing zeros being
suppressed by default (this is how \texttt{l3fp} does it). Sometimes,
perhaps for reasons of presentation like aligning columns of figures,
it may be desirable to pad results with zeros. This is achieved by
inserting an asterisk, {*}, into the final optional argument of the
\verb`\eval` command:
\begin{centred}
\verb`$ \eval{ 1/4 }[4] $` $\Longrightarrow$ $ \eval{ 1/4 }[4] $,

\verb`$ \eval{ 1/4 }[4*] $` $\Longrightarrow$ $ \eval{ 1/4 }[4*] $.
\end{centred}

\subsection{Scientific notation }

\label{subsec:evalScientificNotation} \verb`l3fp` can output numbers
in scientific notation. For example, $1234$ is rendered as $\eval{1234}[e]$,
denoting $1.234\times10^{3}$ , and $0.008$ as $\eval{0.008}[e]$,
denoting $8\times10^{-3}$. The `e' here, the \emph{exponent mark},
separates the \emph{significand} ($1.234$) from the \emph{exponent}
($3$). In scientific notation, the significand always has one \emph{non-zero}
digit before the decimal point.\footnote{Except for $0$ itself.}

For scientific notation rounding still means the number of decimal
places displayed, but it can result in very different numbers of digits
being shown from the number shown in decimal form. To switch on output
in scientific notation in \verb`numerica` enter \verb`e` in the
trailing optional argument: 
\begin{centred}
\verb`$ \eval{ 123.456789 }[e] $` $\Longrightarrow$ $ \eval{ 123.456789 }[e] $. 
\end{centred}
The default rounding value $6$ is in play here, with seven digits
of the significand displayed overall, one preceding the decimal point,
six following it. Compare this with the same number rounded in decimal
form:
\begin{centred}
\verb`$ \eval{ 123.456789012345 } $` $\Longrightarrow$ $ \eval{ 123.456789012345 } $. 
\end{centred}
In this instance, nine digits are displayed, three before the decimal
point and six after. Similarly compare
\begin{centred}
\verb`$ \eval{ 0.0123456789 }[e] $` $\Longrightarrow$ $ \eval{ 0.0123456789 }[e] $ 
\end{centred}
with
\begin{centred}
\verb`$ \eval{ 0.0123456789 } $` $\Longrightarrow$ $ \eval{ 0.0123456789 } $.
\end{centred}
This time scientific notation has gained two extra decimal digits
to display.

Negative rounding values are pointless for scientific notation. A
zero might on occasion be relevant:
\begin{centred}
\verb`$ \eval{ 987654321 }[0e] $` $\Longrightarrow$ $ \eval{ 987654321 }[0e] $.
\end{centred}
Sometimes letters other than `e' are used to indicate scientific
notation, like `E' or `d' or `D'. With a few exceptions, \texttt{numerica}
allows any letter or text character to be used as the exponent marker:
\begin{centred}
\verb`\eval{$ 1/23456789 $}[4d]`\texttt{ $\Longrightarrow$} \eval{$ 1/23456789 $}[4d].
\end{centred}
But when \texttt{x} is inserted in the trailing optional argument,
the output is in the form $d_{0}.d_{1}\ldots d_{m}\times10^{n}$ (except
when $n=0$), where each $d_{i}$ denotes a digit.
\begin{centred}
\verb`\eval{$ 1/23456789 $}[4x]`\texttt{ $\Longrightarrow$ }\eval{$ 1/23456789 $}[4x] .
\end{centred}
The requirements of tables leads to another form of scientific notation.
Placing \texttt{t} in the trailing argument turns on this table-ready
form of notation:
\begin{centred}
\verb`\eval{$ 1/23456789 $}[4t]`\texttt{ $\Longrightarrow$ }\eval{$ 1/23456789 $}[4t].
\end{centred}
This is discussed more fully in the documentation for the \texttt{numerica-tables}
package.

In the next example three options are used in the trailing argument.
The order in which the items are entered does not matter: 
\begin{centred}
\verb`\eval{$ 1/125 $}[*e4]`  $\Longrightarrow$ \eval{$ 1/125 $}[*e4].
\end{centred}
Finally, to illustrate that `any' text character\footnote{Be sensible! An equals sign for instance might confuse \texttt{numerica}
into thinking the number-format option is the vv-list, and will certainly
confuse the reader.} save for \texttt{x} or \texttt{t} can be used to distinguish the
exponent, I use an @ character:
\begin{centred}
\verb`\eval{$ 1/125 $}[@4]` $\Longrightarrow$ \eval{$ 1/125 $}[@4].
\end{centred}

\subsubsection{Numbers in the interval \texttt{{[}1,10)}}

Usually when scientific notation is being used, numbers with magnitude
in the interval $[1,10)$ are rendered in their normal decimal form,
$3.14159$ and the like. Occasionally it may be desired to present
numbers in this range in scientific notation (this can be the case
in tables where the alignment of a column of figures might be affected).
\texttt{numerica} offers a means of extending scientific notation
to numbers in this range by repeating the letter chosen as the exponent
mark in the trailing optional argument.
\begin{centred}
\verb`\eval{$ \pi $}[4tt]` $\Longrightarrow$ \eval{$ \pi $}[4tt] 
\end{centred}

\subsubsection{\textbackslash eval{*} and scientific notation}

Scientific notation can be used for the numerical result output by
\verb`\eval*`:
\begin{centred}
\verb`\eval*{ \pi }[ee]` $\Longrightarrow$ \eval*{ \pi }[ee] 
\end{centred}
There is one catch: if you substitute \texttt{x} for \texttt{e} here,
\LaTeX{} will complain about a missing \verb`$`. An \texttt{x} in
the number-format option produces a \verb`\times` in the output which
requires a math environment. It is up to you, as the user, to provide
the necessary delimiters outside the \verb`\eval*` command. (This
applies even when \verb`\eval*` wraps around math delimiters.)

\subsection{Boolean output}

\label{subsec:evalBoolean-output}\verb`l3fp` can evaluate comparisons,
outputting $0$ if the comparison is false, $1$ if it is true. By
entering a question mark, \texttt{?}, in the trailing optional argument,
you can force \verb`numerica` to do the same depending as the result
of a calculation is zero or not. The expression being evaluated does
not need to be a comparison, \verb`$ \eval{\pi}[?] $` $\Longrightarrow$ $ \eval{\pi}[?]$,
but comparisons are what this is designed for.

Possible comparison relations are \verb`=`, \verb`<`, \verb`>`,
\verb`\ne`, \verb`\neq`, \verb`\ge`, \verb`\geq`, \verb`\le`,
\verb`\leq`. Although programming languages use combinations like
\verb`<=` or \verb`>=`, \texttt{numerica} does \emph{not} accept
these (they are not part of standard \emph{mathematical} usage) and
will generate an error. An example where the relation is equality
exhibits a numerological curiosity:\footnote{The \texttt{{[}p=.{]}} of this and the next example ensures a full
stop appears in the correct place; see §\ref{subsec:settingsPunctuation}.}
\begin{centred}
\verb`\eval[p=.]{\[ \frac1{0.0123456789}=81 \]}[5?]` $\Longrightarrow$
\eval[p=.]{\[ \frac1{0.0123456789}=81 \]}[5?]\smallskip{}
\end{centred}
Notice the $5$ alongside the question mark in the trailing argument.
That is critical. Change the $5$ to a $6$ (or omit it since the
default rounding value is $6$) and the outcome is different: 
\begin{centred}
\verb`\eval[p=.]{\[ \frac1{0.0123456789}=81 \]}[6?]` $\Longrightarrow$
\eval[p=.]{\[ \frac1{0.0123456789}=81 \]}[6?]
\end{centred}
Now the relation is false. Evaluating the fraction to more than $6$
places, say to $9$, we can see what is going on:
\begin{centred}
\verb`\eval{$ 1/0.0123456789 $}[9]` $\Longrightarrow$ \eval{$ 1/0.0123456789 $}[9]. 
\end{centred}

\subsubsection{Outputting \texttt{T} or \texttt{F}}

To my eye, outputting $0$ or $1$ in response to a `question' like
$1/0.0123456789=81$ is confusing. It is easy to change the boolean
output from $0,1$ to a more appropriate $F,T$, or \texttt{$\texttt{F,\texttt{T}}$}
by adding one or two more question marks respectively in the number-format
option.
\begin{centred}
\verb`\eval[p=.]{\[ \frac1{0.0123456789}=81 \]}[6???]` $\Longrightarrow$
\eval[p=.]{\[ \frac1{0.0123456789}=81 \]}[6???]
\end{centred}
The default boolean output format is chosen to be $0,1$ in case an
\verb`\eval` command is used within another \verb`\eval` command
(`nesting'– see Chapter~\ref{chap:Nesting}~). The inner command
needs to output a \emph{numerical} answer.

\subsubsection{Rounding error tolerance}

\label{subsec:evalToleranceRounding}If at least one of the terms
in a comparison is the result of a calculation, then it's value is
likely to contain rounding errors. What level of rounding error can
we tolerate before such errors interfere with the comparison being
made? \texttt{l3fp} tolerates none. It decides the truth or falsity
of a comparison to all $16$ significant figures: 1.000 0000 0000
0000 and 1.000 0000 0000 0001 are \emph{not }equal in \texttt{l3fp}.
But for most purposes this will be far too severe a criterion.

Suppose our comparison relation is $\varrho$, denoting one of =,
<, >, \verb`\le`, etc. If $X\rel Y$ then $X-Y\rel Y-Y$, i.e. $X-Y\rel0$.
This is what \texttt{numerica} does. It takes the right-hand side
of the relation from the left-hand side and then compares the \emph{rounded}
difference under $\varrho$ to $0$. The rounding value used is the
number specified with the question mark in the trailing argument of
the \verb`\eval` command or, if no number is present, the default
rounding value (`out of the box' this is $6$). Thus, in a recent
example, $1/0.0123456789-81$ when rounded to $5$ decimal places
is $0.00000$, indistinguishable from zero at this rounding value;
hence the equality $1/0.0123456789=81$ is true. But when rounded
to $6$ places it is $0.000001$ which \emph{is} distinguishable from
zero and so the equality is false. Truth or falsity depends on the
rounding value. 

When dealing with numbers generated purely mathematically, rounding
values of $5$ or $6$ are likely to be too small. More useful would
be rounding values closer to \texttt{l3fp}'s $16$ – perhaps $14$?
– depending on how severe the calculations are that generate the numbers.
However if the numbers we are dealing with come from outside mathematics,
from practical experiments perhaps, then even a rounding value of
$5$ or $6$ may be too large. 

\noindent\begin{minipage}[t]{1\columnwidth}%
\begin{shaded}%
Mathematically, the claim that $X=Y$ at a rounding value $n$ is
the claim that 
\[
\abs{X-Y}\le5\times10^{-(n+1)}.
\]
since this rounds \emph{down} to zero at $n$ places of decimals.
This gives a more accurate test of equality than doing things in the
opposite order – rounding each number first and then taking the difference.
One might, for instance, have numbers like $X=0.12345$, $Y=0.12335$.
Rounding to $n=4$ places, both round to $0.1234$ and yet the difference
between them is $0.0001$ – they are distinguishable numbers to $4$
places of decimals. This is why \texttt{numerica} forms the difference
\emph{before }doing the rounding.\end{shaded}%
\end{minipage}

\subsubsection{And, Or, Not}

For logical And \LaTeX{} provides the symbols \verb`\wedge` and \verb`\land`,
both displaying as $\land$, but \texttt{numerica} adds thin spaces
( \verb`\,` ) around the symbol for \verb`\land` (copying the package
\texttt{gn-logic14.sty}). For logical Or \LaTeX{} provides the symbols
\verb`\vee` and \verb`\lor`, both displaying as $\lor$, but again
\texttt{numerica} adds thin spaces around the symbol for \verb`\lor`. 
\begin{centred}
\verb`\eval{$ 1<2 \wedge 2<3 $}[??]` $\Longrightarrow$ \eval{$ 1<2 \wedge 2<3 $}[??],

\verb`\eval{$ 1<2 \land 2<3 $}[???]` $\Longrightarrow$ \eval{$ 1<2 \land 2<3 $}[???].
\end{centred}
To my eye the second of these with its increased space around the
wedge symbol displays the meaning of the overall expression better
than the first. Both And and Or have equal precedence; in cases of
ambiguity the user needs to parenthesize as necessary to clarify what
is intended.

\LaTeX{} provides two commands for logical Not, \verb`\neg` and \verb`\lnot`,
both displaying as $\lnot$ . Not binds tightly to its argument:
\begin{centred}
\verb`\eval{$ \lnot A \land B $}[A=0,B=0]` $\Longrightarrow$ \eval{$ \lnot A \land B $}[A=0,B=0].
\end{centred}
Here \verb`\lnot` acts only on the $A$; if it had acted on $A\land B$
as a whole the result would have been different:
\begin{centred}
\verb`\eval{$ \lnot(A \land B) $}[A=0,B=0]` $\Longrightarrow$ \eval{$ \lnot(A \land B) $}[A=0,B=0].
\end{centred}
For a little flourish, I evaluate a more complicated logical statement:\footnote{Quoting from an article in \emph{Quanta Magazine} (August 2020) by
Kevin Hartnett: `Let’s say you and two friends are planning a party.
The three of you are trying to put together the guest list, but you
have somewhat competing interests. Maybe you want to either invite
Avery or exclude Kemba. One of your co-planners wants to invite Kemba
or Brad or both of them. Your other co-planner, with an ax to grind,
wants to leave off Avery or Brad or both of them. Given these constraints,
you could ask: Is there a guest list that satisfies all three party
planners?' I have written $C$ for Kemba, $A$ and $B$ for Avery
and Brad.}
\begin{verbatim}
    \eval{$(A\lor\lnot C)\land(C\lor B)\land
           (\lnot A\lor\lnot B)$}[A=1,B=0,C=1][???]
\end{verbatim}
$\Longrightarrow$ \eval{$(A\lor\lnot C)\land(C\lor B)\land
                (\lnot A\lor\lnot B)$}[A=1,B=0,C=1][???].

\subsubsection{Chains of comparisons}

\texttt{numerica} can handle chains of comparisons like $1<2<1+2<5-1$.
`Behind the scenes' it inserts logical And-s into the chain, $1<2\land2<1+2\land1+2<5-1$,
and evaluates the modified expression:
\begin{centred}
\verb`\eval{$ 1<2<1+2<5-1 $}[?'']` $\Longrightarrow$ \eval{$ 1<2<1+2<5-1 $}[?''].
\end{centred}

\subsubsection{\texttt{amssymb} comparison symbols}

\label{subsec:evalAmssymb-comparisons}\texttt{numerica} accepts some
alternative symbols for the basic comparison relations from the \texttt{amssymb}
package provided that package is loaded, i.e. the preamble of your
document includes the statement
\begin{lyxcode}
\textbackslash usepackage\{amssymb\}
\end{lyxcode}
The variants from this package are: \verb`\leqq` ( $\leqq$ ), \verb`\leqslant`
( $\leqslant$ ), \verb`\geqq` (~$\geqq$~), and \verb`\geqslant`
( $\geqslant$ ).\footnote{No, that is not \texttt{eggplant}.} There
are also negations: \verb`\nless` ( $\nless$ ), \verb`\nleq` (~$\nleq$~),
\verb`\nleqq` ( $\nleqq$ ), \verb`\nleqslant` ( $\nleqslant$ ),
\verb`\ngtr` ( $\ngtr$ ), \verb`\ngeq` ( $\ngeq$ ), \verb`\ngeqq`
( $\ngeqq$ ), \verb`\ngeqslant` ( $\ngeqslant$ ).

\section{Calculational details}

\subsection{Arithmetic}

\label{subsec:evalArithmetic}Addition, subtraction, multiplication,
division, square roots, \emph{$n$}th roots, and exponentiating (raising
to a power) are all available. 

Multiplication can be rendered explicitly with an asterisk, 
\begin{centred}
\verb`\eval{$ 9*9 $}` $\Longrightarrow$ \eval{$ 9*9 $},
\end{centred}
but that's ugly. More elegant is to use \verb`\times`:
\begin{centred}
\verb`\eval{$ 9\times9 $}` $\Longrightarrow$ \eval{$ 9\times9 $}.
\end{centred}
\verb`\cdot` is also available and in many cases juxtaposition alone
suffices: 
\begin{centred}
\verb`\eval{$ \surd2\surd2 $}` $\Longrightarrow$ \eval{$ \surd2\surd2 $},

\verb`\eval{$ ab $}[a=123,b=1/123]` $\Longrightarrow$ \eval{$ ab $}[a=123,b=1/123].
\end{centred}
Division can be rendered in multiple ways too:
\begin{centred}
\verb`\eval{$ 42/6 $}` $\Longrightarrow$ \eval{$ 42/6 $},

\verb`\eval{$ 42\div6 $}` $\Longrightarrow$ \eval{$ 42\div6 $},
\end{centred}
or by using \verb`\frac` or \verb`\tfrac` or \verb`\dfrac` as in
\begin{centred}
\verb`\eval{$ \frac{42}6 $}` $\Longrightarrow$ \eval{$ \frac{42}6 $}.
\end{centred}
But note that since juxtaposition means multiplication, it is also
true that $42\tfrac{1}{6}$\texttt{ }evaluates to $7$ inside an \verb`\eval`
command rather than denoting `forty two and a sixth'. Hence if you
want to use `two and a half' and similar values in \texttt{numerica},
they need to be entered as improper fractions like $\tfrac{5}{2}$
or in decimal form, $2.5$ (as one does automatically in mathematical
expressions anyway because of the ambiguity in a form like $2\tfrac{1}{2}$).

Powers are indicated with the superscript symbol \verb`^`:
\begin{centred}
\verb` \eval{$ 3^{2^2} $}` $\Longrightarrow$ \eval{$ 3^{2^2} $} .
\end{centred}

\subsubsection{Square roots and $n$th roots}

\label{subsec:evalSquareRootsEtc}Let us check that 3, 4, 5 and 5,
12, 13 really are Pythagorean triples (I use \verb`\sqrt` in the
first, \verb`\surd` in the second):
\begin{centred}
\verb`\eval{$ \sqrt{3^2+4^2} $}` $\Longrightarrow$ \eval{$\sqrt{3^{2}+4^{2}}$},

\verb`\eval{$ \surd(5^2+12^2) $}` $\Longrightarrow$ \eval{$ \surd(5^2+12^2) $}.
\end{centred}
The \verb`\sqrt` command has an optional argument which can be used
for extracting $n$th roots of a number. This notation is generally
used when $n$ is a small positive integer like $3$ or $4$. This
practice is followed in \texttt{numerica}: $n$ must be a (not necessarily
small) \emph{positive integer}:
\begin{centred}
\verb`\eval{$ \sqrt[4]{81} $}` $\Longrightarrow$  \eval{$ \sqrt[4]{81} $},

\verb`\eval{$ \sqrt[n]{125} $}[n=\floor{\pi}]` $\Longrightarrow$
\eval{$ \sqrt[n]{125} $}[n=\floor{\pi}].
\end{centred}
If $n$ should not be a positive integer, an error message is generated;
see §\ref{sec:evalErrors}.

For display-style expressions, the \verb`\sqrt` command grows to
accommodate the extra vertical height; the surd doesn't. Here is an
example which anticipates a number of matters not discussed yet. It
shows \verb`\eval` wrapping around a square root containing various
formatting commands (negative spaces, \verb`\left` and \verb`\right`
nested within \verb`\bigg` commands), all digested without complaint
(see §\ref{subsec:evalFormatting-commands}; and see §\ref{subsec:settingsPunctuation}
for the \verb`[p=.]`): \medskip{}
\begin{verbatim}
    \eval[p=.]{\[ \sqrt[3]{\!
      \biggl(\!\left.\frac AD\right/\!\frac BC\biggr)
        }\]}[A=729,B=81,C=9,D=3]
\end{verbatim}
$\Longrightarrow$\eval[p=.]
  {\[ 
    \sqrt[3]
      {\!\biggl(\!\left.\frac AD\right/\!\frac BC\biggr)}
  \]}[A=729,B=81,C=9,D=3]

As implemented in \texttt{numerica}, $n$th roots found using \verb`\sqrt[n]`
are \verb`n=<integer>`\emph{ }roots. This raises an interesting question:
if the `$n$' of an $n$th root is the result of a calculation,
what happens with rounding errors? The calculation may not produce
an \emph{exact} integer. (This problem also arises with factorials;
see §\ref{subsec:evalFactorialsBinom}.) The solution employed in
\texttt{numerica} is to make what is considered an integer depend
on a rounding value. Most calculations will produce rounding errors
in distant decimal places. For `int-ifying' calculations, \texttt{numerica}
uses a rounding value of $14$: a calculation produces an integer
if, when rounded to $14$ figures, the result is an integer. Since
\texttt{l3fp} works to $16$ significant figures, a rounding value
of $14$ allows ample `elbowroom' for rounding errors to be accommodated
when judging what is an integer and what is not. As a practical matter
problems should not arise. 

\subsubsection{\emph{n}th roots of negative numbers}

Odd (in the sense of `not even') integral roots of \emph{negative}
numbers are available with \verb`\sqrt`,
\begin{centred}
\verb`\eval{$ \sqrt[3]{-125} $}` $\Longrightarrow$ \eval{$ \sqrt[3]{-125} $},

\verb`\eval{$ \sqrt[3]{-1.25} $}` $\Longrightarrow$ \eval{$ \sqrt[3]{-0.125} $}.
\end{centred}

\subsubsection{Inverse integer powers }

Of course to find an $n$th root we can also raise to the inverse
power,
\begin{centred}
\verb`\eval{$ 81^{1/4} $}` $\Longrightarrow$ \eval{$ 81^{1/4} $}.
\end{centred}
However, raising a \emph{negative} number to an inverse power generates
an error even when, mathematically, it should not. This matter, which
is a product of floating point representation of numbers, is discussed
below in §\ref{subsec:errorsInverse-powers}.

\subsection{Precedence, parentheses}

The usual precedence rules apply: multiplication and division bind
equally strongly and more strongly than addition and subtraction which
bind equally stongly. Exponentiating binds most strongly. Evaluation
occurs from the left.
\begin{centred}
\verb`\eval{$ 4+5\times6+3 $}` $\Longrightarrow$ \eval{$ 4+5\times6+3 $},

\verb`\eval{$ 6\times10^3/2\times10^2 $}` $\Longrightarrow$ \eval{$ 6\times10^3/2\times10^2 $},
\end{centred}
which may not be what was intended. Parentheses (or brackets or braces)
retrieve the situation:
\begin{centred}
\verb`\eval{$ (4+5)(6+3) $}` $\Longrightarrow$ \eval{$ (4+5)(6+3) $},

\verb`\eval{$ (6\times10^3)/(2\times10^2) $}` $\Longrightarrow$
\eval{$ (6\times10^3)/(2\times10^2) $}.
\end{centred}
Because exponentiating binds most strongly, negative values must be
parenthesized when raised to a power. If not,
\begin{centred}
\verb`\eval{$ -4^2 $}` $\Longrightarrow$ \eval{$ -4^2 $},
\end{centred}
which is clearly not $(-4)^{2}$. But
\begin{centred}
\verb`\eval{$ (-4)^2 $}` $\Longrightarrow$ \eval{$ (-4)^2 $}.
\end{centred}

\subsubsection{Command-form brackets}

\label{subsec:evalCommandBrackets}Note that brackets of all three
kinds are available also in command form: \verb`\lparen \rparen`
(from \verb`mathtools`) for \verb`( )`, \verb`\lbrack \rbrack`
for \verb`[ ]`, and \verb`\lbrace \rbrace` for \verb`\{ \}`.

\subsection{Modifiers\texttt{ (\textbackslash left \textbackslash right}, etc.)}

The  \verb`\left` and \texttt{\textbackslash right} modifiers and
also the series of \verb`\big...` modifiers\texttt{ }(\texttt{\textbackslash}\verb`bigl \bigr`,
\verb`\Bigl \Bigr`, \verb`\biggl \biggr`, \verb`\Biggl \Biggr`)
are available for use with all brackets (parentheses, square brackets,
braces):
\begin{verbatim}
    \eval[p=.]{\[ \exp\left(
        \dfrac{\ln2}{4}+\dfrac{\ln8}{4} 
      \right) \]}
\end{verbatim}
$\Longrightarrow$ \eval[p=.]{\[ \exp\left( \dfrac{\ln2}{4}+\dfrac{\ln8}{4} \right) \]}

\texttt{numerica} also accepts their use with \texttt{.} (dot) and
with \texttt{/} (as noted earlier, the \verb`[p]` and \verb`[p=.]`
are explained at §\ref{subsec:settingsPunctuation}):
\begin{centred}
\verb`\eval[p]{\[ \left.\dfrac{3+4}{2+1}\right/\!\dfrac{1+2}{4+5} \]}`
$\Longrightarrow$ \eval[p=.]{\[ \left. \dfrac{3+4}{2+1} \right/\!\dfrac{1+2}{4+5} \]}
\end{centred}
They can be nested.

\subsection{Trigonometric \& hyperbolic functions}

\LaTeX{} provides all six trignometric functions, \verb`\sin`, \verb`\cos`,
\verb`\tan`, \verb`\csc`, \verb`\sec`, \verb`\cot` and the three
principal inverses \verb`\arcsin`, \verb`\arccos`, \verb`\arctan`.
It also provides four of the six hyperbolic functions: \verb`\sinh`,
\verb`\cosh`, \verb`\tanh`, \verb`\coth`, and \emph{no} inverses.
\texttt{numerica} provides the missing hyperbolic functions, \verb`\csch`
and \verb`\sech`, and all missing inverses, the three trigonometric
and all six hyperbolic: \verb`\arccsc`, \verb`\arcsec`, \verb`\arccot`,
and \verb`\asinh`, \verb`\acosh`, \verb`\atanh`, \verb`\acsch`,
\verb`\asech`, \verb`\acoth`. (\emph{HMF} writes $\text{arcsinh}$,
$\text{arccosh}$, etc. and ISO recommends $\text{arsinh}$, $\text{arcosh}$,
etc. The first seems ill-advised, the second not widely adopted. At
present neither is catered for in \texttt{numerica}.)\emph{ }
\begin{centred}
\verb`\eval{$ \arctan1/1\deg $}` $\Longrightarrow$ \eval{$ \arctan 1/1\deg $} ,

\verb`\eval{$ \atanh\tanh3 $}` $\Longrightarrow$ \eval{$ \atanh\tanh3 $} .
\end{centred}
Inverses can also be constructed using the `$-1$' superscript notation.
Thus
\begin{centred}
\verb`\eval{$ \sin^{-1}(1/\surd2)/1\deg $}` $\Longrightarrow$ \eval{$ \sin^{-1}(1/\surd2)/1\deg $} ,

\verb`\eval{$ \tanh\tanh^{-1}0.5 $}` $\Longrightarrow$ \eval{$ \tanh\tanh^{-1}0.5 $} .
\end{centred}
\noindent\begin{minipage}[t]{1\columnwidth}%
\begin{shaded}%

\subsubsection*{Hyperbolic functions}

Please note that \texttt{l3fp} does not (as yet) provide \emph{any}
hyperbolic functions natively. The values \texttt{numerica} provides
for these functions are \emph{calculated} values using familiar formulas
involving exponentials (for the direct functions) and natural logarithms
and square roots for the inverses. Rounding errors mean the values
calculated may not have $16$-figure accuracy. The worst `offenders'
are likely to be the least used, \verb`\acsch` and \verb`\asech`.
For instance,
\[
\acsch x=\ln\left[\frac{1}{x}+\left(\frac{1}{x^{2}}+1\right)^{1/2}\right],
\]

\begin{centred}
\verb`\eval{$ \csch \acsch 7 $}[16]` $\Longrightarrow$ \eval{$ \csch \acsch 7 $}[16].
\end{centred}
\end{shaded}%
\end{minipage}

\subsection{Logarithms}

The natural logarithm \verb`\ln`, base $10$ logarithm \verb`\lg`,
and binary or base $2$ logarithm \verb`\lb` are all recognized,
as is \verb`\log`, preferably with a subscripted base:
\begin{centred}
\verb`\eval{$ \log_{12}1728 $}` $\Longrightarrow$ \eval{$ \log_{12}1728 $}
\end{centred}
If there is no base indicated, base $10$ is assumed. (The notations
\verb`\ln`, \verb`\lg`, and \verb`\lb` follow ISO 80000-2 recommendation,
which frowns upon the use of the unsubscripted \verb`\log` although
only \verb`\ln` appears widely used.) The base need not be explicitly
entered as a number. It could be entered as an expression or be specified
in the vv-list:
\begin{centred}
\verb`\eval*{$ \log_b c $}[b=2,c=1024]` $\Longrightarrow$ \eval*{$ \log_b c $}[b=2,c=1024],
\end{centred}
the log to base $2$ in this case. It is possible to use the unadorned
\verb`\log` with a base different from $10$; if you wish to do this
only for a particular calculation see §\ref{subsec:settingsLogBase},
or see §\ref{sec:settingsDefaults} if you want to make this default
behaviour.

\subsection{Other unary functions}

Other unary functions supported are the exponential function \verb`\exp`
and signature function \verb`\sgn` (equal to $-1$, $0$, or $1$
depending as its argument is $<0$, $=0$, or $>0$). 

\subsection{Squaring, cubing, \ldots unary functions}

\texttt{numerica} has no difficulty reading a familiar but `incorrectly
formed' expression like 
\[
\sin^{2}1.234+\cos^{2}1.234.
\]
You do not have to render it $(\sin1.234)^{2}+(\cos1.234)^{2}$ or
(heaven forbid) $(\sin(1.234))^{2}+(\cos(1.234))^{2}$. The everyday
usage is fine:
\begin{centred}
\verb`\eval{$ \sin^2\theta+\cos^2\theta $}[\theta=1.234]` $\Longrightarrow$
\eval{$ \sin^2\theta+\cos^2\theta $}[\theta=1.234] .
\end{centred}
Equally \texttt{numerica} has no difficulty reading the `correct'
but pedantic form
\begin{centred}
\verb`\eval{$ (\sin(\theta))^2+(\cos(\theta))^2 $}[\theta=1.234]`
$\Longrightarrow$ \eval{$ (\sin(\theta))^2+(\cos(\theta))^2 $}[\theta=1.234] .
\end{centred}
A hyperbolic identity is corroborated in this example:
\begin{centred}
\verb`\eval{$ \sinh 3x $}[x=1]` $\Longrightarrow$ \eval{$ \sinh 3x $}[x=1],\medskip{}

\verb`\eval{$ 3\sinh x+4\sinh^3x $}[x=1]` $\Longrightarrow$ \eval{$ 3\sinh x+4\sinh^3x $}[x=1].
\end{centred}
In fact all named unary functions in \texttt{numerica} can be squared,
cubed, etc., in this `incorrect' but familiar way, although the
practice outside the trigonometric and hyperbolic context seems (vanishingly?)
rare.

When the argument of the function is parenthesized and raised to a
power – like $\sin(\pi)^{2}$ – it is read by \texttt{numerica} as
the `sine of the square of pi', $\sin(\pi^{2})$, and \emph{not
}as the `square of the sine of pi', $(\sin\pi)^{2}$:
\begin{centred}
\verb`\eval{$ \sin(\pi)^2 $}` $\Longrightarrow$ \eval{$ \sin(\pi)^2  $} .
\end{centred}
Things are done like this in \texttt{numerica} above all to handle
the logarithm in a natural way. Surely $\ln x^{n}=n\ln x$, i.e. $\ln x^{n}=\ln(x^{n})$
rather than $(\ln x)^{n}$? And if we wish to write (as we do) $\ln(1+1/n)^{n}=n\ln(1+1/n)=1-1/2n+1/3n^{2}-\ldots$
to study the limiting behaviour of $(1+1/n)^{n}$, then we cannot
avoid $\ln(x)^{n}=n\ln(x)=\ln(x^{n})$ too.

\subsection{\emph{n}-ary functions}

The functions of more than one variable ($n$-ary functions) that
\texttt{numerica} supports are \verb`\max`, \verb`\min` and \verb`\gcd`,
greatest common divisor. The comma list of arguments to \verb`\max`,
\verb`\min` or \verb`\gcd` can be of arbitrary length. The arguments
themselves can be expressions or numbers. For \verb`\gcd`, \emph{non-integer
arguments are truncated to integers}. Hence both $y$ and $3y$ are
independently truncated in the following example – to $81$ and $243$
respectively:
\begin{centred}
\verb`\eval{$ \gcd(12,10x^2,3y,y,63) $}[y=1/0.0123456789,x=3]` $\Longrightarrow$
\eval{$ \gcd(12,10x^2,3y,y,63) $}[y=1/0.0123456789,x=3] .
\end{centred}
(The truncation occurs in the argument of \verb`\gcd`, not in the
vv-list.) 

For $n$-ary functions, squaring, cubing, etc. follows a different
pattern from that for unary functions. For \verb`\max`, \verb`\min`,
\verb`\gcd` the argument of the function is a comma list. Squaring
the argument makes no sense. We understand the superscript as applying
to the function as a whole. (Consistency is not the point here; it
is what mathematicians do that \texttt{numerica} tries to accommodate.)
\begin{centred}
\verb`\eval{$ \gcd(3x,x,\arcsin 1/\deg)^2 $}[x=24]` $\Longrightarrow$
\eval{$ \gcd(3x,x,\arcsin 1/\deg)^2 $}[x=24] .
\end{centred}

\subsection{Delimiting arguments with brackets \& modifiers }

Arguments of unary and $n$-ary functions can be delimited not only
with parentheses, but also with square brackets and braces, both in
explicit character form and also in the command form of §\ref{subsec:evalCommandBrackets}.
The brackets, of whatever kind, can be qualified with \verb`\left \right`,
\verb`\bigl \bigr`, etc.\footnote{See §\ref{subsec:settingsPunctuation} for the \texttt{{[}p=.{]}}
(which ensures the concluding full stop appears in the correct place.} 
\begin{centred}
\verb`\eval[p=.]{\[ \sin\left\lbrack \dfrac\pi{1+2+3}\right\rbrack \]}`
$\Longrightarrow$\eval[p=.]{\[ \sin\left\lbrack\dfrac\pi{1+2+3}\right\rbrack \]}
\end{centred}

\subsection{Absolute value, floor \& ceiling functions}

It is tempting to use the \texttt{|} key on the keyboard for inserting
an absolute value sign. \texttt{numerica} accepts this usage, but
it is deprecated. The spacing is incorrect – compare $|-l|$ using
\texttt{|} against $\lvert-l\rvert$ using \verb`\lvert \rvert`.
Also, the identity of the left and right delimiters makes nested absolute
values difficult to parse. \texttt{numerica} does not attempt to do
so. Placing an absolute value constructed with \texttt{|} within another
absolute value constructed in the same way is likely to produce a
compilation error or a spurious result. \verb`\lvert \rvert` are
better in every way except ease of writing.\texttt{ }To aid such ease
\texttt{numerica }provides the \verb`\abs` function (using the \texttt{\textbackslash DeclarePairedDelimiter}
command of the \texttt{mathtools} package). This takes a mutually
exclusive star (asterisk) or square bracketed optional argument, and
a mandatory braced argument. The starred form expands to \verb`\left\lvert #1 \right\rvert`
where \verb`#1` is the mandatory argument:
\begin{centred}
\verb`\eval[p=.]{\[ 3\abs*{\frac{\abs{n}}{21}-1} \]}[n=-7]` $\Longrightarrow$
\eval[p=.]{\[ 3\abs*{\frac{\abs{n}}{21}-1} \]}[n=-7]
\end{centred}
The optional argument provides access to the \verb`\big...` modifiers:
\begin{verbatim}
    \eval[p=.]{\[
      \abs[\Big]{\abs{a-c}-\abs[\big]{A-C}} 
    \]}[A=12,a=-10,C=7,c=-5] 
\end{verbatim}
$\Longrightarrow$     \eval[p=.]{\[ 
      \abs[\Big]{\abs{a-c}-\abs[\big]{A-C}} 
    \]}[A=12,a=-10,C=7,c=-5] 

The form without either star or square bracket option dispenses with
the modifiers altogether:
\begin{centred}
\verb`\eval{$ \tfrac12(x+y)+\tfrac12\abs{x-y} $}[x=-3,y=7].` $\Longrightarrow$
\eval{$ \tfrac12(x+y)+\tfrac12\abs{x-y} $}[x=-3,y=7].
\end{centred}
As noted, the star and square bracketed option are mutually exclusive
arguments.

\texttt{numerica} also provides the functions \verb`\floor` and \verb`\ceil`,
defined in the same way, taking a mutually exclusive star or square
bracketed optional argument and for the starred forms expanding to
\verb`\left\lfloor #1 \right\rfloor` and \verb`\left\lceil #1 \right\rceil`
where \verb`#1` is the mandatory argument, and for the square bracket
option forms replacing the \verb`\left` and \verb`\right` with the
corresponding \verb`\big` commands. The form without star or square-bracket
option dispenses with any modifier at all.
\begin{centred}
\verb`\eval{$ \floor{-\pi} $}` $\Longrightarrow$ \eval{$ \floor{-\pi} $},

\verb`\eval{$ \ceil{\pi} $}` $\Longrightarrow$ \eval{$ \ceil{\pi} $}.
\end{centred}
The floor function, $\lfloor x\rfloor$, is the greatest integer $\le x$;
the ceiling function, $\lceil x\rceil$ is the smallest integer $\ge x$.
Like the absolute value, the floor and ceiling functions, can be nested: 
\begin{centred}
\verb`\eval{$ \floor{-\pi+\ceil{e}} $}` $\Longrightarrow$ \eval{$ \floor{-\pi+\ceil{e}} $}.
\end{centred}

\subsubsection{Squaring, cubing, \ldots{} absolute values, etc.}

These three functions can be raised to a power \emph{without} extra
parentheses:
\begin{centred}
\verb`\eval{$ \ceil{e}^2 $},` $\Longrightarrow$ \eval{$ \ceil{e}^2 $},

\verb`\eval{$ \abs{-4}^2 $}.` $\Longrightarrow$ \eval{$ \abs{-4}^2 $}.
\end{centred}

\subsection{Factorials, binomial coefficients}

\label{subsec:evalFactorialsBinom}Factorials use the familiar trailing
\texttt{!} notation: 
\begin{centred}
\verb`\eval{$ 7! $}` $\Longrightarrow$ \eval{$ 7! $},

\verb`\eval{$ (\alpha+\beta)!-\alpha!-\beta! $}[\alpha=2,\beta=3]`
$\Longrightarrow$ \eval{$ (\alpha+\beta)!-\alpha!-\beta! $}[\alpha=2,\beta=3].
\end{centred}
The examples illustrate how \texttt{numerica} interprets the argument
of the factorial symbol:\texttt{ }it `digests' 
\begin{itemize}
\item a preceding (possibly multi-digit) integer, or
\item a preceding variable token, or
\item a bracketed expression, or
\item a bracket-like expression.
\end{itemize}
A bracket-like expression is an absolute value, floor or ceiling function,
since they delimit arguments in a bracket-like way:
\begin{centred}
\verb`\eval{$ \abs{-4}!+\floor{\pi}!+\ceil{e}! $}` $\Longrightarrow$
\eval{$ \abs{-4}!+\floor{\pi}!+\ceil{e}! $}.
\end{centred}
The result of feeding the factorial an expression different in kind
from one of these four cases may give an error message or an unexpected
result. Use parentheses around such an expression; for example write
$(3^{2})!$, rather than $3^{2}!$.

Nesting of brackets for factorials is accepted:
\begin{centred}
\verb`\eval{$ ((5-2)!+1)! $}` $\Longrightarrow$ \eval{$ ((5-2)!+1)! $}.
\end{centred}
The factorials of negative integers or of non-integers are not defined
in \texttt{numerica}, and again there is the problem met in relation
to $n$th roots of what happens if the argument of a factorial is
the result of a calculation and rounding errors mean it is not an
exact integer. This problem is unlikely to be of practical concern
since \texttt{numerica} rounds the result of such a calculation by
default to $14$ significant figures before offering it to the factorial.
Since \texttt{l3fp} works to $16$ significant figures, there is ample
`elbowroom' to accommodate rounding errors before the result of
a calculation ceases to round to an integer. 

\subsubsection{Double factorials}

The double factorial, written $n!!,$ is the product $n(n-2)(n-4)\ldots\times4\times2$
when $n$ is even and the product $n(n-2)(n-4)\ldots\times3\times1$
when $n$ is odd. 
\begin{centred}
\verb`\eval{$ 6!! $}` $\Longrightarrow$ \eval{$ 6!! $},

\verb`\eval{$ n!! $}[n=\sqrt{49}]` $\Longrightarrow$ \eval{$ n!! $}[n=\sqrt{49}],
\end{centred}
Since $n!=n!!(n-1)!!$ it follows that
\[
n!!=\frac{n!}{(n-1)!!}=\frac{(n+1)!}{(n+1)!!}.
\]
Putting $n=0$ in the outer equality shows that $0!!=1$. Now putting
$n=0$ in the left equality gives $(-1)!!=1$. Double factorials therefore
are defined for integers $\ge-1$.

\subsubsection{Binomial coefficients}

Binomial coefficients are entered in \LaTeX{} with the \verb`\binom`\textbf{
}command. It takes two arguments and has a text-style version \verb`\tbinom`
and a display-style version \verb`\dbinom`. As implemented in {\ttfamily\verb`numerica`},
these are \emph{generalised} binomial coefficients:
\[
\binom{x}{k}=\frac{x(x-1)\dots(x-k+1)}{k(k-1)\dots1},\quad(x\in\mathbb{R},~k\in\mathbb{N}),
\]
where $x$ need not be a non-negative integer, and where $\binom{x}{0}=1$
by definition. Although the first (or upper) argument can be any real
number, the lower argument \emph{must} be a non-negative integer.
Thus, \verb`\eval{$ \tbinom53 $}` $\Longrightarrow$ \eval{$\tbinom53$},
\verb`\eval{$ \tbinom70 $}` $\Longrightarrow$ \eval{$\tbinom70$},
\verb`\eval{$ \tbinom{4.2}3 $}` $\Longrightarrow$ \eval{$\tbinom{4.2}3$},
but if the second (or lower) argument of \verb`\binom` is \emph{not}
a non-negative integer, {\ttfamily\verb`numerica`} displays
a message; see §\ref{subsec:errorsIntegerArgs}. 

\subsection{Sums and products}

\verb`numerica` recognizes sums (\verb`\sum` displaying as $\sum$)
and products (\verb`\prod` displaying as $\prod$), and expects both
symbols to have lower and upper summation/product limits specified.
The lower limit must be given in the form \emph{sum/prod variable
= initial value}; the upper limit requires only the final value to
be specified (although it can also be given in the form \emph{sum/prod
variable = final value}). The values may be expressions depending
on other variables and values but must evaluate to integers (or infinity
– see §\ref{sec:settingsInfiniteSumsProds}). Evaluating to an integer
means that they \emph{round} to an integer, using a rounding value
that is set by default to $14$; (recall that \texttt{l3fp} works
to $16$ significant figures). If a limit evaluates to a non-integer
at this `int-ifying' rounding value, an error message results. (To
change this `int-ifying' rounding value, see §\ref{subsec:defaultsIntifyingRounding}.)

As an example of expressions in the limits, this example uses the
floor and ceiling functions to convert combinations of constants to
integers (the \verb`[p]` is explained in §\ref{subsec:settingsPunctuation}),
\begin{centred}
\verb`\eval[p]{\[ \sum_{n=\floor{\pi/e}}^{\ceil{\pi e}}n \]}` $\Longrightarrow$\eval[p]{\[ \sum_{n=\floor{\pi/e}}^{\ceil{\pi e}}n \]}
\end{centred}
\noindent (which is $\sum_{n=1}^{9}n$).\emph{ }If the upper limit
is less than the lower limit the result is zero. Notice that there
is no vv-list. The summation variable does not need to be included
there unless there are other variables that depend on it. However,
in the case
\begin{centred}
\verb`\eval[p]{\[ \sum_{k=1}^N\frac1{k^3} \]}[N=100][4]` $\Longrightarrow$
\eval[p]{\[ \sum_{k=1}^N\frac1{k^3} \]}[N=100][4]
\end{centred}
the upper limit $N$ is necessarily assigned a value in the vv-list. 

To the author it seems natural to enter the lower limit first, immediately
after the \verb`\sum` command (the sum is \emph{from }something \emph{to}
something), but no problem will accrue if the upper limit is placed
first (after all, the appearance of the formula in the pdf is the
same):
\begin{centred}
\verb`\eval[p=.]{\[ \sum^N_{k=1}\frac1{k^3} \]}[N=100][4]` $\Longrightarrow$
\eval[p=.]{\[ \sum^N_{k=1}\frac1{k^3} \]}[N=100][4]
\end{centred}
Another example of a sum, using binomial coefficients this time, is 
\begin{centred}
\verb`\eval[p]{\[ \sum_{m=0}^5\binom{5}{m}x^m y^{5-m} \]}[x=0.75,y=2.25]`
$\Longrightarrow$ \eval[p]{\[ \sum_{m=0}^5\binom{5}{m}x^m y^{5-m} \]}[x=0.75,y=2.25]
\end{centred}
which is just \verb`\eval{$(x+y)^5$}[x=0.75,y=2.25]` $\Longrightarrow$
\eval{$ (x+y)^5 $}[x=0.75,y=2.25], or $3^{5}$.

Now let's calculate a product:
\begin{verbatim}
    \eval[p]{\[
      \prod_{k=1}^{100} 
        \biggl(\frac{x^2}{k^2\pi^2} +1\biggr)
          \]}[x=1][3]
\end{verbatim}
$\Longrightarrow$ \eval[p]{\[\prod_{k=1}^{100} \biggl(\frac{x^2}{k^2\pi^2} +1\biggr)\]}[x=1][3]

\noindent to be compared with \verb`\eval{$ \sinh 1 $}[3]` $\Longrightarrow$
\eval{$ \sinh1 $}[3]. Obviously more terms than $100$ are required
in the product to achieve 3-figure accuracy. 

\subsubsection{Infinite sums and products}

How many more? Let's `go the whole hog' and put $\infty$ in the
upper limit of this product:
\begin{verbatim}
    \eval[p=.]{\[
      \prod_{k=1}^{\infty} 
        \biggl(\frac{x^2}{k^2\pi^2} +1\biggr)
          \]}[x=1][3] 
\end{verbatim}
$\Longrightarrow$ \eval[p=.]
  {\[
    \prod_{k=1}^{\infty} 
      \biggl(\frac{x^2}{k^2\pi^2} +1\biggr)
  \]}[x=1][3] 

\noindent Disappointingly, we still get the same result, deficient
by $1$ in the third decimal place. Obviously \texttt{numerica} has
not multiplied an infinite number of terms and, just as obviously,
the finite number of terms it \emph{has} multiplied are too few. How
\texttt{numerica} decides when to stop evaluating additional terms
in an infinite sum or product is discussed later, §\ref{sec:settingsInfiniteSumsProds}.

For this particular product the problem is that it converges slowly.
Any criterion for when to stop multiplying terms or, for an infinite
sum adding terms, seems bound to fail for some product or series.
Presumably any stopping criterion must measure smallness in some way.
But terms of, for example, the divergent harmonic series $\sum(1/n)$
can always be found smaller than any value we care to specify. It
is not surprising that a sufficiently slowly converging product or
series falls foul of a given criterion.

The default criterion however can be changed. Because this involves
values assigned in the settings option of the \verb`\eval` command,
I discuss infinite sums and products in the section discussing that
optional argument; see §\ref{sec:settingsInfiniteSumsProds}.

Other infinite sums converge more rapidly, and the default settings
work admirably. For example \verb`\eval{$ (1+0.1234)^{4.321} $}`
$\Longrightarrow$ \eval{$ (1+0.1234)^{4.321} $}. Using binomial
coefficients we can express this as an infinite sum:\medskip{}
\begin{verbatim}
    \eval[p=.]{\[
        \sum_{n=0}^{\infty}\binom{\alpha}{n}x^{n}
           \]}[\alpha=4.321,x=0.1234]
\end{verbatim}
$\Longrightarrow$ \eval[p=.]
  {\[
    \sum_{n=0}^{\infty}\binom{\alpha}{n}x^{n}
  \]}[\alpha=4.321,x=0.1234]

\subsection{Formatting commands}

\label{subsec:evalFormatting-commands}There are many formatting commands
which change the layout of a formula on the page but do not alter
its calculational content. \verb`numerica` copes with a great many
of these formatting commands, although there will surely be some that
it has overlooked and which will trigger an `Unknown token' message;
see §\ref{sec:evalErrors}. \footnote{Please contact the author in that case: ajparsloe@gmail.com}

\subsubsection{Spaces, phantoms, struts}

These include cryptic forms like \verb`\,` \verb`\:` and \verb`\>`,
\verb`\;` and the corresponding `verbose' forms, \verb`\thinspace`,
\verb`\medspace` and \verb`\thickspace` and their negative equivalents
\verb`\!` or \verb`\negthinspace`, \verb`\negmedspace` and \verb`\negthickspace`:
\begin{centred}
\verb`\eval{$ 1\negthickspace+\negthickspace 1 $}` $\Longrightarrow$
\eval{$ 1\negthickspace+\negthickspace 1 $}
\end{centred}
which gives the text spacing of 1+1 as against the usual math spacing
$1+1$ but doesn't affect the result of the calculation. 

Other spacing commands are \verb`\quad` and \verb`\qquad`, and \verb`\hspace{arg}`
and \verb`\mspace{arg}`. For \verb`\hspace` there is also a starred
form, \verb`\hspace*{arg}`. Phantoms similarly take an argument:
\verb`\phantom{arg}`, \verb`\hphantom{arg}` and \verb`\vphantom{arg}`.
\begin{centred}
\verb`\eval{$ 1\hphantom{mmm}+\hphantom{mmm}1 $}` $\Longrightarrow$
\eval{$ 1\hphantom{mmm}+ \hphantom{mmm}1 $}.
\end{centred}
Like \verb`\vphantom`, struts allow vertical spacing adjustments.
\verb`numerica` should digest both \verb`\xmathstrut[optarg]{arg}`
from \verb`mathtools` and its `baby cousin' \verb`\mathstrut`
from \TeX . An example from \emph{The \TeX{} book} demonstrating the
use of \verb`\mathstrut` is
\begin{verbatim}
     \eval{$\sqrt{\mathstrut a}+\sqrt{\mathstrut d}+
           \sqrt{\mathstrut y}$}[a=4,d=9,y=16] 
\end{verbatim}
$\Longrightarrow$ \eval{$\sqrt{\mathstrut a}+\sqrt{\mathstrut d}+\sqrt{\mathstrut y}$}[a=4,d=9,y=16], 

And here is an evaluation of an expression from the \verb`mathtools`
documentation using \verb`\xmathstrut`:
\begin{verbatim}
    \eval{\[ \frac{ \frac{ \xmathstrut{0.1} x-1 }
           { \xmathstrut{0.25} x-\sin{ x} } }
             {\xmathstrut{0.4} \sqrt{ 10-x } } \]}
               [x=\pi/6]
\end{verbatim}
$\Longrightarrow$ \eval{\[ \frac{ \frac{ \xmathstrut{0.1} x-1 }
           { \xmathstrut{0.25} x-\sin{ x} } }
             {\xmathstrut{0.4} \sqrt{ 1-x } } \]}
               [x=\pi/6]

\subsubsection{\texttt{\textbackslash splitfrac}}

The \verb`mathtools` package provides \verb`\splitfrac` and \verb`\splitdfrac`
to aid handling of clumsy fractions. The documentation gives an (artificial)
example of use. I've mangled it to produce an even more ridiculous
illustration, adding to the mess an enormous square root, the modifiers
\verb`\left` and \verb`\right`, and the command-form alternatives
to parentheses, \verb`\lparen` and \verb`\rparen`; also the use
of \verb`\dfrac`. A little mental arithmetic will convince that we
are evaluating the square root of $(9\times7)^{2}$ which indeed is
what we get:\footnote{For the \texttt{{[}p=.,vvd={]}} see §\ref{subsec:settingsPunctuation}
and §\ref{subsec:settingsVvDisplayChangeLocal}. The first puts the
concluding full stop in the right place; the second suppresses the
vv-list.} \medskip{}
\begin{verbatim}
	\eval[p=.,vvd=]{\[
    \sqrt{\left\lparen
      \frac{ \splitfrac{xy + xy + xy + xy + xy}
              {+ xy + xy + xy + xy}
           }
           { \dfrac z7}
          \right\rparen \left\lparen
            \frac{ \splitdfrac{xy + xy + xy + xy + xy}
                     {+ xy + xy + xy + xy}
                 }
                 {\dfrac z7}\right\rparen}
      \]}[x=2,y=5,z=10]
\end{verbatim}
$\Longrightarrow$ \eval[p=.,vvd=]{\[
    \sqrt{\left\lparen
      \frac{ \splitfrac{xy + xy + xy + xy + xy}
              {+ xy + xy + xy + xy}
           }
           { \dfrac z7}
          \right\rparen \left\lparen
            \frac{ \splitdfrac{xy + xy + xy + xy + xy}
                     {+ xy + xy + xy + xy}
                 }
                 {\dfrac z7}\right\rparen}
      \]}[x=2,y=5,z=10]

\subsubsection{Colour}

\label{subsec:Colour}(Anglicised spelling at least for the heading!)
If you add to the preamble of your document the line
\begin{lyxcode}
\textbackslash usepackage\{color\}
\end{lyxcode}
two commands become available, \verb`\textcolor[optarg]{arg1}{arg2}`
and the declaration form of command, \verb`\color[optarg]{arg}`.
\verb`numerica` readily accepts the former in a formula to be evaluated:
\begin{centred}
\verb`\eval{$ \sin \tfrac\pi6n\textcolor{red}{T}+1 $}[T=9,n=3]` $\Longrightarrow$
\eval{$ \sin \tfrac\pi6n\textcolor{red}{T}+1  $}[T=9,n=3]
\end{centred}
(assuming you had some wish to highlight the time $T$). 

However there are restrictions on the use of \verb`\color` in \verb`\eval`
commands. \verb`\color` is a \emph{declaration} form of command.
It has effect until the end of the current group or environment. If
you want to restrict it to only part of that group you need to em-brace
the command and what it is to apply to,
\begin{lyxcode}
<pre-stuff>\{\textbackslash color\{red\}<red-stuff>\}<post-stuff\}
\end{lyxcode}
but that is where the problem arises. \verb`numerica` does not check
for `unannounced' brace groups. It expects a brace group to be introduced
by a preceding instruction like \verb`\sqrt` or \verb`\frac` or
\verb`^`. When announced in this way, \verb`numerica` can handle
the brace group appropriately. But the brace group \verb`{\color{red}<red-stuff>}`
is not so announced. \verb`numerica`'s parsing routine will not recognize
what it has just swallowed and a \LaTeX{} error will result. So, \verb`\color`
cannot be used in a formula in a `naked' or unannounced brace group.
Writing \verb`\eval{$ \color{red} \sin \tfrac\pi6nT+1 $}[T=9,n=3]`
is fine, as is 
\begin{centred}
\verb`\eval{$  \sin \tfrac\pi6nT+1 \color{red} $}[T=9,n=3]` $\Longrightarrow$
\eval{$ \sin \tfrac\pi6nT+1 \color{red} $}[T=9,n=3].
\end{centred}
So too, because the \verb`\frac` introduces the confining brace group,
is
\begin{centred}
\verb`\eval{$ \frac{\color[gray]{0.5}A}B $}[A=12,b=4]` $\Longrightarrow$\eval{$ \frac{\color[gray]{0.5}A}B $}[A=12,B=4],
\end{centred}
where both arguments of the \verb`\color` command are used for grayscale
output. 

But trying something like \verb`\eval{$ 3{\color[gray]{0.5}x}+1 $}[x=2]`
will cause a \LaTeX{} error and halt compilation since there is no
command announcing the brace group confining the \verb`\color` command.

\subsubsection{\texttt{\textbackslash text}, \texttt{\textbackslash mbox}, font
commands}

\label{subsec:Text-mbox-fonts}Following a rethink of the behaviour
of a number of font and formatting commands, in version 2 of \verb`numerica`
the content of a \verb`\text` or \verb`\mbox` command is \emph{invisible}
to the \verb`\eval` command. \emph{This behaviour is different from
that of version 1.} Now the content is ignored in a calculation,
\begin{centred}
\verb`\eval*{ 1/0.0123456789 \mbox{approx.} }[5]` $\Longrightarrow$
\eval*{ 1/0.0123456789 \mbox{approx.}}[5],
\end{centred}
even when the \verb`\text` or \verb`\mbox` contains mathematical
content.

Conversely, the content of font commands (like \verb`\mathbf` or
\verb`\mathcal`) is \emph{visible} to \verb`\eval`. This becomes
useful should numbers be input in scientific notation (see §\ref{subsec:settingsInputtingSciNotation}).
As well as the math font commands, \verb`\eval` also accepts \verb`\textrm`,
\verb`\textsf` and \verb`\texttt`. Thus a number in scientific notation
like \verb`2e-1` appearing in the formula or the vv-list can display
correctly by wrapping it in a \verb`\textrm` or \verb`\texttt` command,
rather than displaying inappropriately as the algebraic expression
$2e-1$. 

\subsubsection{\texttt{\textbackslash ensuremath},\texttt{ \$},\texttt{ \textbackslash (},\texttt{
\textbackslash )},\texttt{ \textbackslash{[}},\texttt{ \textbackslash{]}}}

Should \verb`\ensuremath` be included in a formula for evaluation
(but why?) it is digested without demur, irrespective of whether explicit
math delimiters are present or not. More generally, should math delimiters
(through some momentary oversight) be used both within and outside
an \verb`\eval` command, the command is processed as if only the
outside environment is involved; the inner delimiters are ignored:
\begin{centred}
\verb`$ \eval{\[ -4^2 \]} $` $\Longrightarrow$ $ \eval{\[ -4^2 \]} $.
\end{centred}

\section{Error messages }

\label{sec:evalErrors}There are two kinds of error in \texttt{numerica}:
those in the underlying \LaTeX{} which are reported in the \LaTeX{}
log, shown on the terminal, and generally halt compilation, and specifically
\texttt{numerica}-related errors which do not halt compilation and
produce messages displayed in the pdf where one would expect the result
of the calculation to be. The original reason for doing things this
way was to enable \texttt{numerica} to be used effectively with the
instant preview facility of the document processor \LyX . More philosophically,
one might view such errors as similar to errors of grammar or spelling
mistakes in text. It is not clear that they should halt compilation.
Hence strictly \texttt{numerica}-related errors leave brief messages
in the pdf at the offending places.

Before discussing specific error messages, note that there is a debug
facility (of a sort) discussed below in §\ref{subsec:settingsDebug}. 

Error messages are in two parts: a \emph{what} part and a \emph{where}
part.

\subsection{Mismatched brackets}

\label{subsec:errorsMismatched-brackets}An unmatched left parenthesis
or other left bracket (in this case a missing right parenthesis) usually
results in a \texttt{numerica} error:
\begin{centred}
\verb`$\eval{\sin(\pi/(1+x)}[x=1]$` $\Longrightarrow$ $\eval{\sin(\pi/(1+x)}[x=1]$
\end{centred}
For the same error in the vv-list, the what-part remains unchanged
but the where-part is altered:
\begin{centred}
\verb`$\eval{ 1+y }[x=1,y=\sin(\pi/(1+x)]$` $\Longrightarrow$ $\eval{ 1+y }[y=\sin(\pi/(1+x),x=1]$
\end{centred}
The \emph{what} message is the same; the \emph{where} is different.

An unmatched right parenthesis or other right bracket (in this case
a missing \emph{left} parenthesis) usually results in a similar \texttt{numerica}
error:
\begin{centred}
\verb`$\eval{2((x+y)/(y+z)))^2}[x=1,y=2,z=3]$` $\Longrightarrow$
\eval{2((x+y)/(y+z)))^{2}}[x=1,y=2,z=3]
\end{centred}
But note that an unmatched modifier like \verb`\left` or \verb`\right`
is a \LaTeX{} error and is caught by \LaTeX{} before \texttt{numerica}
can respond and so results in a terminal and logfile message. 

\subsection{Unknown tokens}

An `Unknown token' message can arise in a number of ways. If an
expression involves a number of variables, some of which depend on
others, their order in the vv-list matters: 
\noindent \begin{center}
\verb`$\eval{\tfrac12 vt}[t=2,v=gt,g=9.8]$` $\Longrightarrow$ \eval{\tfrac{1}{2}vt}[t=2,v=gt,g=9.8] 
\par\end{center}

The vv-list is evaluated from the \emph{right} so that in this example
the variable \texttt{v} depends on a quantity \texttt{t} that is not
yet defined. Hence the message. The remedy is to move \texttt{t} to
the right of \texttt{v} in the vv-list.

Similarly, if we use a variable in the formula that has not been assigned
a value in the vv-list, we again get the `Unknown token' message,
but this time the location is the formula:
\begin{centred}
\verb`$\eval{\pi r^2h}[r=3]$` $\Longrightarrow$ \eval{\pi r^{2}h}[r=3]
\end{centred}
The remedy obviously is to assign a value to \texttt{h} in the vv-list\texttt{.}

The same message will result if a mathematical operation or function
is used that has not been implemented in \texttt{numerica}:
\begin{centred}
\verb`$\eval{u \bmod v }[v=7,u=3]$` $\Longrightarrow$ \eval{u\bmod v}[v=7,u=3] 
\end{centred}
A missing comma in the vv-list will generally result in an unknown
token message:
\begin{centred}
\verb`$\eval{axy}[a=3 y=2,x=1]$` $\Longrightarrow$ \eval{axy}[a=3y=2,x=1] 
\end{centred}
Because of the missing comma, \verb`numerica` assumes \verb`a` has
the `value' \verb`3y=2`, an expression which it then tries to evaluate,
but the variable \verb`y` in this expression has not been assigned
a value, which generates the message.

\emph{Extra} commas in the vv-list should cause no problems:
\begin{centred}
\verb`$\eval{axy}[,a=3,,y=2,x=1,]$` $\Longrightarrow$ $\eval{axy}[,a=3,,y=2,x=1,]$ 
\end{centred}
The presence of multi-token variables can also cause an unknown token
message if the check for such variables is turned off; see §\ref{subsec:settingsMultitokSwitch}.

\subsection{Overlooked value assignments}

Perhaps if one is evaluating a formula with a number of variables
and assigning different experimental values to them to see the effect,
a variable might be overlooked:
\begin{centred}
\verb`$\eval{axy}[a=3,y=,x=1]$` $\Longrightarrow$ \eval{axy}[a=3,y=,x=1] 
\end{centred}
In the example the variable \verb`y` has been overlooked. The remedy
is obvious – assign a value to it.

\subsection{Integer argument errors}

\label{subsec:errorsIntegerArgs}Some functions require integer arguments
– factorials, the second argument of a binomial coefficient, and (in
\texttt{numerica}) $n$th roots using the optional argument of \texttt{\textbackslash sqrt};
also summation and product variables. If integers are explicitly entered
for these arguments there is no problem, but if the value of the argument
is the result of a calculation, rounding errors require thinking about.
What accumulation of rounding errors is \emph{too} much so that the
result of the calculation \emph{cannot} be considered an integer?
\texttt{numerica} is generous: in the default setup, if a calculation
rounds to an integer at rounding value $14$ the result of the calculation
is considered an integer (obviously, the value resulting from the
rounding). Since \texttt{l3fp} works to $16$ significant figures
that gives ample room for rounding errors to `get lost in' and be
ignored, while still ruling out such things as (recall the example
in §\ref{subsec:evalBoolean-output}),
\begin{centred}
\verb`\eval{\[ \sum_{n=1}^N n \]}[N=1/0.0123456789]` $\Longrightarrow$
\eval{\[ \sum_{n=1}^N n \]}[N=1/0.0123456789]
\end{centred}
where $N$ differs from $81$ not until the seventh decimal place. 

The default rounding value of $14$ for `int-ifying' calculations
can be changed: see §\ref{subsec:defaultsIntifyingRounding}.

\subsection{Comparison errors}

Should a user try to make a comparison using a combination like \verb`>=`
rather than \verb`\geq`, \texttt{numerica} admonishes like this: 
\begin{centred}
\verb`$\eval{ e^\pi >= \pi^e }[?]$` $\Longrightarrow$ $\eval{ e^\pi >= \pi^e }[?]$
\end{centred}
(The relation is true by the way.) The same error is generated by
other multi-token comparisons. They are used in programming languages,
yes, but \emph{not} in mathematics.

\subsection{Invalid base for \texttt{\textbackslash log}}

ISO recommends using \verb`\log` only with a subscripted base specified.
Otherwise how is one to know whether the base is \verb`e` or $10$
or $2$ or whatever? Nonetheless \texttt{numerica} assumes that when
\verb`\log` is used unsubscripted, the base is 10. Suppose you want
to make $12$ the base, but forget to put braces around the $12$:
\begin{centred}
\verb`$\eval{ \log_12 1728 }$` $\Longrightarrow$ $\eval{ \log_12 1728 } $
\end{centred}
Here, \texttt{numerica} has taken \texttt{1} as the base (and $21728$
as the argument) of the logarithm and responds accordingly.

\subsection{\texttt{l3fp} errors}

Some errors arising at the \texttt{l3fp} level are trapped and a message
displayed. 

\subsubsection{Dividing by zero}
\begin{centred}
\verb`$\eval{1/\sin x}[x=0]$` $\Longrightarrow$ \eval{1/\sin x}[x=0] 
\end{centred}
Note however that \verb`$\eval{1/\sin x}[x=\pi]$` $\Longrightarrow\,\eval{1/\sin x}[x=\pi]$,
because of rounding errors in distant decimal places. No doubt this
is true for other functions as well.

\subsubsection{Invalid operation}

\label{subsec:errorsInverse-powers}Finding inverse integer powers
of \emph{positive} numbers should always be possible, but raising
a \emph{negative} number to an inverse power generates an error even
when – mathematically – it should not: 
\begin{centred}
\verb`\eval{$ (-125)^{1/3} $}` $\Longrightarrow$ \eval{$ (-125)^{1/3} $}
\end{centred}
This is a feature of floating point arithmetic. When a number is raised
to a rational power, say $p/q$ where $p$ and $q$ are non-zero integers,
then the result is the $p$th power of the $q$th root of the number.
Can a $q$th root be taken? If our floating point system used (for
ease of illustration) only $4$ significant digits, $p/q=1/3$ would
be the fraction $3333/10^{4}$, an odd numerator over an even denominator.
But a negative number does not possess an even ($10^{4}$th) root. 

Trying to evaluate a function like a factorial or square root or inverse
trig. function outside its domain of definition also produces this
error:
\begin{centred}
\verb`$\eval{\arccos x}[x=2]$` $\Longrightarrow$ \eval{\arccos x}[x=2]
\end{centred}
In this case the inverse cosine, which is defined only on the interval
$[-1,1]$, has been fed the value $2$. 

Trying to evaluate an expression that resolves to $0/0$ also produces
this message:
\begin{centred}
\verb`$\eval{\frac{1-y}{x-2}}[x=2,y=1]$` $\Longrightarrow$ \eval{\frac{1-y}{x-2}}[x=2,y=1]
\end{centred}

\subsubsection{Overflow/underflow}

The factorial (discussed in §\ref{subsec:evalFactorialsBinom}) provides
an example of overflow:
\begin{centred}
\verb`$\eval{3249!}$`\texttt{ }$\Longrightarrow$ \eval{3249!} 
\end{centred}
This is hardly surprising since
\begin{centred}
\verb`$\eval{3248!}[x]$` $\Longrightarrow$ $\eval{3248!}[x]$.
\end{centred}
There is a limit on the size of exponents that \texttt{l3fp} can handle.
A number in the form $a\times10^{b}$ must have $-10001\le b<10000$.
If this is not the case an overflow or underflow condition occurs.
As the examples show, an overflow condition generates a \texttt{numerica}
error. 

For underflow, where the number is closer to $0$ than $10^{-10001}$,
\texttt{l3fp} assigns a zero value to the quantity. \texttt{numerica}
accepts the zero value and the error is ignored.

\chapter{Settings}

\label{chap:Settings}A calculation is effected against a background
of default values for various quantities. For a particular calculation,
these values may not be appropriate; or you may have different preferences.
The way to change settings for a particular calculation is through
the settings option of \verb`\nmcEvaluate` discussed next. The way
to change a \emph{default} setting is by creating a configuration
file \texttt{numerica.cfg} discussed in §\ref{sec:settingsDefaults}.

\section{Settings option}

\label{sec:settingsOption}The second argument of the \verb`\nmcEvaluate`
command is the settings option, delimited by square brackets. This
option is a \emph{key=value} list, hence comma-separated. \emph{Key=value
}lists tend to be wordy. For back-of-envelope calculations one wants
to be able to `dash off' the calculation, hence the short, cryptic
nature of the keys. Most settings are generic, applicable not only
to \verb`\nmcEvaluate` but also to other commands that are available
if the packages \texttt{numerica-plus} or \texttt{numerica-tables}
are loaded; see §\ref{subsec:Related-packages}. 

\subsection{\textquoteleft Debug\textquoteright{} facility}

\label{subsec:settingsDebug}It is rather grandiose to call this a
debug facility, but if a calculation goes wrong or produces a surprising
result, \texttt{numerica} offers a means of examining various quantities
at some intermediate stages on the way to the final result. To use
the facility, enter 
\begin{lyxcode}
dbg~=~<integer>
\end{lyxcode}
into the settings option. (White space around the equals sign is optional.) 
\begin{itemize}
\item \texttt{dbg=0 }turns off the debug function, displays the result or
error message (this is the default);
\item \texttt{dbg=1 }equivalent to \texttt{dbg=2{*}3{*}5{*}7};
\end{itemize}
\begin{table}[t]
\centering
\noindent \centering{}\caption{Settings options}
\noindent \begin{center}
\begin{tabular}{ll>{\raggedright}p{4cm}>{\raggedright}p{4cm}}
\toprule 
{\small key} & {\small type} & {\small meaning} & {\small default}\tabularnewline
\midrule
{\small\texttt{dbg}} & {\small int} & {\small debug `magic' integer} & {\small\texttt{0}}\tabularnewline
{\small\texttt{view}} &  & {\small equivalent to }{\small\texttt{dbg=1}} & \tabularnewline
{\small\texttt{\textasciicircum}} & {\small char} & {\small exponent mark for sci. notation input} & {\small\texttt{e}}\tabularnewline
{\small\texttt{xx}} & {\small int (0/1)} & {\small multi-token variable switch} & {\small\texttt{1}}\tabularnewline
{\small\texttt{()}} & {\small int (0/1/2)} & {\small trig. arg. parsing} & {\small\texttt{0}}\tabularnewline
{\small\texttt{o}} & {\small int (0/1)} & {\small degree switch for trig. functions} & {\small\texttt{1}}\tabularnewline
{\small\texttt{log}} & {\small num} & {\small base of logarithms for }{\small{\small\verb`\log`}} & {\small\texttt{10}}\tabularnewline
{\small\texttt{vv@}} & {\small int (0/1)} & {\small vv-list calculation mode} & {\small\texttt{0}}\tabularnewline
{\small\texttt{vvd}} & {\small token(s)} & {\small vv-list display-style spec.} & {\small\texttt{\{,\}\textbackslash mskip 12mu  6mu minus 9mu(vv)}}\tabularnewline
{\small\texttt{vvi}} & {\small token(s)} & {\small vv-list text-style spec.} & {\small\texttt{\{,\}\textbackslash mskip 36mu minus 24mu(vv)}}\tabularnewline
{*} &  & {\small suppress equation numbering if }{\small\texttt{\textbackslash\textbackslash}}{\small{}
in }{\small\texttt{vvd}} & \tabularnewline
{\small\texttt{p}} & token(s) & {\small punctuation (esp. in display-style)} & {\small\texttt{,}}\tabularnewline
{\small\texttt{reuse}} & {\small int} & {\small form of result saved with }{\small{\small\verb`\nmcReuse`}} & {\small\texttt{0}}\tabularnewline
\bottomrule
\end{tabular}
\par\end{center}
\end{table}

The `magic' integers are the following primes and their products
\begin{itemize}
\item \texttt{dbg=2} displays the vv-list after multi-token variables have
been converted to their single token form, \texttt{\textbackslash nmc\_a},
\texttt{\textbackslash nmc\_b}, etc.;
\item \texttt{dbg=3} displays the formula after multi-token variables have
been converted to their single token form;
\item \texttt{dbg=5} displays the stored variables and their \emph{evaluated}
values (\texttt{dbg=2} lists the values as expressions);
\item \texttt{dbg=7} displays the formula after it has been fp-ified but
before it has been fed to \texttt{l3fp} to evaluate; 
\begin{itemize}
\item should the formula successfully evaluate, the result of the evaluation
is also displayed (but without any formatting).
\end{itemize}
\end{itemize}
To display two or more of the debug elements simultaneously, use the
product of their debug numbers for the magic integer. This can be
entered either as the multiplied-out product, or as the `waiting
to be evaluated' product with asterisks (stars) between the factors.
Thus \texttt{dbg=6} or \verb`dbg=2*3` display both the vv-list and
formula after multi-token variables have been converted to single
token form; \texttt{dbg=10} or \verb`dbg=2*5` display both the vv-list
after multi-token variables have been converted to single token form
and the recorded variables with their evaluated values. And similarly
for the other magic integers listed. For other integers, if they are
divisible by $2$ or $3$ or $5$ or $7$, they will display the corresponding
component. Both \texttt{dbg=210} and \verb`dbg=2*3*5*7` display all
four elements, but rather than remembering this product, it suffices
to put \texttt{dbg=1}. This is equivalent and displays all elements.

The debug option uses an \verb`aligned`\texttt{ }or \verb`align*`
environment to display its wares, depending on the presence or absence
of math delimiters around the \verb`\eval` command. The following
uses \verb`align*` and shows how multi-token variables are handled,
how a chain of comparisons is evaluated (§\ref{subsec:evalBoolean-output})
and how formatting instructions in the number-format option are ignored
in the debug display:
\begin{verbatim}
    \eval[dbg=1]{ a < 2a' < 3a'' }
      [a=\pi,a'=\phi,a''=e\gamma][4???]
\end{verbatim}
$\Longrightarrow$     \eval[dbg=1]{ a < 2a' < 3a'' }
      [a=\pi,a'=\phi,a''=e\gamma][4???]

\noindent The various items are displayed in chronological order.
First comes the vv-list after conversion of multi-token to single-token
variables, then the formula in those single-token variables; these
are created essentially at the same time. Next the stored values of
the variables are displayed. These are the values \emph{after }evaluation.
The fourth element both in the display and chronologically is the
fp-ified formula. Often this can be a thicket of parentheses, especially
if unary functions or fractions are involved. The final element of
both the display and chronologically is the result from evaluating
the formula. This is shown only if $7$ is a factor of the \texttt{dbg}
integer, and there is no error. Despite the appearance of \verb`???`
in the number-format option, the result displays as 1. Results are
never rounded or formatted in the debug display, although as is apparent
here, the rounding number \verb`4` is used in the comparisons.

When interpreting the fp-form, differences in the ways \texttt{numerica}
and \texttt{l3fp} read formulas can lead to more or less parentheses
than seem strictly necessary. In particular be aware that in \texttt{l3fp}
function calls bind most tightly so that, for example, \verb`sin 2pi`
evaluates not to zero but to $(\sin2)\times\pi$, and \verb`sin x^2`
evaluates to $(\sin x)^{2}$. \verb`numerica` takes care of the former
by inserting extra parentheses and exploits the latter by not inserting
parentheses:
\begin{verbatim}
    \eval[dbg=1]{ \sin 2x \cos^2 y }
         [x=\pi/12,y=\pi/4]
\end{verbatim}
$\Longrightarrow$ \eval[dbg=1]{ \sin 2x \cos^2 y }[x=\pi/12,y=\pi/4]Finally,
note that those mathematical operations that have no direct representation
in \verb`l3fp` contribute only their value to the fp-form. This applies
to sums and products, double factorials, partly to binomial coefficients,
and also to \verb`\eval` and other commands when nested one within
another (see Chapter~\ref{chap:Nesting}). The following (ridiculous)
example illustrates the matter:
\begin{verbatim}
    \eval[dbg=1]{\[ 
      \sum_{n=1}^k n + \binom{2k}{m} - \frac1{4k} + 
        \prod_{n=2}^k (1-1/n) + m!! \]}[m=6,k=5]
\end{verbatim}
$\Longrightarrow$     \eval[dbg=1]{\[ 
      \sum_{n=1}^k n + \binom{2k}{m} - \frac1{4k} + 
        \prod_{n=2}^k (1-1/n) + m!! \]}[m=6,k=5]

\noindent ($0$\textdegree{}~C in kelvin!) In the \verb`fp-form`
line, the various contributions to the overall result are displayed
simply as numbers because \verb`l3fp` does not (at least as yet)
handle these elements natively.

\subsection{Negative \texttt{dbg} values}

Negative \texttt{dbg} values are possible: \texttt{dbg=-2}, \texttt{dbg=-3},
etc. (and \texttt{dbg=-1} meaning \texttt{dbg=-210}) have exactly
the same effects as the corresponding positive values except for some
details of display. The display for positive \texttt{dbg} values is
the one evident in the examples above. Lines wrap, the left margin
is not indented and the display occupies the page width. For negative
\texttt{dbg} values, lines do not wrap, the left margin is indented
and the display occupies the text width. An example is presented in
§\ref{subsec:nestDebugging} below where the display for a nested
\verb`\eval` is significantly improved with a negative \texttt{dbg}
value.

\subsection{\texttt{view} setting}

Putting \texttt{dbg=1} may seem a little obscure in order to view
internal values of \verb`numerica`. In that case, simply writing
\verb`view` in the settings option will produce the same effect as
entering \verb`dbg=1`.

\subsection{Inputting numbers in scientific notation}

\label{subsec:settingsInputtingSciNotation}\emph{Outputting} numbers
in scientific notation is controlled by the final trailing argument
of the \texttt{\textbackslash eval} command. Such output is turned
off by default and needs to be explicitly ordered. Similarly, \emph{inputting}
numbers in scientific notation is turned off by default and needs
to be explicitly ordered. To turn it on, write
\begin{lyxcode}
\textasciicircum ~=~<char>
\end{lyxcode}
\noindent in the settings option, where \texttt{<char>} is any single
character, usually \texttt{e} or \texttt{d} or their upper-casings,
but not restricted to them: \texttt{\textasciicircum =@} for instance
is perfectly possible, and has the advantage over \texttt{e} or \texttt{d}
that it doesn't conflict with the use of the character as a variable
or constant. 
\begin{centred}
\verb`$ \eval[^=@]{ 1.23@-1 } $` $\Longrightarrow$ $ \eval[^=@]{ 1.23@-1 } $.
\end{centred}
With letters for the exponent mark – say \verb`d` or \verb`e` –
the problem is interpreting forms like \texttt{8d-3} or \texttt{2e-1}.
Does such a form denote a number in scientific notation or an algebraic
expression? In \texttt{numerica}, if the settings option shows \texttt{\textasciicircum =d},
then a form like \texttt{8d-3} is treated as a number in scientific
notation. Similarly for \texttt{e} or any other letter used as the
exponent marker for the input of scientific numbers. (But only one
character can be so used at a time.) Note that the number \emph{must}
start with a digit: \verb`e-1` for instance does not, and will be
treated as an algebraic expression involving the exponential constant:
\begin{centred}
\verb`$ \eval[^=e]{ x+e-1 }[x=1] $` $\Longrightarrow$ $ \eval[^=e]{ x+e-1 }[x=1] $ 
\end{centred}
but
\begin{centred}
\verb`$ \eval[^=e]{ x+1e-1 }[x=1] $` $\Longrightarrow$ $ \eval[^=e]{ x+1e-1 }[x=1] $. 
\end{centred}
A problem of appearance arises if scientific numbers appear in the
vv-list or formula and either is displayed in the result. A number
like \verb`2e-1` will display as $2e-1$, as if it were an algebraic
expression. In version 1 of \verb`numerica` the cure was to wrap
\verb`2e-1` in a \verb`\text` or \verb`\mbox` command. In version
2 of \verb`numerica` the behaviour of \verb`\text` and \verb`\mbox`
has been re-thought; see §\ref{subsec:Text-mbox-fonts}. Their contents
are now invisible to the \verb`\eval` command. The solution is to
wrap \verb`2e-1` in a \verb`\textrm` or \verb`\textsf` or \verb`\texttt`
command. These commands were not recognized by \verb`\eval` in version
1 but \emph{are} in version 2:
\begin{centred}
\verb`\eval[^=e]{$  5x $ }[x=\texttt{2e-1}]` $\Longrightarrow$ \eval[^=e]{$  5x $ }[x=\texttt{2e-1}]
,

\verb`\eval[^=e]{$ 5\texttt{2e-1} $ }` $\Longrightarrow$ \eval[^=e]{$  5(\texttt{2e-1}) $ }
.
\end{centred}
If you use a particular character as the exponent marker for inputting
numbers in scientific notation, it is good practice \emph{not} to
use that character as a variable, not because it will cause an error
but because it makes expressions harder to read. 

\subsection{Multi-token variables}

\label{subsec:settingsMultitokSwitch}Variables need not consist of
a single character or token (like $x$ or $\alpha$). Multi-token
symbols like $x'$ or $t_{i}$ or $var$ are perfectly acceptable.
For its internal operations, \texttt{numerica} converts such multi-token
names to single tokens (as discussed in §\ref{subsec:evalVariableNames}).
This conversion takes time. Even if there are no multi-token variables
used at all, \texttt{numerica} still needs to check that that is so.
There is a setting that allows a user to turn off or turn on the check
for such variables by entering
\begin{lyxcode}
xx~=~<integer>
\end{lyxcode}
into the settings option. If \texttt{<integer>} is \texttt{0}, the
check for (and conversion of) multi-token variables is turned off;
if \texttt{<integer>} is \texttt{1} (or any other\emph{ non-zero}
integer), the check, and conversion if needed, goes ahead. By default,
checking for multi-token variables and converting them if found is
turned \emph{on}. (The name for the key, \texttt{xx}, is chosen because
\texttt{x} is the most familiar variable of all, introduced in elementary
algebra, and doubling it like this suggests multi-token-ness.)

If checking is turned off when a multi-token variable is present,
an error results. We don't need to enter \texttt{xx=1} in the first
of the following examples because the check for multi-token variables
is on by default. Explicitly turning it off in the second produces
an error.
\begin{centred}
\verb`\eval{$ x_0^{\,2} $}[x_0=5]` $\Longrightarrow$ \eval{$ x_0^{\,2} $}[x_0=5],\medskip{}
\verb`\eval[xx=0]{$ x_0^{\,2} $}[x_0=5]` $\Longrightarrow$ \eval[xx=0]{$ x_0^{\,2} $}[x_0=5]
\end{centred}

\subsection{Parsing arguments of trigonometric functions}

This setting allows a wider range of arguments to trigonometric  functions
to be parsed (think Fourier series) without needing to insert extra
parentheses in order for them to be read correctly by \verb`\eval`;
see §\ref{subsec:parseTrigFns}.

\subsection{Using degrees rather than radians}

\label{subsec:settingsDegrees}You may find it more convenient to
use degrees rather than radians with trigonometric functions. This
can be switched on simply by entering a lowercase \texttt{o} in the
settings option. (The author hopes the charitable eye sees a degree
symbol in the  \texttt{o}.) Thus
\begin{centred}
\verb`\eval[o]{$ \sin 30 $}` $\Longrightarrow$ \eval[o]{$ \sin 30 $},

\verb`\eval[o]{$ \arcsin 0.5 $}` $\Longrightarrow$ \eval[o]{$ \arcsin 0.5 $}.
\end{centred}
This is a \verb`0/1` switch, \verb`0` signifying \verb`off` or
`don't use degrees', \verb`1` signifying \verb`on` or `do use
degrees'. Although the \verb`o` default is \verb`1`, out-of-the-box
\verb`numerica` assumes radians are being used. Thus if \verb`o`
is absent from the settings option of an \verb`\eval` command, the
out-of-the-box setting prevails and radians are used, but if \verb`o`
is present, it is equivalent to \verb`o=1`. To explicitly turn off
the use of degrees requires the full setting, \verb`o=0`. 

If you want to change the out-of-the-box setting you need to put the
line \verb`use-degrees = 1` into a configuration file; see §\ref{sec:settingsDefaults}.

\subsection{Specifying a logarithm base}

\label{subsec:settingsLogBase}If you wish to use \verb`\log` without
a subscripted base in a particular calculation, then add an entry
like 
\begin{lyxcode}
log~=~<positive~number>
\end{lyxcode}
where \verb`<positive number>`~$\ne1$ to the settings option of
the \verb`\eval` command. The \verb`<positive number>` does not
need to be an integer. It could be \verb`e` (if you object to writing
\verb`\ln`) but is more likely to be $2$ or another small integer;
$10$ is the default. If you want to use this changed base not for
one but most calculations, then add an entry with your choice of base
to a configuration file; see §\ref{sec:settingsDefaults}.

\subsection{Calculation mode}

\label{subsec:settingsRecalcMode}A variable may change in the course
of a calculation. This is certainly true of sums and products. If
a parameter in the vv-list depends on the variable then that parameter
will need to be recalculated, perhaps repeatedly, in the course of
a calculation. By entering either
\begin{lyxcode}
vv@~=~<integer>
\end{lyxcode}
or (as in version 1 of \verb`numerica`),
\begin{lyxcode}
vvmode~=~<integer>
\end{lyxcode}
in the settings option it is possible to turn on or off the ability
to repeatedly evaluate the vv-list; \verb`<integer>` here takes two
possible values, \texttt{0} or \texttt{1}. \texttt{vv@=0} (or \texttt{vvmode=0})
means the vv-list is evaluated once at the start of the calculation;
\texttt{vv@=1} (or\texttt{ vvmode=1}) means the vv-list is recalculated
every time the relevant variable changes.\footnote{In version 1 of \texttt{numerica} only the \texttt{vvmode} name for
this setting was available. To the author's eye, the \texttt{@} sign
seems sufficiently close to a symbol like $\circlearrowleft$, suggesting\texttt{
}redo or recalculate, that \texttt{vv@} is now preferred. The \texttt{@}
symbol is – universally? – available on keyboards and \texttt{vv@}
is only half as many keypresses as \texttt{vvmode}.}

For example, in a sum it may be desirable to place the summand, or
some part of it, in the vv-list. Since the summation variable obviously
changes during the course of the calculation, we need to enter \texttt{vv@=1}
in the settings option. Repeating an earlier sum (the seting \verb`p=.`
is discussed in §\ref{subsec:settingsPunctuation}), \medskip{}
\begin{verbatim}
    \eval[p=.,vv@=1]{\[ \sum_{k=1}^N f(k) \]}
      [N=100,f(k)=1/k^3,{k}=1][4] 
\end{verbatim}
$\Longrightarrow$      \eval[p=.,vv@=1]{\[ \sum_{k=1}^N f(k) \]}
        [N=100,f(k)=1/k^3,{k}=1][4]

As you can see, the summand \texttt{f(k)} has been given explicit
form in the vv-list – equated to \texttt{1/k\textasciicircum 3}.
That means we need to give a preceding value to \texttt{k} in the
vv-list to avoid an unknown token message, hence the rightmost entry.
But we don't want \texttt{k=1} appearing in the final display, so
we wrap \texttt{k} in braces (see §\ref{subsec:evalVvSuppressVars}).
Since the value \texttt{k=1} applies only to the first term in the
sum, to ensure it is not used for all terms, we enter \texttt{vv@=1}
in the settings option. This turns vv-recalculation mode on and ensures
\texttt{k=1} is overwritten by \texttt{k=2}, \texttt{k=3} and so on,
and the vv-list recalculated each time. The final result is the same
as before, although recalculating the vv-list at each step is a more
resource-hungry process. The difference may not be marked for this
example; with more complicated expressions it noticeably takes longer.

Because it is necessary to activate this switch when using \emph{implicit}
notations – like $f(k)$ in the example – rather than the explicit
form of the function in the main argument, it seems natural to call
\texttt{vv@=1} \emph{implicit }mode and \texttt{vv@=0} (the default)
\emph{explicit }mode. Most calculations are explicit mode – the vv-list
is evaluated only once.\emph{ }

\subsection{Changing the vv-list display format}

\label{subsec:settingsVvDisplayChangeLocal}In previous formulas with
variables the vv-list has been displayed following the result. It
is wrapped in parentheses following a comma followed by a space. These
formatting elements – comma, space, parentheses – can all be changed
with the settings option.

The default format specification is 
\begin{lyxcode}
\{,\}\textbackslash mskip~12mu~plus~6mu~minus~9mu(vv)
\end{lyxcode}
for a text-style display (an inline formula) and
\begin{lyxcode}
\{,\}\textbackslash mskip~36mu~minus~24mu(vv)
\end{lyxcode}
in a display-style context. The commas are wrapped in braces because
these are items in a comma-separated list. Both entries exhibit the
elements: punctuation (comma), preceding a variable space, preceding
the parenthesized vv-list (the \texttt{vv} placeholder). No full stop
is inserted after the closing parentheses because the \texttt{\textbackslash eval}
command may occur in the middle of a sentence (even in display style).
For inline use, the elasticity of the space becomes relevant when
\TeX{} is adjusting individual lines to fit sentences into paragraphs
and paragraphs into pages. The largest spacing that can be stretched
to is a quad, $18$~mu (mu~=~math unit), and the smallest that
can be shrunk to is a thin space, $3$~mu. In display style,\texttt{
}the largest spacing specified is the double quad, in line with the
recommendation in \emph{The \TeX{} Book}, Chapter~18, but this can
shrink to a single quad, for instance if the vv-list is heavily populated
with variables so that the evaluated result is pushed well to the
left by the vv-list. (But see below, §\ref{subsec:settings New-line-display}.)

If you want to change these defaults, enter in the settings option
\begin{lyxcode}
vvi~=~<new~specification>
\end{lyxcode}
to change the inline display and
\begin{lyxcode}
vvd~=~<new~specification>
\end{lyxcode}
to change the display-style display For example the settings
\begin{lyxcode}
vvi~=~\{,\}\textbackslash quad(vv)

vvd~=~\{,\}\textbackslash qquad(vv)
\end{lyxcode}
would give a comma (in braces since the settings option is a comma-separated
list) and a fixed space (of one or two quads) between the result and
the parenthesized vv-list. 

The vv-list itself in the display specification is represented by
the placeholder \texttt{vv}. If the \texttt{vv} is omitted from the
specification, then the vv-list will not appear at all:
\begin{centred}
\verb`\eval[vvi=?!]{$ \pi $}[\pi=3]` $\Longrightarrow$  \eval[vvi=?!]{$ \pi $}[\pi=3]
\end{centred}
More relevantly, it may well be the case that all variables in the
vv-list are suppressed (wrapped in braces). In that case nothing is
displayed. Compare the last example with
\begin{centred}
\verb`\eval[vvi=?!]{$ \pi $}[{\pi}=3]` $\Longrightarrow$ \eval[vvi=?!]{$ \pi $}[{\pi}=3]
\end{centred}
and
\begin{centred}
\verb`\eval[vvi=?!]{$ \pi $}` $\Longrightarrow$ \eval[vvi=?!]{$ \pi $}
\end{centred}
See also the punctuation setting below, §\ref{subsec:settingsPunctuation}.

\subsection{Displaying the vv-list on a new line }

\label{subsec:settings New-line-display}Display of a long formula
with many variables, hence a full vv-list, may not fit comfortably
on a line. In an earlier example I used Brahmagupta's formula to calculate
the area of a triangle. It squeezed onto a line. I shall now use his
formula for the area of a cyclic quadrilateral:
\[
A=\sqrt{(s-a)(s-b)(s-c)(s-d)}.
\]
The extra side (quadrilateral as against triangle) means there is
a further variable to accommodate, not only in the formula but also
in the vv-list. In the following example, the cyclic quadrilateral
is formed by a 45-45-90 triangle of hypotenuse 2 joined along the
hypotenuse to a 30-60-90 triangle. The sides are therefore $\surd2,\surd2,\surd3,1$.
Adding the areas of the two triangles, the area of the quadrilateral
is $A=1+\tfrac{1}{2}\surd3$, or in decimal form, \verb`$\eval{1+\tfrac12\surd3}$`
$\Longrightarrow$ $\eval{1+\tfrac12\surd3}$. Let's check with Brahmagupta's
formula:
\begin{verbatim}
    \eval[p=.,vvd={,}\\(vv),*]
      {\[ \sqrt{(s-a)(s-b)(s-c)(s-d)} \]}
        [s=\tfrac12(a+b+c+d),
           a=\surd2,b=\surd2,c=\surd3,d=1] 
\end{verbatim}
$\Longrightarrow$     \eval[p=.,vvd={,}\\(vv),*]
      {\[ \sqrt{(s-a)(s-b)(s-c)(s-d)} \]}
        [s=\tfrac12(a+b+c+d), 
           a=\surd2,b=\surd2,c=\surd3,d=1] 

\noindent The values agree. The point to note here is the\texttt{
vvd=\{,\}\textbackslash\textbackslash (vv)} and the \texttt{{*}}
in the settings option. The \texttt{\textbackslash\textbackslash}
in a specification for \texttt{vvd} acts as a trigger for \texttt{numerica}
to replace whatever math delimiters are enclosed by the \verb`\eval`
command with a \verb`multline` environment. As you can see, the specification
inserts a comma after the formula and places the parenthesized vv-list
on a new line. The star \texttt{{*}} if present suppresses equation
numbering by turning the \verb`multline` into a \verb`multline*`
environment. 

Things to note in the use of\texttt{ \textbackslash\textbackslash}
in a \texttt{vvd} specification are that
\begin{itemize}
\item it applies only to the \texttt{vvd} specification, not the \texttt{vvi}
spec.;
\item it applies only when\emph{ }\verb`\eval`\emph{ wraps around }a math
environment of some kind; 
\item it has no effect when the \verb`\eval` command is used \emph{within}
a math environment when the presentation of the result is of the form
\emph{result, vv-list}. The formula is not displayed and so the pressure
on space is less and the `ordinary' vv-list specification is used.
\end{itemize}

\subsection{Punctuation}

\label{subsec:settingsPunctuation}The \verb`\eval` command can be
used within mathematical delimiters or it can be wrapped around mathematical
delimiters. The latter gives a \emph{formula=result} style of display
automatically, which is convenient. One doesn't need to write the
\emph{formula=} part of the expression, but it causes a problem when
\verb`\eval` wraps around a display-style or similar environment:
how to display a following punctuation mark? For an inline display
we can simply follow the \verb`\eval` command with the appropriate
punctuation, for instance: \verb`\eval{$ 1+1 $}.` $\Longrightarrow$
\eval{$ 1+1 $}. But with \verb`\[ \]` delimiters used \emph{within}
the \verb`\eval` command a trailing fullstop will slide off to the
start of the next line, since it is beyond the closing delimiter.
We want it to display as if it were the last element \emph{before}
the closing delimiter.

Explicitly putting it there – \verb`\eval{\[ 1+1. \]}` – means the
punctuation mark becomes part of the formula. Potentially \texttt{numerica}
then needs to check not just for a fullstop but also other possible
punctuation marks like comma, semicolon, perhaps even exclamation
and question marks. All these marks have roles in mathematics or \texttt{l3fp}.
Including them in the formula means distinguishing their punctuation
role from their mathematical role and can only cause difficulties
and slow evaluation.

Instead, \texttt{numerica} uses the setting 
\begin{lyxcode}
p~=~<char(s)>~
\end{lyxcode}
to place the {\ttfamily\verb`<char(s)>`} after the result
but within the environment delimiters. The default punctuation mark
is the comma so that simply entering \texttt{p} will produce a comma
in the appropriate place. This saves having to write \texttt{p=\{,\}}
as would otherwise be required, since the settings option is a \emph{comma}-separated
list.

Nor is one limited to a single punctuation mark: 
\begin{centred}
\verb`\eval[p=\ (but no 8!)]{\[ \frac{1}{81} \]}[9]` $\Longrightarrow$
\eval[p=\ \text{(but no 8!)}]{\[ \frac{1}{81} \]}[9]
\end{centred}

\subsection{Reuse setting}

This setting determines whether the entire display or only the numerical
result is saved to file with the \verb`\nmcReuse` command. See below,
§\ref{subsec:supplReuseEvalSetting}.

\section{Infinite sums and products}

\label{sec:settingsInfiniteSumsProds}There are ways of tweaking various
default settings to nudge infinite sums and products to a correct
limit. These tweaks are applied via the settings option of the \verb`\eval`
command. 

The normal convergence criterion used by \texttt{numerica} to determine
when to stop adding/multiplying terms in an infinite sum/product is
\emph{when the next term added/multiplied leaves the total unaltered
when rounded to 2 more digits than the specified rounding value.}
Suppose $T_{k}$ is the sum/product after the inclusion of $k$ terms,
and $r$ is the rounding value. Denote $T_{k}$ rounded to $r$ figures
by $\left(T_{k}\right)_{r}$. \emph{The infinite sum or product stops
at the $(k+1)$th term (and the value is attained at the $k$th term)
when }$\left(T_{k+1}\right)_{r+2}=\left(T_{k}\right)_{r+2}$. The
hope is that if this is true at rounding value $r+2$ then at rounding
value $r$ the series or product will have attained a stable value
at that level of rounding.

For a series of monotonic terms converging quickly to a limit, this
stopping criterion works well, less so if convergence is slower, as
seen earlier with the infinite product for $\sinh1$. The criterion
can fail completely when terms behave in a non-monotonic manner. Terms
of a Fourier series, for example, may take zero values; the criterion
is necessarily satisfied but the series may still be far from its
limit. In a product the equivalent would be a term taking unit value.
Sometimes the initial terms of series or products are `irregular'
and take these `stopping' values meaning sum or product would stop
after only one or two additions/multiplications and far from any limit. 

\begin{table}[t]
\centering
\noindent \centering{}\caption{Settings for infinite sums \& products}\label{tab:settingsSumsProducts}
\noindent \begin{center}
{\small{}%
\begin{tabular}{ll>{\raggedright}p{4cm}l}
\toprule 
key & type & meaning & default\tabularnewline
\midrule 
\texttt{S+} & int & extra rounding for stopping criterion & \texttt{2}\tabularnewline
\texttt{S?} & $\text{int}\ge0$ & stopping criterion query terms for sums & \texttt{0}\tabularnewline
\texttt{P+} & int & extra rounding for stopping criterion & \texttt{2}\tabularnewline
\texttt{P?} & $\text{int}\ge0$ & stopping criterion query terms for products & \texttt{0}\tabularnewline
\bottomrule
\end{tabular}}
\par\end{center}
\end{table}
To cope with these possibilities, \texttt{numerica} offers two settings
for sums, two for products, summarized in Table~\ref{tab:settingsSumsProducts}.
These are entered in the settings option of the \verb`\eval` command.
\begin{itemize}
\item \texttt{S+=<integer> }or \texttt{P+=<integer>} additional rounding
on top of the specified (or default) rounding for the calculation;
default = $2$
\begin{itemize}
\item the larger the additional \texttt{<integer>} is, the more likely that
sum or product has attained a stable value at the specified rounding
$r$ 
\end{itemize}
\item \texttt{S?=<integer${}\,\mathtt{\ge0}$> }or \texttt{P?=<integer${}\,\mathtt{\ge0}$>}
the number of final terms to query after the stopping criterion has
been achieved to confirm that it is not an `accident' of particular
values; default = $0$
\begin{itemize}
\item a final few terms to be summed/multiplied and the rounded result after
each such operation to be compared with the rounded result at the
time the stopping criterion was achieved. Suppose the additional rounding
(\texttt{S+} or \texttt{P+}) is $n$ on top of the specified rounding
$r$ and let the number of final checking terms be $m$. Suppose $T_{k_{0}}$
is the first term at which the stopping criterion is achieved: $\left(T_{k_{0}}\right)_{r+n}=\left(T_{k_{0}+1}\right)_{r+n}$.
What we require of the final query terms is that $\left(T_{k_{0}}\right)_{r+n}=\left(T_{k_{0}+1+j}\right)_{r+n}$
for $j=0,1,\ldots,m$.
\end{itemize}
\end{itemize}
Previously we found that the infinite product for $\sinh1$ with the
default settings gave the wrong value, $0.174$, deficient by $1$
in the last digit. We now have the means to tweak the stopping criterion
by increasing the additional rounding: 
\begin{verbatim}
    \eval[p,P+=3]{\[
      \prod_{k=1}^{\infty}
      \biggl(\frac{x^2}{k^2\pi^2} +1\biggr) 
    \]}[x=1][3] \nmcInfo{prod}.
\end{verbatim}
\noindent $\Longrightarrow$ \noindent \eval[p,P+=3]{\[
  \prod_{k=1}^{\infty}
  \biggl(\frac{x^2}{k^2\pi^2} +1\biggr) 
\]}[x=1][3]  \nmcInfo{prod}.

\noindent To obtain that last item of information (350 factors), I've
anticipated a little and used the command \verb`\nmcInfo` with the
argument \verb`prod`; see §\ref{sec:supplInfo}. The product now
produces the correct three-figure value, but it takes $350$ factors
to do so. 

Knowing how many terms or factors have been needed helps assess how
trustworthy the result from an infinite sum or product is. For example,
for the exponential series,
\begin{verbatim}
    \eval[p]{\[
      \sum_{k=0}^\infty \frac1{k!} 
    \]}[9] \nmcInfo{sum}.
\end{verbatim}
$\Longrightarrow$ \eval[p]{\[\sum_{k=0}^\infty \frac1{k!} \]}[9] \nmcInfo{sum}. 

To $9$ places of decimals, using the default value \texttt{S+=2},
the exponential series arrives at the right sum after only $15$ terms.
Convergence is rapid. We can trust this result (and it is in fact
the correct nine-figure value). By contrast, if we didn't know the
value of $\sinh1$ beforehand, noting the number of factors required
would make us justly cautious about accepting the result of the infinite
product calculation. 

\noindent\begin{minipage}[t]{1\columnwidth}%
\begin{shaded}%
One way to gain confidence in a result is to choose a possibly unrealistic
rounding value – say, the default $6$ for the infinite product –
then use \emph{negative} values for the extra rounding, \texttt{S+=-5},
\texttt{S+=-4}, \ldots{} , so that the stopping criterion applies at
rounding values $s$ of $6+(-5)=1$, one decimal place, $6+(-4)=2$,
two decimal places, and so on, but the result is always presented
to $6$ decimal places. You can then see how the $6$-figure results
behave relative to the number of terms it takes to meet the stopping
criterion. A little experimenting shows that for the infinite product
for $\sinh1$ the number of factors $N_{s}$ required at a stopping
rounding value $s$ increases in geometric proportion with a scale
factor of about $3$: $N_{s}\approx\text{const}\times3^{s}$. This
rapidly becomes large ($3^{4}=81,3^{5}=243\dots$). For the exponential
series on the other hand $N_{s}=4+s$, the number of terms increases
only slowly, in direct proportion to the stopping rounding value.
Similar experiments with the sums of inverse fourth, third and second
powers of the integers using \verb`\nmcInfo` to find how many terms
are required at each stopping rounding value, show that at least over
the rounding value range $1$ to $8$, for inverse fourth powers $N_{s}\approx\text{const}\times1.7^{s}$,
for inverse third powers $N_{s}\approx\text{const}\times2^{s}$ and
for inverse squares $N_{s}\approx\text{const}\times3^{s}$. All are
geometric rather than arithmetic progressions, but for inverse fourth
powers the scale factor ($\approx1.7$) is sufficiently small that
for these low values of $s$ the number of terms required does not
grow too quickly (e.g. $1.7^6\approx\eval{1.7^6}[0]$). It is a standard
result (Euler) that the series sums to $\pi^{4}/90$: \verb`$ \eval{ \pi^4/90 } $`
$\Longrightarrow$ $ \eval{ \pi^4/90 } $ to six places, and indeed,
with the default \texttt{S+=2},
\begin{centred}
\verb`\eval[p]{\[ \sum_{k=1}^\infty \frac1{k^4} \]}` $\Longrightarrow$
\eval[p=.]{\[ \sum_{k=1}^\infty \frac1{k^4} \]}
\end{centred}
\end{shaded}%
\end{minipage}

\subsection{Premature ending of infinite sums}

All the series considered so far have been monotonic. Trigonometric
series will generally not be so, nor even single-signed. 

Trigonometric sums are computationally intensive and so, for the following
example, I have specified a rounding value of 2. The series 
\[
\sum_{n=1}^{\infty}\frac{4}{n^{2}\pi^{2}}(1-\cos n\pi)\cos2\pi nt
\]
is the Fourier series for the triangular wave function \textbackslash\!/\!\textbackslash\!/\!\textbackslash\!/\!\textbackslash{}
\ldots{} of period 1, symmetric about the origin where it takes its
maximum value 1, crossing the axis at $t=0.25$ and descending to
its minimum $-1$ at $t=0.5$, before ascending to a second maximum
at $t=1$ (and so on). In the interval $[0,0.5)$ the series should
sum to $1-4t$. The problem is that the summand $\frac{4}{n^{2}\pi^{2}}(1-\cos n\pi)\cos2\pi nt$
vanishes both when $n$ is even and when $4nt$ is an odd integer.
If $t=0.1$ then $4nt$ is never an odd integer so the summand vanishes
only for $n$ even, every second term. We expect the result to be
$1-4\times0.1=0.6$.
\begin{verbatim}
    \eval[p]{\[
      \sum_{n=1}^{\infty}
        \frac{4}{n^{2}\pi^{2}}
        (1-\cos n\pi)\cos2\pi nt 
    \]}[t=0.1][2] \nmcInfo{sum}.
\end{verbatim}
$\Longrightarrow$ \eval[p]{\[
  \sum_{n=1}^{\infty}
    \frac{4}{n^{2}\pi^{2}}
    (1-\cos n\pi)\cos2\pi nt 
\]}[t=0.1][2]  \info{sum}.

\noindent Only one term? Of course – since for the second term $n$
is even, the term vanishes and the stopping criterion is satisfied.
The way around this problem is to query terms \emph{beyond} the one
where the stopping criterion is achieved, i.e., to set \texttt{S?}
to a nonzero value. We try \texttt{S?=1}: 
\begin{verbatim}
    \eval[p,S?=1]{\[
      \sum_{n=1}^{\infty}
        \frac{4}{n^{2}\pi^{2}}
        (1-\cos n\pi)\cos2\pi nt 
    \]}[t=0.1][2] \nmcInfo{sum}.
\end{verbatim}
$\Longrightarrow$ \eval[p,S?=1]{\[
  \sum_{n=1}^{\infty}
    \frac{4}{n^{2}\pi^{2}}
    (1-\cos n\pi)\cos2\pi nt 
\]}[t=0.1][2] \info{sum}.

\noindent\begin{minipage}[t]{1\columnwidth}%
\begin{shaded}%
\begin{wraptable}{o}{0.35\columnwidth}%
\centering{}\vspace{-5ex}
\caption{Finite sums}\label{tab:settingsFinite-sums}
\abovetopsep=2ex %
\begin{tabular}{cc}
\toprule 
$N$ & $\Sigma$\tabularnewline
\midrule
$63$ & $0.6001$\tabularnewline
$64$ & $0.6001$\tabularnewline
$65$ & $0.5999$\tabularnewline
$66$ & $0.5999$\tabularnewline
$67$ & $0.5999$\tabularnewline
\bottomrule
\end{tabular}\end{wraptable}%
Table~\ref{tab:settingsFinite-sums} lists the results of evaluating
the \emph{finite }sums from $n=1$ to $N$ for values of $N$ around
$65$. Since the specified rounding value is $2$ for the calculation,
the stopping criterion\emph{ }applies at a rounding value of $2$
more than that, $4$. Since $N=64$ is even, the summand for the $64$th
term is zero and the sum takes the same value as for $N=63$. The
$65$th term is the query term and the sum differs, so the summation
continues. The $66$th term vanishes, so the stopping criterion is
met. This time for the query term, the $67$th, the sum retains the
same $4$-figure value, and the summation stops. The result was attained
at the $65$th term. Should we be confident in the result? Increase
the number of query terms to $3$ (there is no point in increasing
\texttt{S?} to $2$ because of the vanishing of the even terms), the
sum stops after $113$ terms, with the same $0.6$ result.  \end{shaded}%
\end{minipage}

For a final example, consider the error function 
\[
\erf z=\dfrac{2}{\sqrt{\pi}}\int_{0}^{z}e^{-t^{2}}dt
\]
which can also be rendered as an infinite sum (\emph{HMF }7.1.5):
\[
\erf z=\sum_{n=0}^{\infty}(-1)^{n}\frac{z^{2n+1}}{n!(2n+1)}.
\]
(\verb`\erf` expanding to \verb`erf` has been defined in the preamble
to this document using \verb`\DeclareMathOperator`.) We calculate
this sum for $z=2$ to $10$ places of decimals. Although this is
an alternating series, it is obvious that the summand never vanishes
when $z\ne0$ as here. Hence there seems no need to change the default
value \texttt{S?=0}.
\begin{verbatim}
    \eval[p]{\[ 
      \frac2{\sqrt{\pi}}
        \sum_{n=0}^\infty(-1)^n
          \frac{z^{2n+1}}{n!(2n+1)}   
    \]}[z=2][10*] \nmcInfo{sum}.
\end{verbatim}
$\Longrightarrow$ \eval[p]{\[ 
  \frac2{\sqrt{\pi}}
    \sum_{n=0}^\infty(-1)^n
      \frac{z^{2n+1}}{n!(2n+1)}   
\]}[z=2][10*] \nmcInfo{sum}.

According to \emph{HMF }Table 7.1, this calculated value of $\erf2$
is correct to all $10$ places. But beyond $z=2$ errors will begin
to interfere with the result. Note that $26$ terms means $n=26$
was the last value of $n$ for which the summand was evaluated. (The
sum stops at the $26$th term, $n=25$, but the next term $n=26$
needs to be calculated for the stopping criterion.) Fortuitously,
$2^{2\times26+1}=2^{53}$ is the greatest power of $2$ that can be
\emph{exactly} rendered to the $16$ significant figures that \texttt{l3fp}
uses. But $n!$ exceeds the $16$-significant figure limit of \texttt{l3fp}
when $n>21$, so despite the 10-figure result, errors have already
begun to occur in the denominator of the summand and accrue in the
sum when $z=2$. For larger $z$ values the errors can only get worse
and at some point will render the calculated value worthless at any
meaningful rounding value. For example, when $z=7$ the sum apparently
`evaluates' to over $929$ whereas we know that 
\[
\erf z<\dfrac{2}{\sqrt{\pi}}\int_{0}^{\infty}e^{-t^{2}}dt=1.
\]


\subsection{Double sums or products}

Sums or products can be iterate d. For instance, the exponential function
can be calculated this way:
\begin{verbatim}
    \eval[p]
      {\[ \sum_{k=0}^{\infty}
        \prod_{m=1}^{k}\frac{x}{m} \]}[x=2]
\end{verbatim}
$\Longrightarrow$     \eval[p]
     {\[ \sum_{k=0}^{\infty}
       \prod_{m=1}^{k}\frac{x}{m} \]}[x=2]

\noindent which is \verb`\eval{$ e^2 $}` $\Longrightarrow\eval{\ensuremath{e^{2}}}$.

A second example is afforded by Euler's transformation of series (\emph{HMF~}3.6.27).
To calculate $e^{-1}$ we use 
\begin{verbatim}
    \eval[p]
      {\[ \sum_{n=0}^{\infty}
              \frac{(-1)^{n}}{n!} \]}[3] \info{sum}.
\end{verbatim}
$\Longrightarrow$ \eval[p]{\[ \sum_{n=0}^{\infty}\frac{(-1)^{n}}{n!} \]}[3] \info{sum}.

Following Euler, this series can be transformed to the form
\begin{verbatim}
    \eval[p,S?=1]{\[
      \sum_{k=0}^\infty \frac{(-1)^k}{2^{k+1}}
      \sum_{n=0}^k(-1)^n\binom kn \frac1{(k-n)!} 
    \]}[3] \nmcInfo{sum}.
\end{verbatim}
$\Longrightarrow$ \eval[p,S?=1]{\[ \sum_{k=0}^\infty \frac{(-1)^k}{2^{k+1}}\sum_{n=0}^k(-1)^n\binom kn \frac1{(k-n)!} \]}[3] \nmcInfo{sum}.

\noindent Note the setting \verb`S?=1`. Without it, the summation
stops after $1$ term, the $k=0$ term, because the $k=1$ term vanishes.
With \verb`S?=1` it takes $16$ terms of the \emph{outer }sum to
reach the stopping criterion. Since that sum starts at $0$, that
means that changing the upper limit from $\infty$ to $15$ should
give the same result – which it does – but it takes $\tfrac{1}{2}\times16\times17=136$
terms in total to get there, to be compared with the $9$ terms of
the earlier simpler sum, and the terms are more complicated. Obviously
such double sums are computationally intensive.

\section{Changing default values}

\label{sec:settingsDefaults}The settings option enables various settings
to be changed for an individual calculation. You may find yourself
wanting to make such changes sufficiently often that a change of default
value is a better plan than encumbering each calculation with a list
of settings.

The way to do that is to create a \emph{configuration file} named
\texttt{numerica.cfg} in a text editor. Its entries are of the form
\emph{key=value} followed by a comma, and for clarity preferably one
entry per line (although this is not essential).The key names are
noticeably more verbose than the corresponding keys of the settings
option.\emph{ }The possible keys are listed in Table~\ref{tab:settingsDefaults},
together with their current default values. 

\begin{table}[t]
\centering
\noindent \centering{}\caption{Default values, \texttt{\textbackslash eval} command}\label{tab:settingsDefaults}
\noindent \begin{center}
\begin{tabular}{ll}
\toprule 
{\small key} & {\small value}\tabularnewline
\midrule
{\small rounding} & {\small\texttt{6}}\tabularnewline
{\small pad} & {\small\texttt{0}}\tabularnewline
{\small output-sci-notation} & {\small\texttt{0}}\tabularnewline
{\small output-exponent-char} & {\small\texttt{e}}\tabularnewline
{\small\%} & \tabularnewline
{\small input-sci-notation} & {\small\texttt{0}}\tabularnewline
{\small input-exponent-char} & {\small\texttt{e}}\tabularnewline
{\small multitoken-variables} & {\small\texttt{1}}\tabularnewline
{\small use-degrees} & {\small\texttt{0}}\tabularnewline
{\small logarithm-base} & {\small\texttt{10}}\tabularnewline
{\small vv-display} & {\small\texttt{\{,\}\textbackslash mskip 36mu minus 24mu(vv)}}\tabularnewline
{\small vv-inline} & {\small\texttt{\{,\}\textbackslash mskip 12mu 6mu minus 9mu(vv)}}\tabularnewline
{\small\%} & \tabularnewline
{\small sum-extra-rounding} & {\small\texttt{2}}\tabularnewline
{\small sum-query-terms} & {\small\texttt{0}}\tabularnewline
{\small prod-extra-rounding} & {\small\texttt{2}}\tabularnewline
{\small prod-query-terms} & {\small\texttt{0}}\tabularnewline
\% & \tabularnewline
{\small intify-rounding} & {\small\texttt{14}}\tabularnewline
{\small eval-reuse} & {\small\texttt{0}}\tabularnewline
\bottomrule
\end{tabular}
\par\end{center}
\end{table}
Keys taking one of two possible values, \verb`0` (for \verb`false/off`)
or \verb`1` (for \verb`true/on`), are \verb`pad` (the result with
zeros), \verb`output-sci-notation`, \verb`input-sci-notation`, (check
for) \verb`multitoken-variables`, and \verb`use-degrees` (for trig.
functions).

The table is divided into four parts. 
\begin{itemize}
\item The top four rows concern elements that can be changed for individual
calculations with the trailing optional argument of \verb`\eval`:
rounding, padding with zeros, and outputting in scientific notation;
see §\ref{subsec:evalRoundingEtc}. 
\begin{itemize}
\item Note that to output the result always in scientific notation requires
two settings, first setting \texttt{output-sci-notation} to \texttt{1},
and then choosing a character to act as the exponent marker. Because
\texttt{l3fp} uses \texttt{e} for this character, \texttt{numerica}
has made \texttt{e} its default. But this option is turned off by
default (hence the \texttt{0} against this key).
\end{itemize}
\item The next block of rows concern general elements that can be changed
for individual calculations with the settings option of \verb`\eval`;
see §\ref{sec:settingsOption}. The key names are more expansive here
but the effect is the same. 
\begin{itemize}
\item Note that to input numbers in scientific notation requires two settings,
first setting \texttt{input-sci-notation} to \texttt{1}, and then
choosing a character to act as the exponent marker. Because \texttt{l3fp}
uses \texttt{e} for this character, \texttt{numerica} has made \texttt{e}
its default. The option is turned off by default (hence the \texttt{0}
against this key). 
\end{itemize}
\item The third block of rows concern default settings for infinite sums
and products. These correspond to the keys \texttt{S+}, \texttt{S?}
and \texttt{P+}, \texttt{P?} of the settings option that can be used
to tweak the behaviour of the stopping criterion for such sums or
products; see §\ref{sec:settingsInfiniteSumsProds}.
\item The last block is for `left-overs': specifying at what rounding
value a floating point should be considered an integer (see §\ref{subsec:defaultsIntifyingRounding}
below), and specifying what kind of result is saved to file when the
\verb`\nmcReuse` command is used (see §\ref{subsec:supplReuseEvalSetting}).
\end{itemize}
If you are dissatisfied with any of the default values listed, then
in a text editor create a new file called \texttt{numerica.cfg} and
assign \emph{your} values to the relevant keys. For instance, if you
find yourself working to $4$ figures, that rounding to $6$ is too
many, then make the entry \texttt{rounding~=}~4. If also you want
results always presented in proper scientific notation,\emph{ $d.d_{1}d_{2}d_{3}d_{4}\times10^{n}$},
then add a comma after \texttt{4} and enter on a new line (recommended
but not strictly necessary; the comma is the crucial thing), \texttt{output-sci-notation~=~1,}
(note the comma) and on another new line, \texttt{output-exponent-char~=~x}. 

Perhaps you also want a non-zero setting for the final query terms
for infinite sums and products. This makes sense if you are largely
dealing with non-monotonic series – like Fourier series. Even the
Euler transformation of the exponential series for $e^{-1}$ discussed
above required a non-zero \texttt{S?}. If you wish to make this change
then add a comma and on a new line add (for instance) \texttt{sum-query-terms~=~1,}
and again on a new line, \texttt{prod-query-terms~=}~1. If this
is all you wish to change, then no comma is necessary after this final
entry. Your newly created file should look something like
\begin{lyxcode}
rounding~~~~~~~~~~~~~=~4,

output-sci-notation~~=~1,

output-exponent-char~=~x,

sum-query-terms~~~~~~=~1,

prod-query-terms~~~~~=~1
\end{lyxcode}
The white spacing may be different; white space is ignored by \texttt{numerica}
when reading the file. Using it to align the equals signs helps \emph{us}
read the file. Note that the last entry, because it is the last entry,
lacks a comma. Now save the file with the name \texttt{numerica.cfg}.
This file will be read by \texttt{numerica} near the end of its loading
process. These settings will be \texttt{numerica}'s defaults for the
relevant keys. 

\subsection{Location of \texttt{numerica.cfg}}

Save, yes, but where to? If the new settings are likely to apply only
to your current document, then the document's directory is a sensible
place to put it and \texttt{numerica} will certainly find it there
since it is part of \LaTeX 3 file handling that file searches are
not limited to the \TeX{} distribution (including your personal texmf
tree) but also include the current document directory. But what happens
when you start working on another document? Will you remember to copy
\texttt{numerica.cfg} to its new location? That is why your \emph{personal
texmf tree} is a better place. 

\subsubsection{Personal texmf tree? }

\label{subsec:settingsPersonal-texmf-tree}This is a directory for
`waifs and strays' of the \TeX{} system that are not included in
the standard distributions like MiK\TeX{} or \TeX Live. Here you place
personal packages designed for your own particular circumstances.
These may include your own \TeX{} or \LaTeX{} package, say \texttt{mypackage.sty},
achieving some small or singular effect that doesn't warrant wider
distribution on CTAN. Here you might place configuration files for
other packages with your preferences (unless the package requires
some specific location). Here you can put your personal bibliography
files. 

Your personal texmf tree is structured like the standard MiK\TeX{}
or \TeX Live hierarchy but placed in another location so that there
is no chance of its being overwritten when packages in MiK\TeX{} or
\TeX Live are updated. But these distributions need to be alerted
to its existence. 

For example, in the MiK\TeX{} console, click on \textsf{Settings},
and then on the \textsf{Directories} tab of the resulting dialog.
Here you get to add your personal texmf hierarchy to the list of paths
that MiK\TeX{} searches, by clicking on the \textsf{+} button, browsing
to your texmf folder and selecting it. By using the up and down arrow
keys that the MiK\TeX{} console provides, ensure that it lies \emph{above
}the the entry for the main MiK\TeX{} tree. That way, files in your
personal texmf tree will be found first and loaded. Now go to the
\textsf{Tasks} menu and click on \textsf{Refresh the filename database}.
This will let MiK\TeX{} know what is held in your personal texmf tree.
Files there can then be used like standard \LaTeX{} packages.

\subsection{Rounding in \textquoteleft int-ifying\textquoteright{} calculations}

\label{subsec:defaultsIntifyingRounding}Factorials, binomial coefficients,
summation and product variables, and (in \texttt{numerica}) $n$th
roots from the \texttt{\textbackslash sqrt} command, all require
integer arguments. These integers may indeed be entered explicitly
as integers, but they can also be determined as the result of a calculation.
Rounding errors may mean the result is not an exact integer. How much
leeway should be allowed before it is clear that the calculation did
not give an integer result? In the default setup, \texttt{numerica}
is generous. A number is considered an integer if it rounds to an
integer when the rounding value is $14$. Since \texttt{l3fp} works
to $16$ significant figures this provides more than enough `elbowroom'
for innocuous rounding errors to be accommodated. If a calculation
does not round to an integer at a rounding value of $14$ then it
seems reasonable to conclude that it has \emph{really }not given an
integer answer, not just that rounding errors have accumulated. If
you want to change this `int-ifying' value for a particular calculation,
then add a line to \texttt{numerica.cfg} like
\begin{lyxcode}
intify-rounding~=~<integer>
\end{lyxcode}
Since \texttt{l3fp} works to $16$ significant figures, values of
{\ttfamily\verb`<integer>`} greater than $16$ are pointless.
Generally int-ifying rounding values will be less than but close to
$16$ (although when testing the code I used some ridiculous values
like $3$ or $4$). If other entries follow this one in the file,
then conclude the line with a comma.

\section{Parsing mathematical arguments}

\label{sec:Argument-parsing} A main aim of the \texttt{numerica}
package is to require minimal, preferably no, adjustment to the \LaTeX{}
form in which an expression is typeset in order to evaluate it. But
mathematicians do not follow codified rules of the kind programming
languages insist on when writing formulas – like parenthesizing the
arguments of functions, or inserting explicit multiplication signs
({*}) between juxtaposed terms. Hence the question of where the arguments
of mathematical functions end is acute. For a few functions \LaTeX{}
delimits the argument: think of \verb`\sqrt`, \verb`\frac`, \verb`\binom`;
also \verb`^`. But for functions like \verb`\sin` or \verb`\tanh`
or \verb`\ln`, unary functions, this is not so; nor is it for sums
and products, and comparisons.

Before discussing the parsing rules for different groups of functions,
I discuss the means \texttt{numerica} provides to handle exceptions
to those rules, when one \emph{does} need to make some adjustment
to a formula.

\subsection{The cleave commands \texttt{\textbackslash q} and \texttt{\textbackslash Q}}

\label{subsec:parseCleave}The word \emph{cleave} has two opposed
meanings: to adhere or cling to, and to split apart or separate. \texttt{numerica}
defines two commands, \verb`\q` and \verb`\Q` to achieve these opposite
effects. When a mathematical argument is being parsed, the \verb`\q`
command joins the next token to the argument (\emph{cleaves to});
the \verb`\Q` command severs the next token from the argument (\emph{cleaves
apart}). Neither command is added to the argument nor leaves a visible
trace in the output. 

Thus, without \verb`\q`, 
\begin{centred}
\verb`\eval{$ \sin(n+\tfrac12)(x-t) $}[n=3,x=t+\pi,t=1.234]` $\Longrightarrow$
\eval{$ \sin(n+\tfrac12)(x-t) $}[n=3,x=t+\pi,t=1.234],
\end{centred}
which is $(\sin\tfrac{7}{2})\times\pi$. With \verb`\q` between the
bracketed factors,
\begin{centred}
\verb`\eval{$ \sin(n+\tfrac12)\q(x-t) $}[n=3,x=t+\pi,t=1.234]` $\Longrightarrow$
\eval{$ \sin(n+\tfrac12)\q(x-t) $}[n=3,x=t+\pi,t=1.234],
\end{centred}
which is $\sin(\tfrac{7}{2}\pi)$. Similarly, without \verb`\q`,
\begin{centred}
\verb`\eval[p]{\[ \cos\frac{2\pi}{T}n(t+\tfrac12T) \]}[T=2,t=1,n=3]`
$\Longrightarrow$ \eval[p]{\[ \cos\frac{2\pi}{T}n(t+\tfrac12T) \]}[T=2,t=1,n=3]
\end{centred}
which is $(\cos\pi)\times3\times(1+\tfrac{1}{2}\times2)$. With \verb`\q`
used twice, once after the fraction and once before the left parenthesis,
\begin{verbatim}
    \eval[p]{\[ 
      \cos\frac{2\pi}{T}\q n\q(t+\tfrac12T)
        \]}[T=2,t=1,n=3]
\end{verbatim}
$\Longrightarrow$ \eval[p]{\[ \cos\frac{2\pi}{T}\q n\q(t+\tfrac12T) \]}[T=2,t=1,n=3]

\noindent which is $\cos(\pi\times3\times2)$.

It should be noted that for \emph{trigonometric }functions, because
of their use in Fourier series especially, there is another way of
handling arguments that contain parentheses (and fractions). This
is discussed in §\ref{subsec:parseTrigFns} below.

For the \verb`\Q` command which splits an argument we have, without
it,
\begin{centred}
\verb`\eval{$ 1/2e $}` $\Longrightarrow$ \eval{$ 1/2e $}, 
\end{centred}
which is the reciprocal of $2e$, whereas with the \verb`\Q` command
inserted before \verb`e`,
\begin{centred}
\verb`\eval{$ 1/2\Q e $}` $\Longrightarrow$ \eval{$ 1/2\Q e $},
\end{centred}
which is one half of $e$, although it is unlikely to be read that
way. If one half of $e$ is intended then parenthesize the $1/2$
or present as a \verb`\tfrac`.

\subsubsection{Mnemonic}

As mnemonic, best seen in sans serif for the Latin Modern fonts used
in this document, think of the letter \textsf{q} as a circle \emph{cleaving
to} a vertical descender; think of the letter \textsf{Q} as a circle
\emph{cleaved apart} by the diagonal stroke.

\subsection{Parsing groups}

The arguments of different groups of functions are handled in different
ways. The criterion used for deciding when an argument ends for one
group will not be that used for others. Table §\ref{subsec:defaultsIntifyingRounding}
lists the different groups that \texttt{numerica} takes account of.
At the top are functions or operations that have the smallest reach
when determining where their arguments end; at the bottom are operations
that have the greatest reach. The denominator of a slash fraction
is treated as a unary function and is assigned to group II. By default
trigonometric functions are treated the same as other unary functions
but there is a setting which enables the direct (rather than inverse)
trigonometric functions to accept a wider range of arguments, as occurs
in Fourier series. Hence they are separated into their own group.

\begin{table}
\centering

\caption{ Parsing groups}\label{tab:settingsParsing-groups}
{\ttfamily{}%
\begin{tabular}{ll}
\toprule 
{\small\textrm{group}} & {\small\textrm{function/operation}}\tabularnewline
\midrule
{\small\textrm{I}} & {\small\textrm{surd, logical Not}}\tabularnewline
{\small\textrm{II}} & {\small\textrm{unary functions (direct trig. functions default), /}}\tabularnewline
{\small\textrm{III}} & {\small\textrm{direct trig. functions with special setting}}\tabularnewline
{\small\textrm{IV}} & {\small\textrm{sums, products}}\tabularnewline
{\small\textrm{V}} & {\small\textrm{comparisons}}\tabularnewline
{\small\textrm{VI}} & {\small\textrm{logical And, logical Or}}\tabularnewline
\bottomrule
\end{tabular}}

\end{table}
 A formula is a sequence of tokens and brace groups. All parsing occurs
from the left, \LaTeX{} argument by \LaTeX{} argument, where \emph{argument}
means either a token (an N-type argument in \verb`expl3`-speak) or
a brace group (an n-type argument). To distinguish \LaTeX{} arguments
from mathematical arguments I shall when necessary refer to L-args
and M-args. A mathematical argument may end \emph{at} an L-arg, meaning
immediately before the L-arg, or end \emph{with} the L-arg, meaning
immediately after the L-arg. Ending or not will in general depend
on whether the argument is in \emph{first position} – the position
immediately following a function token like \verb`\sin` or \verb`\log`
– or in \emph{general position} – any later position (although for
trigonometric functions we will also need to consider \emph{second}
and even \emph{third }positions). 

For counting position, we need to allow for formatting elements and
multi-token numbers – in both decimal and scientific formats. Formatting
elements do not change the position count. This applies to things
like thin spaces or phantoms (and their arguments) or modifiers like
\verb`\left` or \verb`\biggl`. Multi-token numbers (in decimal or
scientific formats) are treated as single items; they advance the
position count by exactly one. \LaTeX{} functions – like \verb`\frac`
– which take \LaTeX{} arguments again advance the position count only
by one. Mathematically, the fraction is viewed as a single unit.

I shall refer to a token or a token and its \LaTeX{} arguments – like
\verb`\frac` and its arguments – as an \emph{item}. Similarly, a
(possibly multi-token) number is an item. Also it will help to distinguish
tokens within brackets where both brackets lie to the right of a function
from those that do not. The former I call \emph{clothed}; the latter
are \emph{naked}. Thus the plus sign in $(\sin x+y)$ is naked relative
to the sine (one bracket to the left of the function), but is clothed
in $\sin(x+y)$ (both brackets to the right of the function). 

\subsubsection{Parsing group I}

The only functions in this category are the surd and logical Not. 

Why distinguish the surd from other unary functions? Surely we all
agree that \verb`\sin2\pi`, displaying as $\sin2\pi$, vanishes?
The argument of the sine extends beyond the $2$ to include the $\pi$.
But \verb`\surd2\pi`, displaying as $\surd2\pi$, is understood to
be the product $\surd2\times\pi$. The argument of the surd ends with
the $2$. The surd binds more tightly to its argument than is true
of unary functions generally. 

For parsing group I
\begin{enumerate}
\item if a left bracket is in first position, the mathematical argument
ends with the matching right bracket; otherwise
\item the argument ends with the item in first position and any L- or M-args
required by that item.
\end{enumerate}
If the factorial sign \verb`!` \emph{preceded} its argument, it too
would belong to this parsing state, for it also binds tightly like
the surd. This means that an expression like $\surd4!$ is intrinsically
ambiguous. Is it the square root of $24$ or the factorial of $2$?
In \texttt{numerica} it produces the (perhaps rather odd) error
\begin{centred}
\verb`\eval{$ \surd 4! $}` $\Longrightarrow$ \eval{$ \surd4! $}
\end{centred}
The surd has seized the argument; there is nothing for the factorial
to operate on. The same error arises if the $4$ is parenthesized,
but parenthesizing like either \verb`(\surd 4)!` or \verb`\surd(4!)`
repairs the situation. Because other unary functions (like the sine
or logarithm) do not bind as tightly, this ambiguity does not arise
for them.

Exponents cause no problem because taking square roots and raising
to a power are commutative operations – the result is the same whichever
is performed first. 
\begin{centred}
\verb`\eval{$ \surd 3^4 $}` $\Longrightarrow$ \eval{$ \surd 3^4 $}.
\end{centred}

\subsubsection{Parsing group II}

In the default setup this category includes the trigonometric and
hyperbolic functions, their inverses, the various logarithms and the
exponential functions, the signum function \verb`\sgn`, and the denominators
of slash fractions \verb`/`. Note however that there is a setting
switch which enables trigonometric  functions to handle parentheses
in arguments more generally; see §\ref{subsec:parseTrigFns}. 
\begin{itemize}
\item In parsing group II we wish to accommodate usages like $\ln z^{n}=n\ln z$
(\emph{HMF} 4.1.11), or $\gd z=2\arctan e^{z}-\frac{1}{2}\pi$ (\emph{HMF}
4.3.117), defining the Gudermannian. The exponent is included in the
argument. Considering $\ln(1+1/n)^{n}$ exponents must also be part
of parenthesized arguments. 
\item An approximation to Stirling's formula for the factorial is often
written $\ln N!\approx N\ln N-N$ (widely used in texts on statistical
mechanics). Hence the factorial sign should also be considered part
of the argument.
\item $\ln xy=\ln x+\ln y$ means the argument must reach over a product
of variables. Identities like $\sin2z=2\sin z\cos z$ mean the argument
also reaches over numbers, and expressions like $\sin\tfrac{1}{2}\pi x$
(\emph{HMF} 4.3.104) mean that it further reaches over \verb`\tfrac`-s
and constants.
\item Essentially \emph{anything }can be in first position, and without
parentheses; e.g.
\begin{itemize}
\item unary functions: $\ln\ln z$ (\emph{HMF} 4.1.52), $\ln\tan\dfrac{z}{2}$
(\emph{HMF} 4.3.116),
\item fractions: $\ln\dfrac{z_{1}}{z_{2}}$ (\emph{HMF} 4.1.9), $\arcsin\dfrac{(2ax+b)}{(b^{2}-4ac)^{1/2}}$
(\emph{HMF} 3.3.36), $\ln\dfrac{\tan z}{z}$ (\emph{HMF} 4.3.73),
\item absolute values: $\ln\abs*{\dfrac{a+x}{a-x}}$ (\emph{HMF} 3.3.25),
\item square roots: $\arctan\sqrt{\dfrac{\nu_{1}}{\nu_{2}}F}$ (\emph{HMF
}26.6.8)
\end{itemize}
\end{itemize}
With these examples in mind, for parsing group II
\begin{enumerate}
\item if a left bracket is in first position, the mathematical argument
ends with the matching right bracket and any attached exponent, or
factorial or double factorial sign; otherwise
\item the mathematical argument includes the item in first position and
any L- or M-args required by that item; 
\begin{enumerate}
\item if the item in first position is a number, variable, constant or \verb`\tfrac` 
\begin{enumerate}
\item the argument appends the next item if it is a number, variable, constant
or \verb`\tfrac`, and so on recursively; or
\item the argument appends the next item if it is an exponent, or facorial
or double factorial sign, and ends there; otherwise
\item the argument ends.
\end{enumerate}
\item if the item in first position is not a number, variable, constant
or \verb`\tfrac`
\begin{enumerate}
\item the argument appends the next item if it is an exponent, or factorial
or double factorial sign, and ends there; otherwise
\item the argument ends.
\end{enumerate}
\end{enumerate}
\end{enumerate}
An argument may extend over (see 2(a)i) numbers, constants, variables
and \verb`\tfrac`-s, as instanced with $\sin2\tfrac{p}{q}\pi x$
which exhibits all elements.

Illustrating 1, the exponent is included in the argument but not the
following variable:
\begin{centred}
\verb`\eval{$ \log_{10}(1+2+3+4)^3n $}[n=5]` $\Longrightarrow$ \eval{$ \log_{10}(1+2+3+4)^3n $}[n=5].
\end{centred}
For the sake of the reader, and as one naturally does in any case
to avoid ambiguity, the formula should be written with the variable
$n$ preceding the logarithm: $n\log_{10}(1+2+3+4)^{3}$. The way
the example is written suggests that the writer wished the $n$ to
be considered part of the argument of the logarithm. If that is the
case, inserting a \verb`\q` command before \verb`n` would achieve
this, but that would still be confusing for the reader of the pdf.
Inserting parentheses is the only sensible thing to do. 

Illustrating 2(a)ii, again the exponent is included in the argument
but not the following variable:
\begin{centred}
\verb`\eval{$ \log_{10}m^3n $}[m=10,n=5]` $\Longrightarrow$ \eval{$ \log_{10}m^3n $}[m=10,n=5].
\end{centred}
Again, for the sake of the reader and as one naturally does to avoid
ambiguity, the variable $n$ should precede the logarithm. If in fact
the intention was for the $n$ to be included in the argument of the
logarithm, then again the \verb`\q` command could be used or, better
in this case, the $n$ could be shifted to precede the $m$, which
illustrates 2(a)i:
\begin{centred}
\verb`\eval{$ \log_{10}nm^3 $}[m=10,n=5]` $\Longrightarrow$ \eval{$ \log_{10}nm^3 $}[m=10,n=5],
\end{centred}
the logarithm of $5000$, or better still, $m^{3}n$ could (and should)
be parenthesized for the sake of the reader.

Why the difference between $nm^{3}$ where $n$ and $m^{3}$ are included
in the argument, and $m^{3}n$ where $n$ is not? Any criterion is
going to miss some instances where a different outcome might be desirable.
Where an argument ends is affected by visual appearance in the pdf.
It is simple and easy to remember if it is understood that anything
that breaks the `visual flow' of juxtaposed numbers, variables,
constants and \verb`\tfrac`-s ends the argument. An exponent does
just that. If you feel there is ambiguity, parenthesize to clarify.

Illustrating 2(b)ii, the argument stops with the \verb`\dfrac` and
its arguments and does not extend to the following constant:
\begin{centred}
\verb`\eval{$ \sin\dfrac12\pi $}` $\Longrightarrow$ \eval{$ \sin\dfrac12\pi $}.
\end{centred}
Obviously, someone writing an expression like this intends the $\pi$
to be part of the argument. In that case, a \verb`\tfrac` should
be used since the \verb`\dfrac` breaks the `visual flow' of the
argument. 
\begin{description}
\item [{Fractions}]~

But why not a plain \verb`\frac`? After all, for an inline expression
it displays in the same way as a \verb`\tfrac`. I considered making
the argument-behaviour of \verb`\frac` the same as \verb`\tfrac`
for text-style contexts, and the same as \verb`\dfrac` for display-style
contexts, but that would have meant the same expression evaluating
to different results depending on whether it lay between \verb`$ $`
or \verb`\[ \]` delimiters, which ruled it out. Because \verb`\frac`
sometimes displays as \verb`\dfrac`, it is treated like \verb`\dfrac`
(but see §\ref{subsec:parseTrigFns}, specifically \texttt{()=2}). 
\item [{Slash~fractions}]~

It is easy to write ambiguous expressions using the slash $/$ to
indicate fractions or division. How should $\pi/2n$ be interpreted?
With from-the-left evaluation and calculator precedence rules which
give equal precedence to {*} (multiplication) and / (division), this
would be interpreted as $(\pi/2)\times n$, but most people will instinctively
interpret it as $\pi/(2n)$. By placing \verb`/` in parsing group
II, this is what \texttt{numerica} does. It treats the right-hand
argument of the slash \emph{as if it were the argument of a named
function}. This means that $1/2\sin(\pi/6)$ is parsed as $(1/2)\sin(\pi/6)$
rather than as $1/(2\sin(\pi/6))$. It also means that $1/2\exp(1)$
and $1/2e$ give different results, which (in the author's view) is
acceptable since they display differently and are not instinctively
read in the same way. 
\end{description}

\subsubsection{Parsing group III}

\label{subsec:parseTrigFns}By default trigonometric functions are
set to parsing group II. This accommodates many instances of how arguments
are used with these functions, but Fourier series in particular require
more. For them we need to take account of how \emph{parentheses} are
used in arguments. I find $\tan\tfrac{1}{2}(A+B)$ (\emph{HMF }4.3.148),
$\sec\pi(\tfrac{1}{4}+\tfrac{1}{2}az)$ (\emph{HMF }19.3.3), $\cos(2m+p)z$
(\emph{HMF }20.2.3), $\sin(2n+1)v$ (\emph{HMF }16.38.1). Looking
through various texts discussing Fourier series it is easy to find
examples like
\[
\cos\frac{2\pi}{T}nt,\quad\cos\frac{2\pi}{T}n(t+\tfrac{1}{2}T),
\]
and
\[
\cos(N+\tfrac{1}{2})\frac{2\pi\tau}{T},\quad\sin2\pi\left(\frac{x}{\lambda}-\frac{t}{T}\right).
\]
In the last of these \verb`\left` and \verb`\right` have been used
to enlarge the parentheses. 

All these usages can be accommodated by adjusting a setting in the
settings option (§\ref{sec:settingsOption}) of the \verb`\eval`
command:
\begin{lyxcode}
()~=~<integer>
\end{lyxcode}
where \texttt{<integer>} is one of \texttt{0, 1, 2}. 

\noindent\begin{minipage}[t]{1\columnwidth}%
\begin{shaded}%
I remain unsure whether to persist with the \texttt{()} setting. A
principal aim of \texttt{numerica} is to avoid having to modify a
formula to bring it into a form that can be evaluated. The \texttt{()}
setting arose in that context, but it complicates the code and may
well confuse the user. Inserting cleave functions, \texttt{\textbackslash q}
or \texttt{\textbackslash Q}, to achieve the same effects does mean
modifying formulas, but is straightforward and easier to understand.
(And \texttt{\textbackslash q} and \texttt{\textbackslash Q} have
no effect on the visual appearance of formulas.)\end{shaded}%
\end{minipage}

For convenience of statement in what follows call parentheses, square
brackets or braces \emph{brackets}. If preceded by a \verb`\left`
or \verb`\right` or \verb`\biggl` or \verb`\biggr` etc. modifier,
call them \emph{Brackets}, with an uppercase `B'. Modifiers do not
contribute to the position count, so that a left Bracket in first
position means the modifier and left bracket are both considered to
be in first position. When it is immaterial whether it is a bracket
or a Bracket I write b/Bracket. The rules that follow do not prescribe
what mathematicians \emph{ought} to do but are intended to be descriptive
of certain patterns of mathematical practice as discerned in \emph{HMF}
and a number of texts (about half a dozen) on Fourier series.
\begin{description}
\item [{\texttt{()=0}}] is the \emph{default} setting, parsing group II
behaviour; b/Brackets are included in the argument only if
\begin{itemize}
\item the left b/Bracket is in first position;
\begin{itemize}
\item if the first item beyond the matching right b/Bracket is an exponent,
or factorial or double factorial sign, it is appended to the argument,
which ends there, otherwise
\item the argument ends with the right b/Bracket.
\end{itemize}
\end{itemize}
\end{description}
The default setting allows things like $\sin\tfrac{1}{2}a$, $\cos2\pi nt$
and $\tan(A+B)$. It does \emph{not} encompass examples like $\tan\tfrac{1}{2}(A+B)$
or $\cos2(n+\tfrac{1}{2})\pi$. For that we need the setting \verb`()=1`:
\begin{description}
\item [{\texttt{()=1}}] includes a b/Bracketed expression in the argument,
provided 
\begin{itemize}
\item the left Bracket is in first position;
\begin{itemize}
\item if the first item beyond the matching right Bracket is an exponent,
or factorial or double factorial sign, it is appended to the argument,
which ends there, otherwise
\item the argument ends with the right Bracket.
\end{itemize}
\item or the item in first position is a number, variable, constant or \verb`\tfrac`
and the left bracket is in second position;
\begin{itemize}
\item if the first item beyond the matching right bracket is an exponent,
or factorial or double factorial sign, it is appended to the argument,
which ends there, or
\item if the first item beyond the matching right bracket is a number, variable,
constant, or \verb`\tfrac` it is appended to the argument, and so
on recursively, until
\begin{itemize}
\item an exponent, or factorial or double factorial sign is met, which is
appended to the argument which ends there, or
\item an item is met which is \emph{not} an exponent, or factorial or double
factorial sign, or a number, variable, constant or \verb`\tfrac`,
at which point the argument ends, or
\item the end of the formula is reached.
\end{itemize}
\end{itemize}
\end{itemize}
\end{description}
With the \verb`()=1` setting, the arguments of $\tan\tfrac{1}{2}(A+B)$,
$\sec\pi(\tfrac{1}{4}+\tfrac{1}{2}az)$, $\cos(2m+p)z$, $\sin(2n+1)v$
are all accommodated, as is $\sin\tfrac{1}{2}(m+n)\pi$ with items
on both sides of the parentheses. But, note, there must be at most
\emph{one} item before the left parenthesis:
\begin{centred}
\verb`\eval[()=1]{$ \sin\tfrac16(m+n)\pi $}[m=1,n=2]`. $\Longrightarrow$
\eval[()=1]{$ \sin\tfrac16(m+n)\pi $}[m=1,n=2],
\end{centred}
whereas, with two items before the left parenthesis,
\begin{centred}
\verb`\eval[()=1]{$ \sin2\tfrac1{12}(m+n)\pi $}[m=1,n=2]`. $\Longrightarrow$

\eval[()=1]{$ \sin2\tfrac1{12}(m+n)\pi $}[m=1,n=2].
\end{centred}
Whatever the \verb`()` setting, \texttt{numerica} does not check
what is included between the parentheses (or brackets generally) –
it could be anything. However inserting \verb`\left`, \verb`\right`
or other modifiers before the parentheses restricts the argument of
the sine in this example, despite the \verb`()=1` setting, to the
\verb`\tfrac`:
\begin{centred}
\verb`\eval[()=1]{$ \sin\tfrac16\left(m+n\right)\pi $}[m=1,n=2]`
$\Longrightarrow$ \eval[()=1]{$ \sin\tfrac16\left(m+n\right)\pi $}[m=1,n=2].
\end{centred}
Although \verb`()=1` serves well for the kinds of expressions and
identities involved in trigonometry, perusal of any text on Fourier
series will show it does not cover the kinds of expressions met there.
For that we need
\begin{description}
\item [{\texttt{()=2}}] includes a b/Bracketed expression in the argument
provided
\begin{itemize}
\item the left b/Bracket is in first position, or the item in first position
is a number, variable, constant, \verb`\dfrac`, \verb`\frac` or
\verb`\tfrac` and the left b/Bracket is in second position, or the
items in first and second positions are numbers, variables, constants,
\verb`\dfrac`-s, \verb`\frac`-s or \verb`\tfrac`-s and the left
b/Bracket is in third position;
\begin{itemize}
\item if the first item beyond the matching right b/Bracket is an exponent,
or factorial or double factorial sign, it is appended to the argument,
which ends there, or
\item if the first item beyond the matching right b/Bracket is a number,
variable, constant, \verb`\dfrac`, \verb`\frac` or \verb`\tfrac`
it is appended to the argument, and so on recursively, until
\begin{itemize}
\item an exponent, or factorial or double factorial sign is met, which is
appended to the argument which ends there, or
\item an item is met which is \emph{not} an exponent, or factorial or double
factorial sign, or a number, variable, constant, \verb`\dfrac`, \verb`\frac`
or \verb`\tfrac`, at which point the argument ends, or
\item the end of the formula is reached.
\end{itemize}
\end{itemize}
\end{itemize}
\end{description}
{\ttfamily\verb`()=2`} draws no distinction between brackets
and Brackets. It allows all \verb`()=1` possibilities but also \emph{two
}items (of a suitable kind) before a left b/Bracket; it also treats
\verb`\dfrac`-s and \verb`\frac`-s like \verb`\tfrac`-s for determining
the scope of arguments. 

The following examples are taken from different texts on Fourier series.
The first shows a \verb`\frac` being included in the argument, the
second shows \emph{two} items – including a \verb`\frac` – preceding
the left parenthesis, the third shows a \verb`\frac` to the right
of the parentheses, and the fourth shows parentheses using \verb`\left`-\verb`\right`
modifiers with two items preceding them:
\[
\cos\frac{2\pi}{T}nt,\quad\cos\frac{2\pi}{T}n(t+\tfrac{1}{2}T),\quad\text{\ensuremath{\sin(N+\tfrac{1}{2})\frac{2\pi\tau}{T}}\ensuremath{\quad}and}\quad\sin2\pi\left(\frac{x}{\lambda}-\frac{t}{T}\right).
\]
All these usages are accommodated by the \verb`()=2` setting. For
instance
\begin{verbatim}
    \eval[p,()=2]
      { 
        \[ \sin(N+\tfrac12)\frac{2\pi\tau}T \]
      }[N=1,\tau=2,T=3]
\end{verbatim}
$\Longrightarrow$     \eval[p,()=2]
      { 
        \[ \sin(N+\tfrac12)\frac{2\pi\tau}T \]
      }[N=1,\tau=2,T=3]which is the sine of $2\pi=(\tfrac{3}{2})\times(\tfrac{4}{3}\pi)$
where a \verb`\frac` trailing the parentheses has been included in
the argument, and \emph{not }$(\sin\tfrac{3}{2})(\tfrac{4}{3}\pi)$.
Or consider
\begin{verbatim}
    \eval[p,()=2]
      {\[ 
         \sin2\pi\left(\frac{x}{\lambda}
            -\frac{t}{T}\right) 
      \]}[x=1,\lambda=2,t=3,T=4]
\end{verbatim}
$\Longrightarrow$     \eval[p,()=2]
      {\[ 
         \sin2\pi\left(\frac{x}{\lambda}
            -\frac{t}{T}\right) 
      \]}[x=1,\lambda=2,t=3,T=4]which is the sine of $-\tfrac{1}{2}\pi=2\pi\times(-\tfrac{1}{4})$
where there are two items before the parentheses and \verb`\left`
and \verb`\right` modifiers, and \emph{not} $\sin2\pi$ times the
parenthesised expression.

However a usage like $\sin(n+\tfrac{1}{2})(x-t)$, noted in two different
texts, is not available without explicit use of the \verb`\q` command
between the parenthesized groups. 

\subsubsection{Parsing group IV}

The only members of this group are \verb`\sum` and \verb`\prod`.

For parsing group IV
\begin{enumerate}
\item the argument ends 
\begin{enumerate}
\item at the first naked plus or minus sign encountered, or 
\item at the first comparison sign or comparison command encountered, or
\item at the first logical And or logical Or sign encountered, or
\item at the end of the formula.
\end{enumerate}
\end{enumerate}
In practice this means mainly (a) and (d), and seems to be the instinctive
practice. \emph{HMF} has multiple examples in multiple chapters of
the argument to a sum ending at a naked plus sign: 7.3.12 \& 7.3.14,
9.1.11 \& 9.1.77, 9.6.35 \& 9.6.43, 11.1.9, \ldots{} (at that point
I stopped looking). They were all of the form
\[
\sum\text{argument}+\ldots
\]
 A minus sign serving the same purpose was harder to find but \emph{HMF}
10.4.65 \& 10.4.67 are two instances. I considered whether a \verb`\times`
or slash fraction sign \verb`/` might end the argument of a sum,
but surely we need to allow things like $\sum1/n^{2}$ which rules
out the slash and \emph{HMF} 9.9.11 provides two of a number of instances
in \emph{HMF} of sum arguments continuing past explicit \verb`\times`
signs (at line breaks when a summand spills onto a second line).

Because they are evaluated using the same code as sums I (unthinkingly)
placed products with sums but doubts later intruded. In \emph{HMF}
products occur only occasionally and are almost all of the form
\[
\prod\left(\text{argument}\right)
\]
where the argument is bracketed (often with \verb`\left \right` modifiers)
and the multiplicand ends with the right bracket. At least twice (\emph{HMF
}6.1.25 and 24.2.2.1) an exponent ($-1$) is attached to the right
bracket and the argument ends there. Looking further afield, a text
on number theory has examples where the argument of the product extends
to \emph{three} parenthesised factors, $\prod\left(\text{arg}1\right)\left(\text{arg2}\right)\left(\text{arg3}\right)$
and a number of others where it extends to two. A text on theory of
functions has 
\[
\prod_{n=1}^{\infty}\left(1+\frac{z}{n}\right)e^{z/n}
\]
although \emph{HMF}, for the same expression, encloses the two factors
within (large) square brackets, as if some ambiguity existed as to
how far the reach of the \verb`\prod` extended. 

\emph{Tentatively} I retain products here in the same group as sums.

\subsubsection{Parsing group V}

Comparison symbols compose this group: \texttt{=}, \texttt{<}, \texttt{>},
\verb`\ne`, \verb`\le`, \verb`\ge`, \verb`\leq`, \verb`\geq`,
and the various comparison commands from the \texttt{amssymb} package
listed in §\ref{subsec:evalAmssymb-comparisons}. Because of the way
\texttt{numerica} handles comparisons, it is the argument on the right-hand
side of the relation that needs determining.

For parsing group V
\begin{enumerate}
\item the argument ends at
\begin{enumerate}
\item the first logical And or logical Or encountered, or
\item the first comparison sign or command encountered, or
\item the end of the formula.
\end{enumerate}
\end{enumerate}

\subsubsection{Parsing group VI}

Logical And and logical Or are the sole members of this group. It
is the right-hand side of the And or Or command that needs determining.

For parsing group VI
\begin{enumerate}
\item the argument ends at
\begin{enumerate}
\item the first logical And or logical Or encountered, or
\item the end of the formula.
\end{enumerate}
\end{enumerate}

\subsubsection{Disclaimer}

The parsing rules of the different groups are not normative; they
are not statements of how mathematical formulas should be written.
Rather they are attempts to discern regularities in how mathematicians
often do write formulas. It is \emph{how things look in the pdf},
not \LaTeX , that is the guide. You are always free to parenthesize
as you see fit and to insert cleave commands (\verb`\q` or \verb`\Q`)
to force outcomes. 

(But note that parenthesizing has its limits. For sums, writing 
\[
\sum\left(\mathtt{<stuff>}\right)\mathtt{<more\ stuff>}
\]
 does not necessarily end the summand at the right parenthesis: it
ends at the first naked $+$ or $-$ sign, or \verb`\Q` command,
encountered.) 

The rule should always be to write expressions that are clear to the
reader of the pdf. An expression that is ambiguous to the reader,
even if it fits within the parsing rules, is to be deprecated. The
\emph{intent} is that \verb`\eval` can parse unambiguous expressions
correctly. 

\chapter{Supplementary commands}

\label{chap:Supplementary-commands}This chapter introduces four commands,
\verb`\nmcInfo` (which we have already met), \verb`\nmcMacros`,
\verb`\nmcConstants` and \verb`\nmcReuse`, supplementary to the
principal command \verb`\nmcEvaluate`. They use the same machinery
as \verb`\nmcEvaluate` and so have the same syntax. If all arguments
are present it is
\begin{centred}
\noindent \verb`\nmc<cmd>*[settings]{main arg}[vv-list][rounding]`
\end{centred}
where \verb`<cmd>` is one of \verb`Info`, \verb`Macros`, \verb`Constants`
and \verb`Reuse`. All four commands have short-name forms: \verb`\info`,
\verb`\macros`, \verb`\constants`, \verb`\reuse`.

Generally the final two optional arguments will not be used. The user
should be aware of this if following a command with a square bracketed
expression – the expression will be absorbed without trace unless
it is preceded by, for example, an empty brace pair.

Because the commands share the machinery of \verb`\nmcEvaluate`,
the settings discussed earlier (Chapter~\ref{chap:Settings}) for
the \verb`\eval` command are also available for these commands, although
they will, in the main, be irrelevant. The `debug' code has been
used by the \verb`view` setting of some of these supplementary commands
to produce its effects. 

The starred form of command is available in all four cases and in
all cases produces a pure number. If both star and \verb`view` are
used at the same time, the \verb`view` setting prevails over starring.

\section{Feedback on \textquoteleft infinite\textquoteright{} processes:\texttt{ \textbackslash nmcInfo}}

\label{sec:supplInfo}Used after the evaluation of an `infinite'
process, the \verb`\nmcInfo` command, or its short-name form \verb`\info`
will tell you how many terms or factors or other operations\footnote{It also applies to the commands \texttt{\textbackslash nmcIterate
}and \texttt{\textbackslash nmcSolve} from the \texttt{numerica-plus
}package and to derivatives and integrals from the \texttt{numerica-calculus}
package.} were needed to arrive at the result.The main argument contains an
identifier for the `infinite' process: 
\begin{lyxcode}
\textbackslash nmcInfo\{<arg>\}~
\end{lyxcode}
(or \verb`\info{<arg>}`) where, at this stage, \verb`<arg>` is either
\verb`sum` or \verb`prod`. The display, as we have seen in earlier
examples, is a number followed by a space then a descriptor. For \verb`sum`
and \verb`prod` the descriptors are \verb`terms` and \verb`factors`.
Starring \verb`\nmcInfo` – \verb`\nmcInfo*{arg}` or \verb`\info*{arg}`
– suppresses the descriptor and leaves only the number. This allows
the starred form to be nested in an \verb`\eval` command, which might
sometimes be convenient. 

As an example, let's test `the hard way' a standard identity, $\cosh^{2}x-\sinh^{2}x=1$.
We know that $\cosh x=\sum_{n=0}^{\infty}\frac{x^{2n}}{(2n)!}$ and
$\sinh x=x\prod_{k=1}^{\infty}\left(1+\frac{x^{2}}{k^{2}\pi^{2}}\right)$.
The difference of their squares should be $1$:
\begin{verbatim}
    \eval{\[  
      \left[\sum_{n=0}^{\infty}
        \frac{x^{2n}}{(2n)!}
      \right]^2-
        \left[x\prod_{k=1}^{\infty}
          \left(1+\frac{x^{2}}{k^{2}\pi^{2}}\right)
        \right]^2
      \]}[x=1][3] \info{sum},\quad \info{prod}
\end{verbatim}
$\Longrightarrow$     \eval{\[  
      \left[\sum_{n=0}^{\infty}
        \frac{x^{2n}}{(2n)!}\right]^2-
      \left[x\prod_{k=1}^{\infty}
        \left(1+\frac{x^{2}}{k^{2}\pi^{2}}\right)
    \right]^2
        \]}[x=1][3] \info{sum},\quad\info{prod}.

Nearly right. Obviously the product converges only slowly which is
where the error comes from (see the discussion in §\ref{sec:settingsInfiniteSumsProds},
where we needed the extra rounding setting \texttt{P+=3} and $350$
factors to get a correct 3-figure value). The point of the example
is to show the information command being used for both sum and product
in the one evaluation. One does not exclude the other. 

\subsection{Suppressing the descriptor: \texttt{\textbackslash nmcInfo{*}}}

The starred form of the \verb`\info` command suppresses the descriptor
(`terms', `factors') and gives a purely numerical result: 
\begin{verbatim}
    \eval{$
      \sum_{k=0}^{\infty}\binom \alpha k x^k
        $}[x=1/2,\alpha=3], 
    requiring $ \info*{sum}-1 $ additions. 
\end{verbatim}
$\Longrightarrow$     \eval{$
      \sum_{k=0}^{\infty}\binom \alpha k x^k
         $}[x=1/2,\alpha=3], 
requiring $ \info*{sum}-1 $ additions. (Four terms added, therefore $3$ additions.)

\subsection{Errors}

Should the \emph{wrong} argument be used in the \verb`\nmcInfo` command,
no harm is done:
\begin{verbatim}
    \eval{$
      \sum_{k=0}^{\infty}\binom \alpha k x^k
         $}[x=1/2,\alpha=3], \ \info{prod} 
\end{verbatim}
$\Longrightarrow$  \eval{$ \sum_{k=0}^{\infty}\binom \alpha k x^k $}[x=1/2,\alpha=3],\ \info{prod}.\medskip{}

$119$ \emph{factors}? The information command is remembering a previous
result, the last time \verb`prod` was used as its argument. Changing
the argument from \verb`prod` to \verb`sum` reveals the correct
number of \emph{terms}.

Should a non-existent argument be used, an error message is generated:
\begin{verbatim}
    \eval{$ 
      \sum_{k=0}^{\infty}\binom \alpha k x^k 
         $}[x=1/2,\alpha=3], \\ \info{Fred} 
\end{verbatim}
$\Longrightarrow$  \eval{$ \sum_{k=0}^{\infty}\binom \alpha k x^k $}[x=1/2,\alpha=3],\\  \info{Fred}

\subsection{\texttt{view} setting}

As noted at the start of this chapter, \verb`\nmcInfo` uses the `machinery'
of \verb`\nmcEvaluate`. Most of the settings available for \verb`\eval`
are also available for \verb`\info` but of these only one seems relevant:
the \verb`dbg` setting. However, rather than use the obscure \verb`dbg=<integer>`
(which is possible), it suffices to enter \verb`view` in this argument:
\begin{centred}
\verb`\info[view]{}` $\Longrightarrow$ \info[view] {}
\end{centred}
The result is a display of all the current values of all the `infinite'
processes available. All such values are initialized to $0$. (Further
processes \verb`iter` and \verb`solve` become relevant if the \verb`numerica-plus`
package is used; \verb`deriv` and \verb`integ` become relevant if
the \verb`numerica-calculus` package, currently under development,
is used.)

\section{User-defined macros: \texttt{\textbackslash nmcMacros}}

\label{sec:supplMacros}The \verb`\nmcMacros` and \verb`\nmcConstants`
commands were prompted by a question on \TeX{} Stack Exchange.\cprotect\footnote{A question from Giacomo Petrillo on \TeX{} Stack Exchange and a response
by `egreg' prompted the introduction of the \cprotect\texttt{\textbackslash nmcMacros}
and \cprotect\texttt{\textbackslash nmcConstants} commands. See \url{https://tex.stackexchange.com/questions/602993/use-macros-in-numerica-vv-list/602998#602998}} Some time later the maintainer of the \verb`mandi` package\footnote{The maintainer is Joe Heafner, who explains that `mandi' is an abbreviation
of `matter and interactions' after a physics textbook of that name
(by different authors). Among other things, the package defines a
long list of macros, each containing the value of a physical constant.} approached me with a similar problem. Suppose one has defined a macro
to contain a value, say 
\begin{itemize}
\item \verb`\def\myvalue{0.35}`, or
\item \verb`\newcommand\myvalue{0.35}`, or
\item \verb`\NewDocumentCommand\myvalue{}{0.35}`, if you're using \verb`xparse`.
\end{itemize}
(If you're using the document processor \LyX{} then there is good reason
to prefer \verb`\gdef` to define your macro, \verb`\gdef\myvalue{0.35}`;
see Chapter~\ref{chap:LyX}). After one of these commands, \verb`\myvalue`
is now known to \LaTeX , but it is not known to \verb`numerica`.
The quantities \verb`numerica` \emph{does }know about are variables
in the vv-list of an \verb`\eval` command, and those \LaTeX{} (and
\verb`amsmath` and \verb`mathtools`) commands used for writing mathematical
expressions. These quantities are stored in \verb`numerica` in structures
called property lists. Since \verb`\myvalue` is not recorded in these
lists yet, putting \verb`x=\myvalue` in the formula or vv-list of
an \verb`\eval` command will produce an `Unknown token' error message:
\begin{verbatim}
    \NewDocumentCommand \myvalue {} {0.35}
    \eval{ \myvalue }
\end{verbatim}
$\Longrightarrow$     \NewDocumentCommand \myvalue {} {0.35}
    \eval{ \myvalue }

With version 2 of \verb`numerica` a command is now available, \verb`\nmcMacros`,
to register macros and their values with the property lists used internally
by \verb`numerica`. (This command was not available in version 1.)
The macro must have been defined earlier in the document or in a supporting
package. 

The basic usage is simple. If you have a list of macros you wish to
make available to \verb`\nmcEvaluate`, enter them in a comma list
in the mandatory argument of \verb`\nmcMacros`:
\begin{lyxcode}
\textbackslash nmcMacros\{~\textbackslash macro1,~\textbackslash macro2,~\ldots ~\}
\end{lyxcode}
There is an equivalent short-name form of the command, \verb`\macros`. 

Multiple \verb`\nmcMacros` commands can be used in a document. If
the command is placed in the preamble (\emph{after} the definition
of the macros)\emph{ }then the user-defined macros and their values
are available throughout the document, otherwise they are available
from the position of the \verb`\macros` statement. However, macros
do not need to be defined in your current document provided they are
defined and accessible from elsewhere – for example from a loaded
\LaTeX{} package. But always an \verb`\nmcMacros` command is required
to `register' them with \verb`numerica` for use in an \verb`\eval`
command. 

\subsection{What can be stored in a macro?}

Generally a user-defined macro will store a number. This macro might
well be defined in an external package – for example the \verb`mandi`
package defines a large number of macros containing the values of
physical constants, some fundamental like the speed of light, others
contingent like the earth–moon distance. If the \verb`mandi` package
is loaded then writing, for instance, 
\begin{verbatim}
    \macros{ \electronmassprecisevalue,
             \protonmassprecisevalue }
\end{verbatim}
will make these two macros available for use in \verb`numerica`.
One could then write in the vv-list of an \verb`\eval` command
\begin{verbatim}
m_e=\electronmassprecisevalue,m_p=\protonmassprecisevalue
\end{verbatim}
which would allow (among other things) calculation of the mass ratio
$m_{p}/m_{e}$ of proton to electron. (The length of name of some
of the macros in the \verb`mandi` package has a pedagogical purpose,
but makes them unwieldy for direct use in mathematical expressions.)

\subsubsection{Macros containing formulas }

Numbers are not the only quantities that can be stored in a macro
for use in \verb`numerica`. In fact any mathematical expression that
can be \verb`\eval`-uated can be stored in a macro:
\begin{verbatim}
    \NewDocumentCommand \mysumC {}
        { \sum_{n=1}^{100}1/n - \ln 100 }
    \macros{ \mysumC }
    \eval{$ \mysumC $}[4]
\end{verbatim}
$\Longrightarrow$     \NewDocumentCommand \mysumC {}
        { \sum_{n=1}^{100}1/n - \ln 100 }
    \macros{ \mysumC }
    \eval{$ \mysumC $}[4], \medskip{}

\noindent (to be compared with Euler's constant \eval{$ \gamma $}[4]
– obviously many more terms are needed). The \verb`\eval` command
wraps around math delimiters in the example. Hence the result is presented
in the form \emph{formula=result}. In that presentation, note how
\verb`\mysumC` displays as the formula it contains.

\paragraph{The essential space: }

But the critical thing to notice in the example is \emph{the space
preceding }\verb`\sum`\emph{ in the definition of }\verb`\mysumC`.
When a formula starts with an expandable token, \emph{this space is
essential}. For macros to register successfully with \verb`numerica`,
the first character in their definition must be \emph{un}expandable.
Thus a digit is fine: storing a number in a macro is straightforward
and you don't need to fuss about such niceties. But a control sequence
like \verb`\sum` does expand (to $\sum$ ). If it is the initial
token of the formula, then it will cause a possibly obscure error
– see §\ref{subsec:supplMacrosErrors} – unless preceded by an unexpandable
token. Hence the space before \verb`\sum` in the \verb`\NewDocumentCommand`
statement. (On the other hand the spacing in the \verb`\macros` statement
is purely aesthetic.)

When using macros from another package, this is a matter to be aware
of. If the macros contain only numbers, there should be no problem,
but if they contain more complicated expressions, the absence of an
initial space could make them unusable in \verb`numerica`.

\subsubsection{Vv-list}

In the example it would be nice to be able to vary the number of terms
summed. This is easily done by using a vv-list in the \verb`\macros`
statement:
\begin{verbatim}
    \NewDocumentCommand \mysumN {}
        { \sum_{n=1}^{N}1/n - \ln N }
    \macros{ \mysumN }[N=150]
    \eval{$ \mysumN $}
\end{verbatim}
$\Longrightarrow$     \NewDocumentCommand \mysumN {}
        { \sum_{n=1}^{N}1/n - \ln N }
    \macros{ \mysumN }[N=150]
    \eval{$ \mysumN $}.\medskip{}

\noindent \verb`numerica` needs a definite value to store; it does
not store the formula as such. To give \verb`\mysumN` a definite
value, give the variable \verb`N` a value. This is done in the vv-list
added to the \verb`\macros` statement: \verb`N=150`. In this way
a definite value is stored in \verb`numerica` against the macro \verb`\mysumN`.
The definition of the macro is unaffected. If a new value is given
to \verb`N` in the \verb`\macros` statement (which is the point
of using a variable), the old value is overwritten and the new value
is used in subsequent calculations.

\subsection{Seeing what macros are available}

Perhaps your document has a number of \verb`\nmcMacros` statements
scattered through it and you want to remind yourself of what exactly
has been stored. \verb`\nmcMacros` has the \verb`view` setting for
this purpose. Writing
\begin{centred}
\verb`\macros[view]{}` $\Longrightarrow$ \macros[view]{}
\end{centred}
produces a list of all macros registered with \verb`numerica` and
their values, as you can see. 

If the braced argument is not empty, the display is slightly modified:
\begin{verbatim}
    \def\mydef{ \sin(m\pi/n) }
    \newcommand\mynewcmd{ \cos(m\pi/n) }
    \macros[view]{ \mydef,\mynewcmd }[m=3,n=18]
\end{verbatim}
$\Longrightarrow$     \def\mydef{ \sin(m\pi/n) }
    \newcommand\mynewcmd{ \cos(m\pi/n) }
    \macros[view]{ \mydef,\mynewcmd }[m=3,n=18]

\noindent \verb`\mydef` and \verb`\mynewcmd` have been added to
those available for use in \verb`numerica`. 

\subsubsection{Freeing macros from storage}

Rather than cluttering \verb`numerica`'s property lists with no-longer-needed
macros, it is possible to remove them from there with the \verb`free`
setting. This has no effect on the \LaTeX{} definition of the macro.
It merely `de-registers' the macro with \verb`numerica`.
\begin{centred}
\verb`\macros[free,view]{ \mysumC }` $\Longrightarrow$ \macros[free,view]{ \mysumC }
\end{centred}
If you want to free \emph{all} macros registered with \verb`numerica`
use an empty main argument with the \verb`free` setting. For an example,
see just below.

\subsubsection{Counting how many macros are available}

You can count how many macros are currently registered with \verb`numerica`
by starring the \verb`\nmcMacros`~command:
\begin{centred}
\verb`\macros*{}` $\Longrightarrow$ \macros*{}.
\end{centred}
If the braced argument is not empty, the list of macros it contains
will be added to those registered with \verb`numerica` and included
in the overall count.

Note that the \verb`view` setting prevails over starring if both
are used. 

The star can also be used with the \verb`free` setting. As mentioned
above, if the main argument is empty, then \emph{all} macros are freed:
\begin{centred}
\verb`\macros*[free]{}` $\Longrightarrow$ \macros*[free]{ }
\end{centred}

\subsection{Errors}

\label{subsec:supplMacrosErrors}If a macro is used in a \verb`\macros`
statement and the macro has not been defined in the document or a
supporting package it will cause an error:
\begin{verbatim}
    \macros{ \mymacro }
\end{verbatim}
$\Longrightarrow$ \macros{\mymacro }

\noindent As noted in the introduction to this section, an undefined
macro used in an \verb`\eval`-uation will cause an `Unknown token'
message in \verb`numerica`. The solution in this and the preceding
case is (obviously) to define the macro.

If a macro contains a formula which begins with an expandable token
and a preceding space is omitted (see above), then entering that macro
in a \verb`\macros` statement to register it with \verb`numerica`
will generally cause a puzzling error:
\begin{verbatim}
    \newcommand\mysin{\sin(\pi/7)}
    \macros{ \mysin }
\end{verbatim}
$\Longrightarrow$     \newcommand\mysin{\sin(\pi/7)}
    \macros{ \mysin }

\noindent The \verb`\protect` seems to be plucked from nowhere. In
fact it comes from the expansion of \verb`\sin`. If \verb`\sum`
had been the first token in the macro definition, again with no preceding
space, then \verb`\protect` would have been replaced by the even
more puzzling \verb`\DOTSB`. The solution is to insert a space as
the first token in the macro definition.

If a macro is defined but the \verb`\macros` statement is overlooked,
and the macro is used in an \verb`\eval`-uation, it will generate
an `Unknown token' message.

If your macro stores a formula with variables, and you forget to give
those variables values in the \verb`\macros` statement that will
produce a message:
\begin{verbatim}
    \def\mysumk{ \sum_{n=1}^k n }
    \macros{ \mysumk }
\end{verbatim}
$\Longrightarrow$     \def\mysumk{ \sum_{n=1}^k n } 
    \macros{ \mysumk }

\noindent The `where' part of the message is specific in this case,
but is generally `\verb`\nmcMacros` command'. 

And of course there can be `all the usual suspects' discussed at
§\ref{sec:evalErrors} in the evaluation of the vv-list or the formula.

\noindent\begin{minipage}[t]{1\columnwidth}%
\begin{shaded}%

\subsubsection{Display of macros}

\label{subsec:supplMacrosDisplay}As shown in earlier examples, macros
display as their content. Thus \verb`\mysumC` displayed as $\sum_{n=1}^{100}1/n-\ln100$.
But once a macro is known to \LaTeX{} (not necessarily to \verb`numerica`)
it can be used as a variable name. This has the same potential for
abuse as noted earler for multi-token variables (§\ref{subsec:evalDon't-do-this!}).
In the following example note that there is no \verb`\macros` statement.
It suffices for the macro to be known to \LaTeX . 
\begin{verbatim}
    \def\mymac{1}
    \eval[vvi=\,???]{$ \mymac+\mymac $}[\mymac=2]
\end{verbatim}
$\Longrightarrow$     \def\mymac{1}
    \eval[vvi=\,???]{$ \mymac+\mymac $}[\mymac=2]

The value assigned to a variable name – in this case \verb`\mymac`
– by \verb`numerica` for \cprotect\emph{calculational} purposes and
how that variable name \cprotect\emph{displays} in \LaTeX{} are two
separate things. One relies on the user not to do something deliberately
deceptive.\end{shaded}%
\end{minipage}

\subsection{Rounding value}

\label{subsec:supplMacrosRounding}Values are stored to $16$ significant
figures (if available). In most cases appending a rounding value to
a \verb`\macros` statement has no effect on the value stored. In
the following example note the \verb`o` setting, meaning the sine
reads angles in degrees:
\begin{verbatim}
    \NewDocumentCommand\testi{}{ \sin 60 }
    \NewDocumentCommand\testii{}{ \sin 60 }
    \macros[o]{ \testi }[10]
    \macros[o]{ \testii }[3]
    \macros[view]{}
\end{verbatim}
$\Longrightarrow$     \NewDocumentCommand\testi{}{ \sin 60 }
    \NewDocumentCommand\testii{}{ \sin 60 }
    \macros[o]{ \testi }[10]
    \macros[o]{ \testii }[3]
    \macros[view]{}

\noindent Despite the different rounding values the same $16$ figures
are stored in both \verb`\testi` and \verb`\testii`. 

For the \verb`\eval` command, rounding values specify how results
are \emph{displayed}. The rounding value matters only \emph{after},
not during, the calculation. Only for infinite sums or products is
this otherwise. There the rounding value is used to determine when
to stop adding further terms or factors. The same is true of the \verb`\macros`
command. Only if a macro contains an infinite sum or product does
the rounding value become relevant. Sixteen figures are still stored,
but most of them will be `wrong' since the infinite sum or product
has stopped early, after only a finite number of terms or factors.
Exactly how many of the first few figures are correct depends on the
rounding value. An example may clarify the matter.
\begin{verbatim}
    \macros[free]{}
    \def\zetaiii{ \sum_{n=1}^\infty 1/n^3 }
    \macros[view]{ \zetaiii }[3]
    \info{sum}
    \macros[view]{ \zetaiii }[6]
    \info{sum}
\end{verbatim}
$\Longrightarrow$     \macros[free]{} 
    \def\zetaiii{ \sum_{n=1}^\infty 1/n^3 }
    \macros[view]{ \zetaiii }[3]
    \info{sum}
    \macros[view]{ \zetaiii }[6]
    \info{sum}

\emph{HMF }Table 23.3 tells me that $\zeta(3)=1.202056903159594\dots$
The different rounding numbers have restricted the infinite sums to
the very finite $47$ and $468$ terms respectively. Although $16$
figures are stored, only the first few are correct. Just how many
depends on the number of terms summed which depends on when the stopping
criterion is met which depends on the rounding value.

\section{User-defined constants:\texttt{ \textbackslash nmcConstants}}

\label{sec:supplConstants}As noted much earlier in this document
(§\ref{subsec:Built-in-Constants}), there are five built-in constants:
\verb`\pi`, \verb`e`, \verb`\phi`, \verb`\gamma` and \verb`\deg`,
but a user may well want to define their own constant or constants.
There are contexts where it would make sense to permanently record
fundamental constants like the speed of light or Planck's constant,
or more down-to-earth constants like the acceleration due to gravity
or the viscosity of water, rather than having to enter them in the
vv-list for each calculation. Or a parameter might be held constant
for a particular problem or class of problems where other variables
change – for example triangles of constant perimeter but varying sides.
This is the purpose of the \verb`\nmcConstants` command.

The symbols used to denote constants are subject to exactly the same
constraints and freedoms as the symbols used to denote variables.
They might be single latin letters like \verb`c` (e.g. $c=3\times10^{8}$),
or greek letters like \verb`\alpha` (e.g. $\alpha=1/137)$, or multitoken
combinations like the Rydberg constants \verb`R_\infty` or \verb`R_{\mathrm{H}}`
from atomic physics, or \verb`\mu_0` and \verb`\epsilon_0` used
to denote the permeability and permittivity of free space, or personal
constants like \verb`total` of no wider significance. \verb`numerica`
handles all these different forms of constant with the command \verb`\nmcConstants`:
\begin{verbatim}
    \nmcConstants{  const-n=value-n, ... ,
                    const2=value2, const1=value1 }
\end{verbatim}
This is the simplest use – each constant is assigned a (numerical)
value. But it is easy to envisage situations where it would be convenient
to have a constant with value $1/\sqrt{2\pi}$ say, or another with
value $e^{\tfrac{\pi}{2}}$, and so on. That is easy: simply put the
expession for the value on the right:
\begin{verbatim}
    \constants{ a=1/\sqrt{2\pi},b=e^{\tfrac\pi2} } 
\end{verbatim}
Or the values could be expressions depending on parameters:
\begin{verbatim}
    \constants{ s=\tfrac12(a+b+c) }[a=3,b=5,c=7]
\end{verbatim}
Some constants might depend on earlier constants in the list:
\begin{verbatim}
    \constants{ A=\sqrt{s(s-a)(s-b)(s-c),
                s=\tfrac12(a+b+c) }[a=3,b=5,c=7]
\end{verbatim}
Or the values could involve an `infinite' process, requiring a rounding
number:
\begin{verbatim}
    \constants{ \zeta=\sum_{n=1}^\infty(1/n^k) }[k=4][5]
\end{verbatim}
In this, although $16$ figures will be stored, only the first few
will be accurate, the precise number depending on the value of \verb`k`
and the rounding number (\verb`5` in the example); see the discussion
on this issue for user-defined macros, §\ref{subsec:supplMacrosRounding}.

\subsection{New list replaces old by default}

A particular group of constants may be relevant only to a particular
part of a document. Another part of the document may use other constants.
By default, a second list of constants \emph{replaces }the first list.
Thus each of the \verb`\constants` statements above would replace
the previous one.

There is a technical reason why replacing rather than appending is
the default. For each calculation all \emph{multi-token} constants
(e.g. \verb`R_\infty`, \verb`N_0`, \ldots ) are added internally
to the start of the vv-list of the \verb`\eval` command. Even if
the vv-list is empty, this is still the case since the formula might
well use constants. Like multi-token variables and for the same reason
(see §\ref{subsec:evalVariableNames}), multi-token constants are
mapped internally to single tokens. This occurs afresh for each calculation.
If there are a lot of multi-token constants then each calculation
is going to involve not only this mapping from multi- to single tokens
but the evaluation of a long vv-list. In that case it seems better
to make the default behaviour replacement of one constant list by
another, rather than appending them. 

\subsection{Adding constants to a list}

Despite the default behaviour, there will be occasions when you want
to add a new constant or constants to the current list. This is easily
done with the \verb`add` setting. For instance,
\begin{verbatim}
    \nmcConstants[add]{ \sigma=5.67\times10^{-8},
                          k_B = 1.381\times10^{-23} }
\end{verbatim}
would add \verb`\sigma` and \verb`k_B` to the current list. The
presence of the \verb`add` setting triggers appending rather than
replacement.

\subsection{Examples of use}

\subsubsection{Example 1: atomic constants}

In the following example, the values of various atomic constants are
taken from the \verb`mandi` package. I use two \verb`\constants`
statements in order to show the use of the \verb`add` setting. I've
also included a \verb`view` setting in the second \verb`\constants`
statement.

The constants are used to calculate the fine-structure constant \verb`\alpha`
in the vv-list of the \verb`\eval` command, and its well-known reciprocal
(close to $137)$ in the main argument. Note that the constants do
not need to be entered in the vv-list of the \verb`\eval` command.
Their values are available from the \verb`\constants` statements.
\begin{verbatim}
    \constants{ c=2.99792458\times10^{8},
                h=6.62607015\times10^{-34},
                e=1.602176634\times10^{-19} }
    \constants[view,add]
      { \epsilon_0=8.854187817\times10^{-12} }
    \eval{$ 1/\alpha $}[\alpha=e^2/2\epsilon_0hc]
\end{verbatim}
$\Longrightarrow$     \constants{ c=2.99792458\times10^{8},
                h=6.62607015\times10^{-34},
                e=1.602176634\times10^{-19} }
    \constants[view,add]{
       \epsilon_0=8.854187817\times10^{-12} }
    \eval{$ 1/\alpha $}[\alpha=e^2/2\epsilon_0hc].

The \verb`view` setting produces a now familiar kind of display.
It shows that the three-token \verb`\epsilon_0` (the control sequence
\verb`\epsilon`, the underscore \verb`_` and the digit \verb`0`)
has been replaced by \verb`\nmc_q` – which may look as if it is also
three tokens but is in fact a single control sequence.

\subsubsection{Example 2: local constants}

Long ago, when there were such creatures as reference librarians,
I was asked about a school physics project along these lines.

A\emph{ car is travelling at 50 km/hr when it hits a lamppost. The
bonnet crumples 1 metre and the car comes to an immediate halt. Although
she herself is wearing a seat-belt, a woman in the passenger cabin
is holding her 5 kg baby. Does the baby survive?}

The enquirer was familiar with the equations describing constant acceleration,
\[
x=ut+\tfrac{1}{2}at^{2},\quad\text{and}\quad v^{2}-u^{2}=2ax,
\]
and Newton's second law, $F=ma$, force equals mass times acceleration.
The question was really about understanding these laws and how to
think with them. Here, $s$ is the distance travelled in time $t,$
with initial speed $u$ at $t=0$, speed $v$ at time $t$, and constant
acceleration $a$ – a deceleration in this case. 

The given data provide our constants: distance $x=1$ metre, initial
speed $u=1000*50/(60*60)=(10/36)*50$ metres per second, final speed
$v=0$. To estimate whether the woman can hold on to her baby, we
will need to make a comparison with forces we have personally experienced.
Most of us have tried lifting someone else, so let's use a characteristic
human weight as our test mass. Thus, we have the (baby's) mass $m=5$
kilograms, and a test mass, $M$ say, which we will leave as a variable.
But dealing with weight, we will need the acceleration due to gravity.
For the kind of rough estimating we are doing, $g=10$ metres per
second per second will be an adequate approximation.
\begin{verbatim}
    \constants{ x=1,v=0,u=(10/36)50,m=5,g=10 }
\end{verbatim}
The deceleration experienced by the woman is found from the second
equation of constant acceleration, $a=(v^{2}-u^{2})/2x$. Even if
the deceleration isn't constant this will give an estimate of its
magnitude. (If some of the deceleration is less than this $a$, some
must be greater.) This is also the deceleration experienced by the
baby as long as the woman holds onto her. Hence the magnitude of the
force exerted by the baby on the woman's arms is $ma=m(v^{2}-u^{2})/2x=-mu^{2}/2x$
which we want to compare with our test force, say that required to
lift $M=70$ kilograms, which was once considered the mass of an average
western adult male (but is doubtless a considerable underestimate
now). Hence the test force is $Mg$. Let's do the calculations. (I
have altered the \verb`\constants` statement to allow for a later
comparison with the effect of a small increase in speed.)
\begin{verbatim}
    \constants{ x=1,u=(10/36)U,m=5,g=10 }[U=50]
    \eval{$ mu^2/2x $}[0], \par
    \eval{$ Mg $}[M=70].
\end{verbatim}
$\Longrightarrow$     \constants{ x=1,u=(10/36)U,m=5,g=10 }[U=50]
    \eval{$ mu^2/2x $}[0], \par
    \eval{$ Mg $}[M=70].

The force required to hold on to the baby is noticeably less than
that required to lift a $70$~kg person – in fact about the same
as that needed to lift a $50$~kg person. But we have ignored the
force experienced by the mothers forearms – perhaps doubling $m$
(baby plus forearms) would give a better estimate of the force she
experiences. In that case $mu^{2}/2x$ obviously doubles and the total
force required by the woman to retain her baby – now $964$ newtons
– is significantly more than that required to lift a $70$~kg person.
I think it almost certain that the baby is torn from her arms. 

What difference does increasing the speed to 60 km/hr make? 
\begin{verbatim}
    \constants{ x=1,u=(10/36)U,m=5,g=10 }[U=60]
    \eval{$ mu^2/2x $}[1], \par
    \eval{$ Mg $}[M=70].
\end{verbatim}
$\Longrightarrow$     \constants{ x=1,u=(10/36)U,m=5,g=10 }[U=60]
    \eval{$ mu^2/2x $}[0], \par
    \eval{$ Mg $}[M=70].

Now the force of baby alone is comparable to that required to lift
a $70$ kg person. Including the woman's forearms in $m$, doubling
$m$ say, will result in a force twice as great – like that required
to lift two $70$~kg people or one $140$~kg person. There is no
chance of the woman holding on to her baby. The force is too great. 

\subsubsection{Example 3: macros and constants}

Constants can depend on previously defined and registered user macros.
Suppose I have defined two macros
\begin{verbatim}
    \NewDocumentCommand\electronmassprecisevalue {}
                       {9.1093837015\times10^{-31}}
    \NewDocumentCommand\protonmassprecisevalue {}
                       {1.672621898\times10^{-27}}
\end{verbatim}
(I have taken both the names and the values from the \verb`mandi`
package.) The long explicit names of the macros has a pedagogic purpose,
but they are too cumbersome to use in calculations. For that purpose
we need, first, a \verb`\macros` statement registering the two macros
with \verb`numerica`, and then a \verb`\constants` statement like
\begin{verbatim}
    \nmcConstants{ m_e=\electronmassprecisevalue,
                   m_p=\protonmassprecisevalue }
\end{verbatim}
With that \verb`m_e` and \verb`m_p` could be entered in formulas,
taking the values contained in the macros. Let's do it:
\begin{verbatim}
    \NewDocumentCommand\electronmassprecisevalue {}
                       {9.1093837015\times10^{-31}}
    \NewDocumentCommand\protonmassprecisevalue {}
                       {1.672621898\times10^{-27}}
    \nmcMacros{ \electronmassprecisevalue,
                \protonmassprecisevalue }
    \nmcConstants{ m_e=\electronmassprecisevalue,
                   m_p=\protonmassprecisevalue }
    \eval{$ m_p/m_e $}
\end{verbatim}
$\Longrightarrow$ \NewDocumentCommand\electronmassprecisevalue {}
                       {9.1093837015\times10^{-31}}
\NewDocumentCommand\protonmassprecisevalue {}
                       {1.672621898\times10^{-27}}
\nmcMacros{ \electronmassprecisevalue,
                \protonmassprecisevalue }
\nmcConstants{ m_e=\electronmassprecisevalue,
                   m_p=\protonmassprecisevalue }
\eval{$ m_p/m_e $},

\noindent the familiar mass ratio of proton and electron.

\subsection{Viewing, counting constants}

To see all constants currently `in play', use the \verb`view` setting
in the \verb`\constants` command. The main argument can be empty,
\begin{centred}
\verb`\constants[view]{}` $\Longrightarrow$ \constants[view]{}
\end{centred}
or contain a list of constants. In the latter case, the display is
of the above form but featuring the constants of the new list or,
if the \verb`add` setting is used, featuring the joined lists, old
and new:
\begin{centred}
\verb`\constants[view,add]{X=42}` $\Longrightarrow$ \constants[view,add]{X=42}
\end{centred}
To count how many constants are currently in play, star the \verb`\constants`
command. The number will depend on whether the main argument is empty
or not, and whether the \verb`add` setting is active:
\begin{centred}
\verb`\constants*{}` $\Longrightarrow$ \constants*{}.
\end{centred}
If the \verb`view` setting is being used at the same time as the
star, the \verb`view` prevails.

\subsection{Errors}

When contemplating error messages from \verb`numerica` it needs to
be remembered that \emph{multi-token} constants are added to the vv-list
for every calculation. Hence an error may not be in the vv-list as
indicated in the message but in the \verb`\constants` statement,
specifically, the multi-token constants. 

\section{Saving and reusing results: \texttt{\textbackslash nmcReuse}}

\label{sec:supplReuse}You may want to use at some place in a document
a result calculated earlier. It would be good to be able to do so
without having to do the calculation again at the new location. \texttt{numerica}
offers a command \verb`\nmcReuse` (short-name form, \verb`\reuse`)
which saves a result to a control sequence that can then be used elsewhere
in the document, expanding to the saved result. The control sequence
and its content are also saved to file, allowing the possibility of
using the result in other documents. 

\noindent\begin{minipage}[t]{1\columnwidth}%
\begin{shaded}%
The \texttt{\textbackslash nmcReuse} command in version 2 of \texttt{numerica
}has been completely rewritten. Its use is not compatible with how
the command was used in version 1. I found that I could bring \texttt{\textbackslash nmcReuse}
along with \texttt{\textbackslash nmcMacros} and \texttt{\textbackslash nmcConstants}
into the coding scheme used for \texttt{\textbackslash nmcEvaluate}
and the reasons for doing so were too compelling. \end{shaded}%
\end{minipage}

\subsection{Use of \texttt{\textbackslash nmcReuse}}

As noted, all the supplementary commands share the syntax of the \verb`\eval`
command, so that \verb`\nmcReuse` has an optional settings argument
preceding a mandatory main argument, followed by two trailing optional
arguments. \verb`\nmcReuse` does not use the last two. The command
is used mainly in two ways:
\begin{enumerate}
\item {\small\verb`\nmcReuse{}`}{\small , which loads the saved control
sequences from file, if not already loaded; and}{\small\par}
\item {\small\verb`\nmcReuse{csname}`}{\small , which loads the saved
control sequences from file, if not already loaded, assigns the latest
result from }{\small\verb`\eval`}{\small{} to the control sequence
}{\small\verb`\csname`}{\small , and saves }{\small\verb`\csname`}{\small{}
to file.}{\small\par}
\end{enumerate}
You may wish to put \verb`\nmcReuse{}` in the preamble of your document
(\emph{after} \verb`\usepackage{numerica}` of course). In that way,
saved control sequences are available from the start. Indeed, control
sequences saved from later in the document can be used in earlier
sections in a later \LaTeX{} run.

Note that only the \emph{name}, \verb`csname`, of the control sequence
is supplied to \verb`\reuse`, not the control sequence (\verb`\csname`).
The name should be composed of letters only. If the name has already
been defined in \LaTeX{} a \verb`numerica` error is produced, see
below §\ref{subsec:reuseDeletingOverwriting:}, although if you want
to save a \emph{new} value in a previously saved control sequence,
that can be done without invoking a message; see §\ref{subsec:reuseDeletingOverwriting:}.

Once defined {\small with a }{\small\verb`\nmcReuse{csname}`}
command, \verb`\csname` becomes available for use elsewhere in the
document.

\subsubsection{What is saved?}

What is saved is the most recent result of an \verb`\eval`-uation.
This is the \emph{full} result. It may include the vv-list; it may
include formatting elements; it may include math delimiters. Thus,
using \verb`\csname` in your document (after the command \verb`\nmcReuse{csname}`)
may not be straightforward – simply writing \verb`\csname` where
you want the value it expands to, may produce a \LaTeX{} error and
halt compilation. You may have to write \verb`$ \csname $` or provide
some other math environment in order for the control sequence to display
correctly.

It can be helpful to see \emph{exactly} what has been saved; to do
that see §\ref{subsec:reuseSeeingSaved}.

\paragraph{Use of \texttt{\textbackslash eval{*}}}

Users will make life simpler for themselves if they make a habit of
using the starred form \verb`\eval*` to produce the results to save.
\verb`\eval*` produces solely a number with no formatting or delimiters;
even a negative result uses a hyphen for the minus sign, just as one
would type it. In this case \verb`\csname` can be used freely in
both text and math environments.

\subsubsection{The \texttt{.nmc} file}

\label{subsec:suppleReuse.nmc-file}The file that control sequences
are saved to has a filename composed of the document name with the
extension \verb`.nmc`. If your document is \verb`mydoc.tex` (so
that the \LaTeX{} command {\small\verb`\jobname`} expands to
\verb`mydoc`) then the file to which results are saved is \verb`mydoc.nmc`,
located in the document directory. 

\verb`mydoc.nmc` is a comma list of pairs of the form \verb`\csname {value}`.
Thus, the contents of \verb`mydoc.nmc` might be \texttt{\textbackslash csname1
\{value1\},\textbackslash csname2 \{value2\},..., \textbackslash csname$n$
\{value$n$\}}. If \verb`mydoc.nmc` does not already exist then it
is created in the document directory, and \verb`\csname {value}`
becomes its first element.

\paragraph{Editing the \texttt{.nmc} file externally}

The \verb`.nmc` file is a text file and can be edited in a text editor.
Thus it is possible to externally add control sequences and values
to it provided the structure of the file is strictly adhered to. It
is also possible to delete items from it or rename control sequences
or edit values by the same mechanism. Editing the file externally
like this, or renaming it, or transferring items from one \verb`.nmc`
file to another, provides a way of using saved values in multiple
documents.

\subsubsection{Messages}

If a control sequence \verb`\csname` is already known to \LaTeX ,
then writing \verb`\reuse` \verb`{csname}` will produce a \verb`numerica`
message and the result of the latest \verb`\eval`-uation will \emph{not}
be saved:
\begin{verbatim}
    \eval*{\sum_{n=1}^{10}n}\par
    \reuse{sigma}
\end{verbatim}
$\Longrightarrow$ \eval*{\sum_{n=1}^{10}n} \par \reuse{sigma}

If there is no result to save – perhaps an \verb`\eval`-uation produces
an error message instead – then another message is generated:
\begin{verbatim}
    \eval*{1/0}\par
    \reuse{oops}
\end{verbatim}
$\Longrightarrow$     \eval*{1/0}\par
    \reuse{oops}

\subsubsection{Deleting and renewing}

\label{subsec:reuseDeletingOverwriting:}There may be occasions when
you wish to change a previously saved value and yet, irritatingly,
the control sequence name will now be known to \LaTeX{} and so will
generate an `already known' message. If you choose a different name
for the control sequence to save the new value to, do you want the
old name cluttering the \verb`.nmc` file? Deleting and renewing the
values of saved control sequences are controlled by the settings \verb`delete`
and \verb`renew`.

Entering \verb`delete` in the settings option \emph{deletes} a control
sequence and its value from the \verb`.nmc` file and undefines it
in \LaTeX{} terms. Thus \verb`\reuse[delete]` \verb`{csname}` would
delete \verb`\csname` and its value from the \verb`.nmc` file and
undefine \verb`\csname`. If \verb`\csname` is not present in the
file, nothing happens. Entering \verb`renew` replaces the value of
a saved control sequence with a new value. If there is no such \emph{saved}
control sequence but the control sequence is otherwise known to \LaTeX{}
the `already defined' message will still be generated. This prevents
giving control sequences like \verb`\sin`~or \verb`\frac` new meanings
with the \verb`renew` setting.
\begin{itemize}
\item \verb`\reuse[delete]{csname}` deletes \verb`\csname` and its value
from the \verb`.nmc` file and from memory if present; otherwise has
no effect;
\item \verb`\reuse{csname}` (the default) saves the result of the latest
\verb`\eval` command to \verb`\csname`, provided \verb`\csname`
is not already defined; in that case a warning message is presented
and the result is not saved;
\item \verb`\reuse[renew]{csname}` behaves like the default mode unless
\verb`\csname` is already a saved control sequence in the \verb`.nmc`
file, in which case its previous value is replaced by the result of
the latest \verb`\eval` command; 
\item if \verb`delete` and \verb`renew` are used together, whichever occurs
second prevails.
\end{itemize}
In the following example, the first \verb`\reuse` deletes \verb`\suma`
should it be present in the \verb`.nmc` file, the second saves the
result, $55$, of the latest \verb`\eval`-uation (in fact an \verb`\eval*`-uation)
and the third overwrites that saved value with the new value, $210$. 
\begin{verbatim}
    \reuse[delete]{suma}
    \eval*{\sum_{n=1}^{10}n} \par
    \reuse{suma}
    \eval*{\sum_{n=1}^{20}n} \par 
    \reuse[renew]{suma}
\end{verbatim}
$\Longrightarrow$ \reuse[delete]{suma}
\eval*{\sum_{n=1}^{10}n} \par 
\reuse{suma}
\eval*{\sum_{n=1}^{20}n} \par 
\reuse[renew]{suma}

\subsubsection{Viewing what has been saved}

\label{subsec:reuseSeeingSaved}It would be good in this example to
see that the new value $210$ has in fact been saved. That is easy.
Simply enter \verb`view` in the settings option of \verb`\nmcReuse`
(I've removed the now unnecessary \verb`\par` tokens from the example.)
\begin{verbatim}
    \reuse[delete]{suma}
    \eval*{\sum_{n=1}^{10}n}
    \reuse[view]{suma}
    \eval*{\sum_{n=1}^{20}n} 
    \reuse[renew,view]{suma}
\end{verbatim}
$\Longrightarrow$     \reuse[delete]{suma}
    \eval*{\sum_{n=1}^{10}n}
    \reuse[view]{suma}
    \eval*{\sum_{n=1}^{20}n}
    \reuse[renew,view]{suma}

\noindent First the original value $55$ was saved to \verb`\suma`
but then the value was overwritten by the new value $210$. 

The \verb`view` setting allows us to see how formatting is stored
if the \emph{un}starred form of the \verb`\eval` command is used.
In the following example, \verb`\eval` wraps around math delimiters:
\begin{centred}
\verb`\eval{$ 1+1 $} \reuse[view,renew]{two}` $\Longrightarrow$
\eval{$ 1+1 $} \reuse[view,renew]{two}
\end{centred}
The full \emph{formula=result} display has been captured in \verb`\two`
along with the math delimiters. If a vv-list is also involved, things
become messy (but informative):
\begin{verbatim}
    \eval{$ x+y $}[x=1,y=2]
    \reuse[view,renew]{three}
\end{verbatim}
$\Longrightarrow$     \eval{$ x+y $}[x=1,y=2]
    \reuse[view,renew]{three}

You may want to see \emph{all} saved control sequences. In that case
use an \emph{empty} main argument: \verb`\nmcReuse[view]{}`. We now
have enough saved control sequences to make this worthwhile: 
\begin{verbatim}
    \reuse[view]{}
\end{verbatim}
$\Longrightarrow$ \reuse [view]{} 

\noindent (The \verb`\seven` that appears here is defined shortly.
Its appearance \emph{before} definition is presumably due to \LaTeX{}
making a number of passes when compiling this document.)

\subsubsection{Counting saved control sequences: \texttt{\textbackslash nmcReuse{*}}}

Because \verb`\nmcReuse` uses the same machinery as \verb`\nmcEvaluate`,
it has a starred form, \verb`\nmcReuse*`, which produces a purely
numerical result (just like \verb`\eval*`, \verb`info*`, \verb`\macros*`
and \verb`\constants*`). In this case, the number is the count of
how many control sequences have been saved:
\begin{centred}
\verb`\reuse*{}` $\Longrightarrow$ \reuse*{ }.
\end{centred}

\subsection{\texttt{reuse} setting of \texttt{\textbackslash eval} command}

\label{subsec:supplReuseEvalSetting}Using \verb`\eval*` for a calculation
ensures a purely numerical result, with no vv-list or formatting in
the display of the result. But sometimes we might want the full display
yet wish to save only the numerical result. This is the point of the
\verb`reuse` setting of the \verb`\eval` command. 

For the \emph{starred} form of the \verb`\eval` command it is always
\emph{only the numerical result} that is saved, whatever the value
of the \texttt{reuse} key in the settings option of the \verb`\eval`
command.

For the \emph{unstarred} form of the \verb`\eval` command exactly
what is saved with \verb`\nmcReuse` depends on the \texttt{reuse}
setting:
\begin{lyxcode}
reuse~=~<integer>
\end{lyxcode}
where \verb`<integer>` can take one of two values,
\begin{itemize}
\item \texttt{reuse=0} (the default) saves\emph{ the form that is displayed}
including the vv-list if there is one and possibly a formatting component
(like math delimiters). Note that if the vv-list is empty, a formatting
component (math delimiters) may still be present in the saved result; 
\item \texttt{reuse=1} (or, indeed, any non-zero integer) saves only the
numerical result with no other elements of the display (no vv-list,
no formatting component, no math delimiters).
\end{itemize}
As we saw earlier, saving the result from \verb`\eval{$ x+y $}[x=1,y=2]`,
corresponding to \verb`reuse=0`, means the full display is saved.
Check by writing \verb`\three` $\Longrightarrow$ \three. The full
display was saved (including math delimiters).

On the other hand, with {\ttfamily\verb`reuse=1`} only
the numerical value is saved:
\begin{centred}
\verb`\eval[reuse=1]{$ x + y $}[x=3,y=4] \reuse[renew]{seven}` $\Longrightarrow$
\eval[reuse=1]{$ x + y $}[x=3,y=4] \reuse[renew]{seven}.
\end{centred}
The numerical result only of the calculation should be saved, although
the formula and vv-list are displayed as the result of the \verb`\eval`-uation.
We can easily check: \verb`\seven` $\Longrightarrow$ \seven. Indeed,
only the numerical result was saved.

\chapter{Nesting  commands}

\label{chap:Nesting} The \verb`\eval` command and the supplementary
commands of the previous chapter can be \emph{nested} –\emph{ }used
within other \verb`\eval` or supplementary commands. Nesting may
occur in the main argument, or the vv-list, or the settings option,
or some combination of all three. With the commands currently introduced,
nesting is unlikely to be a major concern, but it becomes significant
for the commands defined in the associated package \texttt{numerica-plus}
(see §\ref{subsec:Related-packages}). Since those additional commands
are not available for this document, the examples below use the commands
introduced earlier: \verb`\eval`, \verb`\info`, \verb`\macros`,
\verb`\constants` and \verb`\reuse`.

\section{Nesting in the formula}

Consider a statement like \verb`\eval{...\eval...}`. There is an
inner \verb`\eval` and an outer \verb`\eval`. The inner \verb`\eval`
`digests' \emph{its} \LaTeX{} formula to produce an \verb`l3fp`-readable
expression which is fed to \verb`l3fp` to evaluate. The result is
then returned to (the inner) \verb`\eval` to display. In version
1 of \verb`numerica` that meant the inner command \emph{had} to be
starred, \verb`\eval*`, so that no display formatting was fed to
the outer command to try to digest (and cause an error).\emph{ }In
version 2 of \verb`numerica` this is no longer the case. \verb`numerica`
detects whether a command is inner or outer, and if inner, suppresses
all display formatting, producing only a number, as if the command
had been starred: 
\begin{centred}
\verb`\eval{$ \sin(\eval{\sin x}[x=\pi/6]\pi) + 1 $}` $\Longrightarrow$
\eval{$ \sin(\eval{\sin x}[x=\pi/6]\pi) + 1$}.
\end{centred}
In the presentation of the overall result, the inner \verb`\eval`
command is evaluated, displaying as a number.

In this example, the \verb`x=\pi/6` could be removed from the inner
\verb`\eval` and placed in the vv-list of the outer command since
outer variables are available to the inner command:
\begin{centred}
\verb`\eval{$ \sin(\eval{\sin x}\pi) + 1 $}[x=\pi/6]` $\Longrightarrow$
\eval{$ \sin(\eval{\sin x}\pi) + 1$}[x=\pi/6].
\end{centred}
Just to show that it is possible, the next example shows \verb`\eval`
being used in a \verb`\constants` command. The \verb`o` setting
in the \verb`\constants` command pervades its argument; hence it
needs to be explicitly turned off for the \verb`\eval` if \verb`\sin(\pi/6)`
is to evaluate as expected.
\begin{verbatim}
    \constants[o]{ y=\sin 30,x=\eval[o=0]{\sin(\pi/6)} }
    \eval{$ x+y $}
\end{verbatim}
$\Longrightarrow$     \constants[o]{ y=\sin 30,x=\eval[o=0]{\sin(\pi/6)} }
    \eval{$ x+y $}.

\subsection{Math delimiters and double evaluations}

Any math delimiters in the inner \verb`\eval` are ignored. (This
also differs from version 1 of \verb`numerica` where they caused
an error.) Obviously it is simpler to omit them as I have done in
the examples. 

However, math delimiters in the \emph{outer} \verb`\eval` command
still have their normal effect and produce a \emph{formula = result,
(vv-list)} display. One consequence of such a display is that the
formula in the \emph{inner} \verb`\eval` command is evaluated \emph{twice}
– once when the overall result is being calculated (i.e. the formula
of the \emph{outer} \verb`\eval`) and later when the overall display
of the result is created. In the \emph{formula} part of the \emph{formula
= result, {[}vv-list{]}} display, the tokens in the \emph{formula}
are expanded to their display form. For example, \verb`\sin` is expanded
to $\sin$, \verb`\pi` is expanded to $\pi$ – and the inner \verb`\eval`
is expanded to the numerical result of its evaluation – a second evaluation.
If the inner formula is simple, this will be of little moment, but
should the inner formula contain, say, a slowly converging infinite
series, then evaluating it twice is a bad idea and it would be better
to remove the delimiters from the outer \verb`\eval`. That prevents
the second evaluation. 

The problem does not arise if the outer \verb`\eval` lies within
a math environment (e.g. \verb`$ \eval{...} $`) since that produces
a display of the form \emph{result, {[}vv-list{]}.} The formula is
not displayed and so the second evaluation does not occur. The inner
\verb`\eval` is evaluated once only  to calculate the result.

\section{Nesting in the vv-list}

The inner \verb`\eval` can be placed in the vv-list of the outer
command. If the vv-list of the inner \verb`\eval` contains a comma
(meaning there are at least two variables), then the entire inner
\verb`\eval` and its \LaTeX{} arguments needs to be wrapped in braces
to hide the comma or commas of its vv-list from the outer \verb`\eval`.
To show the effect of not doing so, I have slightly complicated the
previous example by adding a second (unnecessary) variable. The first
example is with braces, the second without:
\begin{centred}
\verb`\eval{$ \sin k\pi + 1 $}[k={\eval{y\sin x}[x=\pi/6,y=1]}]`
$\Longrightarrow$ \eval{$ \sin k\pi + 1 $} [k={\eval{y\sin x}[x=\pi/6,y=1]}].

\verb`\eval{$ \sin k\pi + 1 $}[k=\eval{y\sin x}[x=\pi/6,y=1]]` $\Longrightarrow$
\eval{$ \sin k\pi + 1 $}[k=\eval{y\sin x}[x=\pi/6,y=1]].
\end{centred}
The vv-list of the outer \verb`\eval` is parsed as containing two
entries, \verb`k=\eval` \verb`{y\sin x}[x=\pi/6` and \verb`y=1]`.
Both will cause errors but since the vv-list is evaluated from the
right, it is \verb`y=1]` which actually does so. 

\section{Nesting in the settings option}

This will be rare, but commands can occur in the settings option of
the outer command. The \verb`\info` command provides a good example.
I have included it in the punctuation setting of an \verb`\eval`-uation. 
\begin{verbatim}
    \eval[p=\mbox{,\qquad\info{sum} terms.}]
      {\[ \sum_{n=0}^{\infty}\frac{(-1)^{n}}{n!} \]}[3] 
\end{verbatim}
$\Longrightarrow$ \eval[p=\mbox{,\qquad\info{sum} terms.}]{\[ \sum_{n=0}^{\infty}\frac{(-1)^{n}}{n!} \]}[3]

Because of the \verb`\[ \]` math delimiters, if the \verb`\info`
command had been placed \emph{after} the \texttt{\textbackslash eval}
command, it would have slid down to the next line. Used in the settings,
as here, the display is \emph{inside} the \verb`\[ \]` delimiters,
on the same line as the expression. This may be significant for adjusting
vertical spacing of later parts of the document – widow and orphan
control for instance.

A point to note is the explicit writing of the `terms' descriptor.
Normally \verb`\info{sum}` would automatically supply the descriptor,
but as noted earlier, nesting of one command in another suppresses
all elements of display of the inner command beyond the numerical
result. It is as if the inner command is starred. Because the \verb`\info`
command is nested in the \verb`\eval` command, the `terms' descriptor
is suppressed and has had to be explicitly supplied by hand.

\section{Rounding and display}

In the display of the overall result, it is the result of the inner
command which is shown, not the formula that the inner command acts
on. How that number is displayed is determined by the number-format
specification of the \emph{inner} command. Note however that this
specification affects only how the result of the inner command is
shown. Always 16 figures are passed from the inner command to the
outer, as you can see in this example:
\begin{verbatim}
    \eval{$ \pi - \eval{ \pi }[4] $}[15]
\end{verbatim}
$\Longrightarrow$ \eval{$ \pi - \eval{ \pi }[4] $}[15].

\noindent The outer result would not be zero to $15$ places of decimals
if the inner result were restricted to $4$ decimal places. It is
only the \emph{display} of the inner result which is so restricted.

For infinite sums and products (and for \verb`\nmcIterate` and \verb`\nmcSolve`
of the \verb`numerica-plus` package), the rounding value is not just
for display purposes but is also used to determine the result. This
may require judicious use of the extra rounding setting to get a sensible
display. In the first instance below, the second sum stops at an effective
rounding value of $5+2=7$, since the default extra rounding is $+2$,
and the first sum (the inner one) also stops at $7=4+3$. No surprise
then that the overall result is $0$. In the second instance, the
inner sum stops at a rounding value of $4+2=6$. Although the left-hand
side of the display is unaltered, the result is no longer $0$.
\begin{verbatim}
     \eval{$ \eval[S+=3]{\sum_{n=1}^\infty 1/n^3}[4*]
       - \sum_{n=1}^\infty 1/n^3$}[5]
\end{verbatim}
$\Longrightarrow$      \eval{$ \eval[S+=3]{\sum_{n=1}^\infty 1/n^3}[4*]
       - \sum_{n=1}^\infty 1/n^3$}[5]

\noindent whereas
\begin{verbatim}
    \eval{$ \eval[S+=2]{\sum_{n=1}^\infty 1/n^3}[4*]
       - \sum_{n=1}^\infty 1/n^3$}[5]
\end{verbatim}
$\Longrightarrow$      \eval{$ \eval[S+=2]{\sum_{n=1}^\infty 1/n^3}[4*]
       - \sum_{n=1}^\infty 1/n^3$}[5].

\section{Error messages}

Errors in an inner command create a small change in error message
display.
\begin{centred}
\verb`\eval{ 1 + \eval{ 1 + \eval{ k } } }` $\Longrightarrow$ \eval{ 1 + \eval{ 1 + \eval{ k } } }

\verb`\eval{ x + \eval{ k }[k=\arcsin 2] }[x=1]` $\Longrightarrow$
\eval{ x + \eval{ k }[k=\arcsin 2] }[x=1]
\end{centred}
An integer is added to the `where' part of the error message. The
integer indicates the \emph{level of nesting} where the error occurs.

If there is no nesting where the error occurs, the integer is suppressed,
even though there may be nesting elsewhere in the overall expression.
This is in the interests of straightforwardness when nesting is absent,
which will be overwhelmingly the most common situation.
\begin{centred}
\verb`\eval{ k + \eval{ x }[x=1] }[k=\arcsin 2]` $\Longrightarrow$
\eval{ k + \eval{ x }[x=1] }[k=\arcsin 2]
\end{centred}

\section{Debugging}

\label{subsec:nestDebugging}It is worth looking at the debug display
when \verb`\eval` commands are nested. For the outer \verb`\eval`
command:
\begin{centred}
\verb`\eval[dbg=1]{$ \sin \eval*{\sin x}[x=\pi/6]\pi + 1 $}` $\Longrightarrow$
\eval[dbg=210]{$ \sin \eval*{\sin x}[x=\pi/6]\pi + 1$}
\end{centred}
There is no vv-list for the outer command whence the two empty slots
in the display but when the inner \verb`\eval` is in the vv-list,
they are filled:
\begin{centred}
\verb`\eval[dbg=1]{$ \sin k\pi + 1 $}[k=\eval*{\sin x}[x=\pi/6]]`
$\Longrightarrow$ \eval[dbg=1]{$ \sin k\pi + 1$} [k=\eval*{\sin x}[x=\pi/6]]
\end{centred}
For the inner \verb`\eval` command debugging may still work but in
an idiosyncratic way. To clarify exactly what is going on I have added
a \verb`\left( \right)` pair around the entire inner \verb`\eval`
command. Note that I have also used a \emph{negative} \texttt{dbg}
value. With a positive value, the right parenthesis is pressed toward
the right margin of the page. The negative value limits the display
to the text width and gives the much neater result shown.
\begin{verbatim}
    \eval[()=2]{$
      \sin\left( 
            \eval*[dbg=-1]{ \sin x }[x=\pi/6]
          \right)\pi + 1 $}
\end{verbatim}
$\Longrightarrow$     \eval[()=2]{$
      \sin\left( 
            \eval*[dbg=-1]{ \sin x }[x=\pi/6]
          \right)\pi + 1 $}

\medskip{}
The debug display from the inner \verb`\eval` command has been inserted
into the formula of the outer \verb`\eval` in the position occupied
by the inner \verb`\eval`. I did not deliberately code for this,
but have decided to leave it as is despite the potential for some
rather odd displays, since there can be no confusion about which \verb`\eval`
command is being `debugged'. In this last example, in order to both
use \verb`\left(...\right)` and have the calculation give the previous
result I have employed the setting \verb`()=2` in the outer \verb`\eval`;
see §\ref{subsec:parseTrigFns}.

\chapter{Using \texttt{numerica} with \protect\LyX}

\label{chap:LyX}The document processor \LyX{} has a facility that
enables snippets from a document to be compiled separately and the
results presented to the user without having to compile the entire
document. The present document was written in \LyX . The demonstration
calculations were evaluated using this \emph{instant preview} facility. 

To use \texttt{numerica} in \LyX{} go to \textsf{Document \lyxarrow{}
Settings \lyxarrow{} LaTeX Preamble} and enter
\begin{lyxcode}
\textbackslash usepackage\{numerica\}
\end{lyxcode}
then click \textsf{OK}. You may wish to follow the above line in the
preamble with \verb`\nmcReuse{}`:
\begin{lyxcode}
\textbackslash usepackage{[}lyx{]}\{numerica\}

\textbackslash nmcReuse\{\}
\end{lyxcode}
In that case, type the extra line and \emph{then} click \textsf{OK}.
The additional line ensures all saved values are available in your
document from the outset.

\section{Instant~preview}

The instant preview facility of \LyX{} performs mini-\LaTeX{} runs on
selected parts of a document (for instance, the mathematical parts)
and displays the results in \LyX{} while the user continues to work
on the surrounding document.\texttt{ numerica} uses these mini-\LaTeX{}
runs to do its evaluations and display their results. That means you
get feedback on your calculations almost immediately.

To use this facility first ensure that instant preview is turned on.
This means selecting \textsf{Tools \lyxarrow Preferences \lyxarrow Look
\& Feel \lyxarrow{} Display}, ensuring that the \textsf{Display graphics}
checkbox is checked, and against \textsf{Instant preview} selecting
\textsf{On}, then clicking \textsf{OK}.

\subsection{Document location}

It also matters where your document is located. You may have your
own local or personal texmf tree (see §\ref{subsec:settingsPersonal-texmf-tree}).
If your document is located there, perhaps in the \verb`doc` folder,
then not all features of preview will work as expected. Presumably
this is because both \LyX{} and your \LaTeX{} distribution (e.g. \TeX Live
or MiK\TeX ) are interacting with the location and interfere. Move
your document to another location which your \LaTeX{} distribution
has no interest in, and open it in \LyX{} there.

\subsection{Global vs local previewing}

\label{subsec:LyXGlobal-vs-local}Compilation of previews occurs in
two distinct modes. 

\paragraph{Global preview generation:}

When a document is opened (and preview is \emph{on}), all previews
in the document are formed in sequence in the one \LaTeX{} run. This
is the global mode. The mini-\LaTeX{} run may well be substantial.
It compiles a \verb`.tex` file that begins with the document's preamble
with some additions then comes \verb`\begin{document}`. That is followed
by a sequence of preview environments,
\begin{lyxcode}
\textbackslash begin\{preview\}

<stuff>

\textbackslash end\{preview\}
\end{lyxcode}
one for each preview in the document. Finally there is an \verb`\end{document}`
statement. The critical point is that all previews are between the
same \verb`\begin{document}`, \verb`\end{document}` statements,
and so earlier previews in the sequence can communicate with later
ones.

\paragraph{Local preview generation:}

The other mode in which preview operates is local. Suppose you have
your document open and want to add to it, for instance with a simple
evaluation, \verb`\eval{x+y}[x=1,y=2]` in an ERT inset in a preview
inset. The resulting mini-\LaTeX{} run is of the form
\begin{lyxcode}
<preamble>

\textbackslash begin\{document\}

\textbackslash begin\{preview\}

\textbackslash eval\{x+y\}{[}x=1,y=2{]}

\textbackslash end\{preview\}

\textbackslash end\{document\}
\end{lyxcode}
The preamble is as before but there is only \emph{one} preview between
the \verb`\begin{document}`, \verb`\end{document}` statements. That
preview is isolated from all other, previous previews and will be
isolated from all other, later previews. 

This has implications for the supplementary commands of the previous
chapter and means that if you want to transfer information (a macro,
a constant, a result) from one preview to another, you need to do
it through the preamble or by means of an external file or, in some
cases, by forcing a global preview run in which all previews are recompiled
between the same \verb`\begin{document}`, \verb`\end{document}`
statements.

\subsubsection{Forcing a global preview run}

Closing then opening a document is one way to force a global preview
compilation. Another is to change the zoom level. This causes \LyX{}
to recompile all previews at the new zoom level. But you may not want
to work at the new zoom level. Going back to the old zoom level will
force a second recompilation of all previews. For a large document
\emph{two} recompilations is too heavy a burden. The secret is to
combine a zoom in and a zoom out into one command and attach it to
a shortcut.

If you go to \textsf{Tools \lyxarrow{} Preferences \lyxarrow{} Editing
\lyxarrow{} Shortcuts}, click on the \textsf{New} button and enter
\begin{lyxcode}
command-sequence~buffer-zoom-in;~buffer-zoom-out
\end{lyxcode}
then assign a shortcut to it (\verb`Alt+Z` for zoom?) you will gain
a simple means of forcing a global recompilation of previews. 

\section{Mathed}

(Mathed = the \LyX{} mathematics editor.) If you have instant preview
\emph{on} then one way to use \texttt{numerica} in \LyX{} is to enter
an \verb`\eval` command in mathed. Clicking the cursor outside the
editor with the mouse or moving it outside with the arrow keys will
then trigger formation of a preview of the editor's contents – a snippet
of what will be shown in the pdf. This will be displayed in mathed's
place after a generally short `pause for thought' as the mini-\LaTeX{}
run progresses behind the scenes. 

The original expression can be recovered by clicking on the preview.
The content of mathed is immediately displayed and can be edited.

\subsection{\protect\LaTeX{} braces~\{~~\}}

\LyX{} does not support \texttt{numerica}'s \verb`\eval` command `out
of the box' as it does, say, \verb`\frac` or \verb`\sqrt`. To use
the \verb`\eval` command in mathed you will need to supply the braces
used to delimit its mandatory argument. (For \verb`\frac` and \verb`\sqrt`
by contrast, \LyX{} supplies these automatically in the form of blue-outlined
boxes.) Unfortunately the \verb`{` key\footnote{\textsf{Shift+{[}} on my keyboard.}
does not insert a left brace into the document but rather an escaped
left brace \verb`\{` as you can see by looking at \textsf{View \lyxarrow{}
Code Preview Pane}. Escaped braces like this are used for grouping
terms in \emph{mathematics}; they are not the delimiters of a \LaTeX{}
argument. 

The brace delimiters for \LaTeX{} arguments are entered in mathed by
typing a backslash \textsf{\textbackslash{} }then a left brace\textsf{
\{} – two separate key presses rather than a single combined press.
This enters a balanced pair of (unescaped) braces with the cursor
sitting between them waiting for input. Alternatively, if you have
already written an expression that you want to place between braces,
select it, then type \textsf{\textbackslash{} }then\textsf{ \{}.

\section{Preview insets}

There are problems with using mathed for calculations.
\begin{itemize}
\item Expressions entered in mathed are necessarily of the form \verb`$ \eval... $`
or more generally \verb`delimiter` \verb`\eval...` \verb`delimiter`.
But you may wish to wrap the \verb`\eval` command \emph{around} the
math delimiters to produce a \emph{formula=result} form of display.
In mathed the only way to effect such a display is to write the \emph{formula=
}part yourself – which may involve no more than copy and paste but
is still additional mouse work/key pressing. 
\item Mathed does not accept carriage returns. If you want to format a complicated
expression for readability by breaking it into separate lines, you
can't. The expression is jammed into the one line, along with the
settings option content and the vv-list, often extending well beyond
the edge of the screen. 
\end{itemize}
For these reasons I have come to prefer \emph{not} using mathed for
calculations but instead to use preview insets wrapped around \TeX -code
(ERT) insets. \LyX{} uses the shortcut \textsf{Ctrl+L} to insert an
ERT inset. Since \LyX{} now does no printing itself, the shortcut \textsf{Ctrl+P}
that was formerly used for printing is available for other purposes.
On my keyboard, the \textsf{P} key lies diagonally up and to the right
but adjacent to the \textsf{L} key. I suggest assigning \textsf{Ctrl+P}
to inserting a preview inset. Then typing \textsf{Ctrl+P Ctrl+L} –
which means holding the \textsf{Ctrl} key down and tapping two diagonally
adjacent keys, \textsf{P} followed immediately by \textsf{L} – will
insert an ERT inset inside a preview inset with the cursor sitting
inside the ERT inset waiting for input. In the ERT inset you can enter
carriage returns, and so format complicated expressions. You can place
the vv-list on a separate line or onto consecutive lines. And when
you have finished, clicking outside the preview inset will trigger
preview into doing its thing and present the result `before your
eyes'. 

To assign the suggested shortcut, go to \textsf{Tools \lyxarrow{} Preferences
\lyxarrow{} Editing \lyxarrow{} Shortcuts}. Under \textsf{Cursor, Mouse
and Editing Functions} in the main window on the right, scroll down
until you come to \textsf{preview-insert}, select it, then click \textsf{Modify}.
Now press \textsf{Ctrl+P}. The shortcut will magically appear in the
greyed, depressed key.\textsf{ }Click \textsf{OK} and then \textsf{OK}
in the \textsf{Preferences} window to close it. (Most of the examples
in this document have been evaluated in this way, using \textsf{Ctrl+P
Ctrl+L.)}

\section{Errors }

Instant preview will display error messages generated by \verb`numerica`
in \LyX{} just as it does the results of calculations. Clicking on
the message will show the underlying expression which can then be
edited. However \LaTeX{} errors will \emph{not} produce a preview;
formation of the preview will stall. To find precisely what has gone
wrong, you will need to look at the \LaTeX{} log, but not the log of
the overall document; rather the \emph{preview} log. 

\subsection{Temporary directory of \protect\LyX}

Unfortunately this is tucked away in a temporary directory and is
not immediately accessible in \LyX{} (unlike the main \LaTeX{} log from
\textsf{Document \lyxarrow{} \LaTeX{} Log}). When \LyX{} is started,
it sets up a temporary directory in which to perform various tasks.
On Windows systems this will be located in \texttt{C:\textbackslash Users\textbackslash <your
name>\textbackslash AppData\textbackslash Local\textbackslash Temp}
and will have a name like \texttt{lyx\_tmpdir.XOsSGhBc1344}. 

One of the tasks \LyX{} uses this temporary directory for is to create
preview images when a document is opened. If you look inside \LyX 's
temporary directory when a document is first loaded, you will see
a subdirectory created, with a name like \texttt{lyx\_tmpbuf0}. There
may already be such directories there, in which case the number on
the end will be greater than \texttt{0} – it depends on whether other
documents are or have been open in the current instance of \LyX .
Inside the appropriate \texttt{lyx\_tmpbuf}\texttt{\emph{n}} folder
will be the preview log with a name like \texttt{lyxpreviewZL1344.log}.
It will usually be accompanied by other files with extensions like
\texttt{.dvi}, \texttt{.tex}, and – depending on the number of previews
in your document – a number, perhaps a lot, of image files with the
extension \texttt{.png}, each one of which is a preview. For a document
just loaded there will be only the one preview log, but if you have
added preview insets or math insets to your document\textsf{ }in the
current editing session there will be a number of such logs and you
will need to determine the relevant one by the time stamp. 

The log files are text files and can be opened in a text editor. The
relevant part of the log is towards the end (just before the final
statistical summary) where you will find a list of entries like \texttt{Preview: Snippet
1 641947 163840 7864588}. If there is an error, it will be noted here
among these snippets and will generally make clear what needs remedying.

\subsection{CPU usage, \protect\LaTeX{} processes}

It is possible when a preview stalls that the \LaTeX{} process associated
with the preview will continue to run, using CPU cycles, slowing overall
computer performance, and perhaps resulting in extra fan use giving
a different sound to the computer. In Windows 10, the \textsf{Task
Manager} (\textsf{Ctrl+Shift+esc}) under the \textsf{Details} tab
shows the current executables running. The \textsf{CPU} column will
show which processes are preoccupying the CPU. Check whether one or
more of these processes looks \LaTeX -related (e.g. \texttt{latex.exe}
or \texttt{pdflatex.exe}, or \texttt{miktex-pdftex.exe} if using MiK\TeX ).
Click the \textsf{Name} column to sort the processes by name and look
for the relevant name in the list, select it, and end the process
(click the \textsf{End Task} button). 

I am not familiar with the corresponding situation on Linux or Mac.

\section{Hyperref support vs speed}

If you want the pdf produced from your document to support hyperref
links and show an outline window in your pdf viewer (generally placed
on the left in the viewer) then you need to ensure the checkbox at
\textsf{Document Settings \lyxarrow{} PDF Properties \lyxarrow{} Use
Hyperref Support} is indeed checked. But you don't need to do this
until the final compilation of the document. The advantage of leaving
this until the last is that in a large document with many previews
the time for preview generation is essentially halved. If hyperref
support is enabled, preview generation not only creates all the individual
image files that are the previews (files of extension \verb`.png`)
but also requires the compilation of a single pdf document showing
all the previews in sequence. (Like the previews, the pdf document
`hides' in the termporary directory where \LyX{} does its work.)
In other words, \emph{two} images are created for each preview, the
\verb`.png` image which is the one \LyX{} displays, and another image
buried inside the pdf of all images. That second step does not occur
if hyperref support is disabled. In a small document, this is not
going to matter; in a large document it becomes significant. It is
well worth temporarily turning off hyperref support and then, when
the time for final compilation comes, turning it back on.

\section{Supplementary commands in \protect\LyX}

There are some difficulties using the supplementary commands successfully
with instant preview. 

\subsection{Reuse of earlier previews}

One is that whenever \LyX{} has generated a preview image for a particular
\LaTeX{} expression, it will use that same image whenever it meets
that same \LaTeX{} expression later. That means that a statement like
\verb`\macros[view]{}` and the same statement later will display
the same image, even though there may have been macros defined or
freed in between. The same goes for all the other supplementary functions,
including, for example, \verb`\info{sum}`. A second instance of \verb`\info{sum}`
will display the image generated by the first instance even though
further infinite sums may have been evaluated between the \verb`\info`
statements. 

The remedy is to make some small but insignificant difference to the
\LaTeX{} expression  in the second instance – generally a change in
white space will do. For example: first time \verb`\macros[view]{}`,
second time \verb`\macros[view]{ }` where a space has been inserted
between the braces; or: first time \verb`\info{sum}`, second time
\verb`\info{ sum}` where a space has been inserted before \verb`sum`.
This will ensure \LyX{} doesn't fall back on the previously generated
image.

\subsection{\textquoteleft Stalled\textquoteright{} previews}

It is possible to put content into an ERT inset inside a preview inset
(\textsf{Ctrl+P Ctrl+L}) and for nothing to happen. The preview has
apparently stalled. Certainly this can be the case if there is an
error in the input (e.g. a missing brace) but it also occurs if there
is no output to display. For instance \verb`\constants{ c=300000000 }`
does not produce any visual output. There is nothing for the preview
to display and so the preview inset sits there, apparently stalled.
This is a security measure for previews in \LyX{} to provide at least
some guard against malicious code being run in the preview. If the
preview resolved, it would disappear completely from view in the \LyX{}
window. 

If you find the visual appearance of such apparently stalled previews
distracting, the addition of some displayable content to the preview
will result in it resolving to that content; the content could be
as small as a full stop.

\subsection{Using \texttt{\textbackslash nmcMacros}}

As noted earlier, previews are mini-\LaTeX{} runs, either local or
global. Each local preview is of the form<preamble>
\begin{lyxcode}
\textbackslash begin\{document\}

\textbackslash begin\{preview\}

<stuff>

\textbackslash end\{preview\}

\textbackslash end\{document\}
\end{lyxcode}
Whatever goes into or comes out of the preview is isolated from any
other local preview, unless it is through the preamble or an external
file. Sometimes a global preview run can overcome this problem for
then all the previews lie between the same \verb`\begin{document}`,
\verb`end{document}` statements. However, this does not help with
macro definitions. \verb`\def`, \verb`\newcommand`, \verb`\NewDocumentCommand`
all provide \emph{local} definitions which remain trapped within their
own \verb`\begin{preview}`, \verb`\end{preview}`) statements. Another
preview, say containing an \verb`\eval` command, between a different
pair of \verb`\begin{preview}`, \verb`\end{preview}`) statements,
will not know about the macro definition.

There are (at least) three ways out:
\begin{enumerate}
\item Confine everything to the same preview inset: the definition of a
macro, the \verb`\macros` statement, and the use of the macro in
an \verb`\eval` command. 
\item Confine macro definitions to the preamble (\textsf{Document \lyxarrow{}
Settings \lyxarrow{} \LaTeX{} Preamble}).
\item Within previews use \verb`\gdef` (or \verb`\global\def`) exclusively
for making your macro definitions; this makes the macro available
to all later previews.
\end{enumerate}

\subsection{Using \texttt{\textbackslash nmcConstants}}

Because \verb`\nmcConstants` doesn't use \verb`\def` or \verb`\newcommand`
or \verb`\NewDocumentCommand` it is not subject to the same localisation
problem as \verb`\nmcMacros`, but the reach of a \verb`\constants`
command will still be confined to its own preview unless a \emph{global}
preview run is forced; see above §\ref{subsec:LyXGlobal-vs-local}.

\subsection{Using \texttt{\textbackslash nmcReuse}}

As noted earlier, \LyX{} creates its previews in a temporary directory,
not the document directory. If you want to save values from your current
document – say, \texttt{mydoc.lyx} – to \texttt{mydoc.nmc} then you
do so as described earlier (§\ref{sec:supplReuse}), but the file
\texttt{mydoc.nmc} containing the saved results will be located in
the temporary directory. When \LyX{} is closed the file will be deleted
along with all the other contents of that directory.

Fortunately \LyX{} has a copying mechanism for getting files out of
the temporary directory and into the document directory. When a document
is exported – say to pdf – it is possible to specify a \emph{copier}
to automatically copy back to the document directory or subdirectory
various files in the temporary directory. We want the \texttt{.nmc}
file containing the saved values to be copied back. Go to \textsf{Tools
\lyxarrow{} Preferences \lyxarrow{} File Handling \lyxarrow{} File Formats}
and find \textsf{PDF (pdflatex)} (assuming export to \texttt{pdf}
by this route) in the list of formats. In\textsf{ }the \textsf{Copier}
slot of the dialogue insert the following line of code:
\begin{lyxcode}
{\small python~-tt~\$\$s/scripts/ext\_copy.py~-e~nmc,pdf~-d~\$\$i~\$\$o}{\small\par}
\end{lyxcode}
\verb`ext_copy.py` is a python script that is supplied with \LyX .
The \texttt{-e nmc,pdf -d} part of the line tells \texttt{ext\_copy.py}
that on export to \texttt{pdf} by the \texttt{pdflatex} route\texttt{
}to copy any files with the extensions \texttt{.nmc} or \texttt{.pdf}
from the temporary directory where \LyX{} does its work back to the
document directory – the \verb`-d` option (which became available
with \LyX{} 2.3.0). 

But if you have a complex document, it may take too much time to want
to export to pdf before closing \LyX , particularly if there are a
lot of evaluations in the document. Much faster is to export to \emph{plain
text}, not because you want a plain text version of your document
but because it too can be used to trigger the copier mechanism. Go
to \textsf{Tools \lyxarrow{} Preferences \lyxarrow{} File Handling \lyxarrow{}
File Formats} and find \textsf{Plain text} in the list of formats.
In the \textsf{Copier} slot enter
\begin{lyxcode}
{\small python~-tt~\$\$s/scripts/ext\_copy.py~-e~nmc~-d~\$\$i~\$\$o}{\small\par}
\end{lyxcode}
The only difference from the previous copier command is the absence
of \texttt{pdf}.\footnote{I'm assuming that you don't actually want the plain text version of
the file copied back. If you do, then change \texttt{-e nmc} to \texttt{-e
nmc,txt}.} This will copy \texttt{mydoc.nmc} with its saved values from the
temporary directory back to the document directory. To effect the
export, go to \textsf{File \lyxarrow{} Export }and find \textsf{Plain
text} in the list of formats and click on it. 

A shortcut would be nice. For that go to \textsf{Tools \lyxarrow{}
Preferences \lyxarrow{} Editing \lyxarrow{} Shortcuts}, click on \textsf{New},
enter \texttt{buffer-export text} in the \textsf{Function:} slot,
click on the blank key against \textsf{Shortcut:} and type your shortcut.
You may have to try a number before you find one that hasn't already
been assigned. (I'm using \textsf{Ctrl+}; for no particular reason
beyond the fact that it fits under the fingers easily and saving values
to the document directory has a punctuation-like feel to it, a pause
in the process of writing.) It is now an easy matter to press the
shortcut at the end of a \LyX{} session to copy all the values saved
in \texttt{mydoc.nmc} back to a file of the same name in the document
directory. And it is brisk, not least because plain text export ignores
ERT insets (and hence preview insets wrapped around ERT insets), nor
does it evaluate \verb`\eval` commands in math insets.

\subsubsection{A final tweak?}

But one still needs to \emph{remember} to press the shortcut. The
thought arises: can \emph{closing} the current document trigger the
copying process? \LyX{} provides a means of linking two commands and
assigning a keyboard shortcut to them with its \texttt{command-sequence}
\LyX{} function. I suggest assigning a shortcut\textsf{ }to
\begin{lyxcode}
command-sequence~buffer-export~text;~view-close
\end{lyxcode}
Indeed, why not reassign the current shortcut for \texttt{view-close},\texttt{
}which is \textsf{Ctrl+W} on my system, to this command sequence?
(I use the \texttt{cua} key bindings – check the \textsf{Bind file:}
slot in \textsf{Tools \lyxarrow{} Preferences \lyxarrow{} Editing \lyxarrow{}
Shortcuts}.)

Please note, however, that \emph{this will work as intended only from
\LyX{} 2.4.0}.\footnote{Due for release in 2021.} For \LyX{} 2.3 and
earlier, the command sequence will generally fail because of `asynchronous'
processing – \texttt{buffer-export }and \texttt{view-close} use different
threads and the latter may well start before the former is complete.
From \LyX{} 2.4.0 this defect has been fixed. You press your shortcut,
the export to plain text occurs and the \texttt{.nmc} file is copied
back to the document directory, then the current view is closed.

Note that in the other direction, the \verb`.nmc` file in your document
directory is \emph{automatically} copied to the temporary directory
when needed. Nothing needs to be done by you, the user.

\subsubsection{Use of \protect\LyX{} notes}

The central fact about a \LyX{} note is that it does not contribute
to the pdf. But instant preview still works there. This suggests a
possibility: that a calculation be performed within a \LyX{} note and
the result saved using \verb`\nmcReuse` within the same note. The
saved value is now available \emph{from file} for use elsewhere in
the document. In this way, some selected content from a LyX note \emph{can}
find its way into the pdf when the document is compiled.

\chapter{Reference summary}

\section{Commands defined in \texttt{numerica}}
\begin{enumerate}
\item \texttt{\textbackslash nmcEvaluate, \textbackslash eval }
\item \texttt{\textbackslash q, \textbackslash Q }(`cleave' commands)
\item \texttt{\textbackslash nmcInfo, \textbackslash info}
\item \texttt{\textbackslash nmcMacros, \textbackslash macros}
\item \texttt{\textbackslash nmcConstants, \textbackslash constants}
\item \texttt{\textbackslash nmcReuse, \textbackslash reuse}
\end{enumerate}
Provided they have not already been defined when \texttt{numerica}
is loaded, the following commands are defined in \texttt{numerica}
using \verb`\DeclareMathOperator` from \texttt{amsmath} :
\begin{enumerate}
\item \texttt{\textbackslash arccsc, \textbackslash arcsec, \textbackslash arccot}
\item \texttt{\textbackslash csch, \textbackslash sech}
\item \texttt{\textbackslash asinh, \textbackslash acosh, \textbackslash atanh,
\textbackslash acsch, \textbackslash asech, \textbackslash acoth}
\item \texttt{\textbackslash sgn, \textbackslash lb}
\end{enumerate}
Provided they have not already been defined, the following commands
are defined in \texttt{numerica} using \verb`\DeclarePairedDelimiter`
from \texttt{mathtools}:
\begin{lyxcode}
\textbackslash abs,~\textbackslash ceil,~\textbackslash floor
\end{lyxcode}
The following commands have been redefined in \texttt{numerica} to
give more spacing around the underlying \verb`\wedge` and \verb`\vee`
symbols:
\begin{lyxcode}
\textbackslash land,~\textbackslash lor
\end{lyxcode}

\section{\textquoteleft Digestible\textquoteright{} content}

\texttt{numerica} knows how to deal with the following content, meaning
that any of these elements occurring within an \verb`\eval` command
should not of itself cause a \texttt{numerica} error. Not all formatting
commands affect display of the output.
\begin{enumerate}
\item variable names (sequences of tokens given values in the variable~=~value
list)
\item digits, decimal point
\begin{enumerate}
\item \texttt{1, 2, 3, 4, 5, 6, 7, 8, 9, 0, .}
\end{enumerate}
\item constants
\begin{enumerate}
\item \texttt{e, \textbackslash pi, \textbackslash gamma, \textbackslash phi,
\textbackslash deg, \textbackslash infty}
\end{enumerate}
\item arithmetic operators
\begin{enumerate}
\item \texttt{+, -, {*}, /, \textasciicircum , \textbackslash times, \textbackslash cdot,
\textbackslash div}
\end{enumerate}
\item logical operators
\begin{enumerate}
\item \texttt{\textbackslash wedge, \textbackslash land, \textbackslash vee,
\textbackslash lor, \textbackslash neg, \textbackslash lnot}
\end{enumerate}
\item comparisons
\begin{enumerate}
\item \texttt{=, <, >, \textbackslash ne, \textbackslash neq, \textbackslash le,
\textbackslash leq, \textbackslash ge, \textbackslash geq}
\item (if \texttt{amssymb} loaded) \texttt{\textbackslash nless, \textbackslash ngtr,
\textbackslash geqq, \textbackslash geqslant, \textbackslash leqq,
\textbackslash leqslant, \textbackslash ngeq, \textbackslash ngeqq,
\textbackslash ngeqslant, \textbackslash nleq, \textbackslash nleqq,
\textbackslash nleqslant}
\end{enumerate}
\item brackets, bracket-like elements, modifiers
\begin{enumerate}
\item \texttt{( ), {[} {]}, \textbackslash\{ \textbackslash\}}
\item \texttt{\textbackslash lparen \textbackslash rparen} (from \texttt{mathtools})\texttt{,
\textbackslash lbrack \textbackslash rbrack, \textbackslash lbrace
\textbackslash rbrace}
\item \texttt{\textbackslash lvert \textbackslash rvert, \textbackslash lfloor
\textbackslash rfloor, \textbackslash lceil \textbackslash rceil}
\item \texttt{| |} (no nesting, deprecated)
\item \texttt{\textbackslash left \textbackslash right, \textbackslash bigl
\textbackslash bigr, \textbackslash Bigl \textbackslash Bigr, \textbackslash biggl
\textbackslash biggr, \textbackslash Biggl \textbackslash Biggr} 
\item \texttt{.} \texttt{/ |} (used with a modifier)
\item \texttt{\textbackslash abs{[}{]}\{\}, \textbackslash abs{*}\{\},
\textbackslash floor{[}{]}\{\}, \textbackslash floor{*}\{\}, \textbackslash ceil{[}{]}\{\},
\textbackslash ceil{*}\{\}} 
\end{enumerate}
\item unary functions (in the mathematical sense)
\begin{enumerate}
\item \texttt{\textbackslash sin, \textbackslash cos, \textbackslash tan,
\textbackslash csc, \textbackslash sec, \textbackslash cot}
\item \texttt{\textbackslash arcsin, \textbackslash arccos, \textbackslash arctan,
arccsc, \textbackslash arcsec, \textbackslash arccot }
\item \texttt{\textbackslash sin\textasciicircum\{-1\}, \textbackslash cos\textasciicircum\{-1\},
\textbackslash tan\textasciicircum\{-1\}, \textbackslash csc\textasciicircum\{-1\},
\textbackslash sec\textasciicircum\{-1\}, \textbackslash cot\textasciicircum\{-1\}}
\item \texttt{\textbackslash sinh, \textbackslash cosh, \textbackslash tanh,
\textbackslash csch, \textbackslash sech, \textbackslash coth }
\item \texttt{\textbackslash asinh, \textbackslash acosh, \textbackslash atanh,
\textbackslash csch, \textbackslash sech, \textbackslash acoth}
\item \texttt{\textbackslash sinh\textasciicircum\{-1\}, \textbackslash cosh\textasciicircum\{-1\},
\textbackslash tanh\textasciicircum\{-1\}, \textbackslash csch\textasciicircum\{-1\},
\textbackslash sech\textasciicircum\{-1\}, \textbackslash acoth\textasciicircum\{-1\}}
\item \texttt{\textbackslash exp, \textbackslash lb, \textbackslash lg,
\textbackslash ln, \textbackslash log, \textbackslash log\_\{\},
\textbackslash sgn, \textbackslash surd}
\item \texttt{\textbackslash sqrt\{\}, \textbackslash abs{[}{]}\{\}, \textbackslash abs{*}\{\},
\textbackslash floor{[}{]}\{\}, \textbackslash floor{*}\{\}, \textbackslash ceil{[}{]}\{\},
\textbackslash ceil{*}\{\}}
\item \texttt{!, !! }(prepended argument)
\end{enumerate}
\item binary functions 
\begin{enumerate}
\item \texttt{\textbackslash tfrac\{\}\{\}, \textbackslash frac\{\}\{\},
\textbackslash dfrac\{\}\{\}}
\item \texttt{\textbackslash tbinom\{\}\{\}, \textbackslash binom\{\}\{\},
\textbackslash dbinom\{\}\{\}}
\item \texttt{\textbackslash sqrt{[}{]}\{\}}
\end{enumerate}
\item $n$-ary functions
\begin{enumerate}
\item \texttt{\textbackslash min, \textbackslash max, \textbackslash gcd}
\end{enumerate}
\item sum, prod
\begin{enumerate}
\item \texttt{\textbackslash sum\_\{\}\textasciicircum , \textbackslash prod\_\{\}\textasciicircum{} }
\end{enumerate}
\item formatting commands
\begin{enumerate}
\item \texttt{,} (comma, in $n$-ary functions)
\item \texttt{\{\}, \textbackslash\textbackslash , \&, \textbackslash to}
\item \texttt{\textbackslash begin\{\}, \textbackslash end\{\}, \$, \textbackslash{[},
\textbackslash{]}}
\item \texttt{\textbackslash dots, \textbackslash ldots, \textbackslash cdots}
\item \texttt{\textbackslash{} , \textbackslash ,{}, \textbackslash ;,
\textbackslash :, \textbackslash !, \textbackslash >}
\item \texttt{\textbackslash thinspace, \textbackslash medspace, \textbackslash thickspace,}
\item \textbackslash\texttt{negthinspace, \textbackslash negmedspace,
\textbackslash negthickspace,}
\item \textbackslash\texttt{hspace{*}\{\}, \textbackslash mspace\{\},}
\item \texttt{\textbackslash quad, \textbackslash qquad , \textbackslash hfill,
\textbackslash hfil}
\item \texttt{\textbackslash phantom\{\}, \textbackslash vphantom\{\},
\textbackslash hphantom\{\}}
\item \texttt{\textbackslash xmathstrut{[}{]}\{\}} \texttt{, \textbackslash splitfrac\{\}\{\},
\textbackslash splitdfrac\{\}\{\} }(from \texttt{mathtools}), \texttt{\textbackslash mathstrut}
\item \texttt{\textbackslash displaystyle, \textbackslash textstyle, \textbackslash scriptstyle,
\textbackslash scriptscriptstyle}
\item \texttt{\textbackslash label\{\}, \textbackslash ensuremath\{\},
\textbackslash text\{\}, \textbackslash mbox\{\}, \textbackslash smash\{\}}
\item \texttt{\textbackslash color{[}{]}\{\}, \textbackslash textcolor{[}{]}\{\}\{\}}
\end{enumerate}
\item font commands
\begin{enumerate}
\item \texttt{\textbackslash mathrm\{\}, \textbackslash mathit\{\}, \textbackslash mathcal\{\},
\textbackslash mathtt\{\}, \textbackslash mathbf\{\}, \textbackslash mathbb\{\},
\textbackslash mathsf\{\}, \textbackslash mathfrak\{\}, \textbackslash mathscr\{\}}
\item \texttt{\textbackslash mathnormal\{\}, \textbackslash boldsymbol\{\}}
\item \texttt{\textbackslash textrm, \textbackslash textsf, \textbackslash texttt}
\end{enumerate}
\end{enumerate}

\section{Settings}

\subsection{Available \texttt{\textbackslash nmcEvaluate} settings}
\begin{center}
\begin{tabular}[t]{>{\raggedright}p{1.5cm}l>{\raggedright}m{4cm}>{\raggedright}m{4cm}}
\toprule 
{\small key} & {\small type} & {\small meaning} & {\small default}\tabularnewline
\midrule
{\small\texttt{dbg}} & {\small int} & {\small debug `magic' integer} & {\small\texttt{0}}\tabularnewline
{\small\texttt{view}}  &  & {\small equivalent to }{\small\texttt{dbg=1}} & \tabularnewline
{\small\texttt{\textasciicircum}} & {\small char} & {\small exponent mark for sci. notation input} & {\small\texttt{e}}\tabularnewline
{\small\texttt{xx}} & {\small int (0/1)} & {\small multi-token variable switch} & {\small\texttt{1}}\tabularnewline
{\small\texttt{()}} & {\small int (0/1/2)} & {\small trig. function arg. parsing} & {\small\texttt{0}}\tabularnewline
{\small\texttt{o}} & {\small int (0/1)} & {\small degree switch for trig. funcions} & {\small\texttt{1}}\tabularnewline
{\small\texttt{log}} & {\small num} & {\small base of logarithms for }{\small{\small\verb`\log`}} & {\small\texttt{10}}\tabularnewline
{\small\texttt{vv@}} & {\small int (0/1)} & {\small vv-list calculation mode} & {\small\texttt{0}}\tabularnewline
{\small\texttt{vvmode}} & {\small int (0/1)} & {\small equivalent to }{\small\texttt{vv@}} & {\small\texttt{0}}\tabularnewline
{\small\texttt{vvd}} & {\small token(s)} & {\small vv-list display-style spec.} & {\small\texttt{\{,\}\textbackslash mskip 12mu plus 6mu minus 9mu(vv)}}\tabularnewline
{\small\texttt{vvi}} & {\small token(s)} & {\small vv-list text-style spec.} & {\small\texttt{\{,\}\textbackslash mskip 36mu minus 24mu(vv)}}\tabularnewline
{*} &  & {\small switch to suppress equation numbering (if }{\small\texttt{\textbackslash\textbackslash}}{\small{}
in }{\small\texttt{vvd}}{\small )} & \tabularnewline
{\small\texttt{p}} & char(s) & {\small punctuation (esp. in display-style)} & {\small\texttt{,}}\tabularnewline
{\small\texttt{S+}} & {\small int} & {\small extra rounding for stopping criterion, sums} & {\small\texttt{2}}\tabularnewline
{\small\texttt{S?}} & {\small$\text{int}\ge0$} & {\small query stopping with these final terms, sums} & {\small\texttt{0}}\tabularnewline
{\small\texttt{P+}} & {\small int} & {\small extra rounding for stopping criterion, products} & {\small\texttt{2}}\tabularnewline
{\small\texttt{P?}} & {\small$\text{int}\ge0$} & {\small query stopping with these final terms, products} & {\small\texttt{0}}\tabularnewline
{\small\texttt{reuse}} & {\small int} & {\small form of result saved with }{\small{\small\verb`\nmcReuse`}} & {\small\texttt{0}}\tabularnewline
\bottomrule
\end{tabular}
\par\end{center}

\subsection{Available settings for supplementary commands}

All settings for \verb`\nmcEvaluate`, the \verb`view` setting in
particular (although most will be irrelevant), plus for
\begin{itemize}
\item \verb`\nmcMacros`
\begin{itemize}
\item \verb`free` `deregister' a macro from \verb`numerica`
\end{itemize}
\item \verb`\nmcConstants`
\begin{itemize}
\item \verb`add` add the new list of constants to the current one
\end{itemize}
\item \verb`\nmcReuse`
\begin{itemize}
\item \verb`delete` remove the listed control sequences from the \verb`.nmc`
file
\item \verb`renew` change the value of a control sequence in the \verb`.nmc`
file
\end{itemize}
\end{itemize}

\subsection{Available configuration file settings}
\begin{center}
\bigskip{}
\begin{tabular}{ll}
\toprule 
key & default\tabularnewline
\midrule
{\small\texttt{rounding}} & {\small\texttt{6}}\tabularnewline
{\small\texttt{pad}} & {\small\texttt{0}}\tabularnewline
{\small\texttt{output-sci-notation}} & {\small\texttt{0}}\tabularnewline
{\small\texttt{output-exponent-char}} & {\small\texttt{e}}\tabularnewline
{\small\texttt{input-sci-notation}} & {\small\texttt{0}}\tabularnewline
{\small\texttt{input-exponent-char}} & {\small\texttt{e}}\tabularnewline
{\small\texttt{multitoken-variables}} & {\small\texttt{1}}\tabularnewline
{\small\texttt{use-degrees}} & {\small\texttt{0}}\tabularnewline
{\small\texttt{logarithm-base}} & {\small\texttt{10}}\tabularnewline
{\small\texttt{intify-rounding}} & {\small\texttt{14}}\tabularnewline
{\small\texttt{vv-display}} & {\small\texttt{\{,\}\textbackslash mskip 36mu minus 24mu(vv)}}\tabularnewline
{\small\texttt{vv-inline}} & {\small\texttt{\{,\}\textbackslash mskip 12mu  6mu minus 9mu(vv)}}\tabularnewline
{\small\texttt{sum-extra-rounding}} & {\small\texttt{2}}\tabularnewline
{\small\texttt{sum-query-terms}} & {\small\texttt{0}}\tabularnewline
{\small\texttt{prod-extra-rounding}} & {\small\texttt{2}}\tabularnewline
{\small\texttt{prod-query-terms}} & {\small\texttt{0}}\tabularnewline
{\small\texttt{eval-reuse}} & {\small\texttt{0}}\tabularnewline
\bottomrule
\end{tabular} 
\par\end{center}
\end{document}
